% Generated mechanically from frozen input p12/ko/build/ch106/stacks-more-morphisms_full.tex.
% Do not edit this generated file; edit the bound locale source instead.

% OK, start here.
%

\setcounter{chapter}{105}
\chapter{스택의 사상에 관한 추가 내용}



\phantomsection
\label{stacks-more-morphisms-section-phantom}


\section{서론}
\label{stacks-more-morphisms-section-introduction}

\noindent
이 장에서는 대수적 스택의 사상이 지니는 성질을 계속 연구한다.
대각사상이 표현가능한 준분리 대수적 스택의 경우에는
\cite{LM-B}를 참고할 수 있다.





\section{규약과 용어의 관용적 사용}
\label{stacks-more-morphisms-section-conventions}

\noindent
『스택의 성질』의 절
\ref{stacks-properties-section-conventions}에서 도입한 규약과
용어의 관용적 사용을 계속 따른다.







\section{두껍게 하기}
\label{stacks-more-morphisms-section-thickenings}

\noindent
다음 용어가 완전히 표준적이지는 않을 수 있지만 편리하다.
대수적 스택 $\mathcal{X}$의 닫힌 부분스택이 $\mathcal{Y}$이면,
사상 $\mathcal{Y} \to \mathcal{X}$는 표현가능하다.

\begin{definition}
\label{stacks-more-morphisms-definition-thickening}
두껍게 하기.
\begin{enumerate}
\item $\mathcal{X}$가 $\mathcal{X}'$의 닫힌 부분스택이고
그에 대응하는 위상공간들이 같을 때, 대수적 스택 $\mathcal{X}'$를
대수적 스택 $\mathcal{X}$의 {\it 두껍게 하기}라 한다.
\item 두 두껍게 하기 $\mathcal{X} \subset \mathcal{X}'$와
$\mathcal{Y} \subset \mathcal{Y}'$가 주어졌다고 하자.
{\it 두껍게 하기의 사상}이란 $f'|_\mathcal{X}$가 닫힌 부분스택
$\mathcal{Y}$를 통해 분해되는 대수적 스택의 사상
$f' : \mathcal{X}' \to \mathcal{Y}'$을 말한다. 이 상황에서
$f = f'|_\mathcal{X} : \mathcal{X} \to \mathcal{Y}$로 두고,
$(f, f') : (\mathcal{X} \subset \mathcal{X}') \to
(\mathcal{Y} \subset \mathcal{Y}')$를 두껍게 하기의 사상이라 한다.
\item $\mathcal{Z}$를 대수적 스택이라 하자. 마찬가지로
{\it $\mathcal{Z}$ 위의 두껍게 하기}와
{\it $\mathcal{Z}$ 위의 두껍게 하기의 사상}을 정의한다.
이는 대수적 스택 $\mathcal{X}'$와 $\mathcal{Y}'$에
$\mathcal{Z}$로 가는 구조 사상이 주어져 있고, $f'$가 대수적 스택의
적절한 $2$-가환 도식에 들어맞는다는 뜻이다.
\end{enumerate}
\end{definition}

\noindent
$\mathcal{X} \subset \mathcal{X}'$를 대수적 스택의 두껍게 하기라 하자.
$U'$을 스킴이라 하고 $U' \to \mathcal{X}'$를 전사 매끄러운
사상이라 하자. $U = \mathcal{X} \times_{\mathcal{X}'} U'$로 두면
두껍게 하기의 사상
$$
(U \subset U') \longrightarrow (\mathcal{X} \subset \mathcal{X}')
$$
을 얻고, $U \to \mathcal{X}$는 전사 매끄러운 사상이다.
두껍게 하기 $U \subset U'$가 지니는 대응하는 성질로부터
두껍게 하기 $\mathcal{X} \subset \mathcal{X}'$의 성질을 흔히 추론할 수 있다.
때로는 용어를 관용적으로 사용하여, 사상
$\mathcal{X} \to \mathcal{X}'$가 닫힌 몰입이고 전단사
$|\mathcal{X}| \to |\mathcal{X}'|$를 유도할 때 이 사상을
두껍게 하기라고 한다.

\begin{lemma}
\label{stacks-more-morphisms-lemma-thickening}
$i : \mathcal{X} \to \mathcal{X}'$를 대수적 스택의 사상이라 하자.
다음은 동치이다.
\begin{enumerate}
\item $i$는 대수적 스택의 두껍게 하기이다(위와 같은 용어의 관용적 사용).
\item $i$는 대수공간으로 표현가능하고, 『스택의 성질』의 절
\ref{stacks-properties-section-properties-morphisms}의 의미에서
두껍게 하기이다.
\end{enumerate}
이 경우 $i$는 닫힌 몰입이자 보편 위상동형이다.
\end{lemma}

\begin{proof}
『대수공간의 사상에 관한 추가 내용』의 보조정리
\ref{spaces-more-morphisms-lemma-descending-property-thickening}와
\ref{spaces-more-morphisms-lemma-base-change-thickening}에 따르면,
대수공간의 사상이 (일차) 두껍게 하기라는 성질 $P$는 밑에서 fpqc 국소적이고
밑변환 아래 안정적이다. 따라서 『스택의 성질』의 절
\ref{stacks-properties-section-properties-morphisms}의 논의를 실제로 적용할 수 있다.
이제 (1)과 (2)의 동치는 $P = P_1 + P_2$라는 사실에서 따른다.
여기서 $P_1$은 닫힌 몰입이라는 성질이고 $P_2$는 전사라는 성질이다.
(엄밀히 말하면, 모든 개념이 서로 들어맞는다는 것을 확인하려면
『대수공간의 사상에 관한 추가 내용』의 정의
\ref{spaces-more-morphisms-definition-thickening},
『스택의 성질』의 정의 \ref{stacks-properties-definition-immersion}와
그 뒤의 논의, 『대수공간의 사상』의 보조정리
\ref{spaces-morphisms-lemma-surjective-representable},
그리고 『스택의 성질』의 절 \ref{stacks-properties-section-surjective}도
참고해야 한다.) 마지막 명제는 앞의 내용에서 명백하다.
\end{proof}

\noindent
앞으로 이 보조정리를 별도의 언급 없이 사용한다. 이 보조정리에서 사용한 것과
같은 『대수공간의 사상에 관한 추가 내용』의 보조정리
\ref{spaces-more-morphisms-lemma-descending-property-thickening}와
\ref{spaces-more-morphisms-lemma-base-change-thickening}를 이용하면
일차 두껍게 하기를 다음과 같이 정의할 수 있다.

\begin{definition}
\label{stacks-more-morphisms-definition-first-order-thickening}
$\mathcal{X}$가 $\mathcal{X}'$의 닫힌 부분스택이고
$\mathcal{X} \to \mathcal{X}'$가 『스택의 성질』의 절
\ref{stacks-properties-section-properties-morphisms}의 의미에서
일차 두껍게 하기일 때, 대수적 스택 $\mathcal{X}'$를
대수적 스택 $\mathcal{X}$의 {\it 일차 두껍게 하기}라 한다.
\end{definition}

\noindent
$(U \subset U') \to (\mathcal{X} \subset \mathcal{X}')$가 위와 같은
스킴에 의한 매끄러운 덮개라면, 이는 단순히 $U \subset U'$가
일차 두껍게 하기라는 뜻이다. 이제 빠질 수 없는 보조정리들을 서술한다.

\begin{lemma}
\label{stacks-more-morphisms-lemma-base-change-thickening}
$\mathcal{Y} \subset \mathcal{Y}'$를 대수적 스택의 두껍게 하기라 하자.
$\mathcal{X}' \to \mathcal{Y}'$를 대수적 스택의 사상이라 하고
$\mathcal{X} = \mathcal{Y} \times_{\mathcal{Y}'} \mathcal{X}'$로 두자.
그러면
$(\mathcal{X} \subset \mathcal{X}') \to (\mathcal{Y} \subset \mathcal{Y}')$
는 두껍게 하기의 사상이다. $\mathcal{Y} \subset \mathcal{Y}'$가
일차 두껍게 하기이면 $\mathcal{X} \subset \mathcal{X}'$도
일차 두껍게 하기이다.
\end{lemma}

\begin{proof}
위의 논의와 『스택의 성질』의 절
\ref{stacks-properties-section-properties-morphisms}, 그리고
『대수공간의 사상에 관한 추가 내용』의 보조정리
\ref{spaces-more-morphisms-lemma-base-change-thickening}를 보라.
\end{proof}

\begin{lemma}
\label{stacks-more-morphisms-lemma-composition-thickening}
$\mathcal{X} \subset \mathcal{X}'$와 $\mathcal{X}' \subset \mathcal{X}''$가
대수적 스택의 두껍게 하기이면,
$\mathcal{X} \subset \mathcal{X}''$도 그러하다.
\end{lemma}

\begin{proof}
위의 논의와 『스택의 성질』의 절
\ref{stacks-properties-section-properties-morphisms}, 그리고
『대수공간의 사상에 관한 추가 내용』의 보조정리
\ref{spaces-more-morphisms-lemma-composition-thickening}를 보라.
\end{proof}

\begin{example}
\label{stacks-more-morphisms-example-reduction-thickening}
$\mathcal{X}'$를 대수적 스택이라 하자. 그러면 $\mathcal{X}'$는
축약 $\mathcal{X}'_{red}$의 두껍게 하기이다. 『스택의 성질』의 정의
\ref{stacks-properties-definition-reduced-induced-stack}를 보라.
더욱이 $\mathcal{X} \subset \mathcal{X}'$가 대수적 스택의 두껍게 하기이면
$\mathcal{X}'_{red} = \mathcal{X}_{red} \subset \mathcal{X}$이다.
달리 말하면, $\mathcal{X} = \mathcal{X}'_{red}$일 필요충분조건은
$\mathcal{X}$가 축약 대수적 스택인 것이다.
\end{example}

\begin{lemma}
\label{stacks-more-morphisms-lemma-reduced-diagonal}
$(f, f') : (\mathcal{X} \subset \mathcal{X}') \to
(\mathcal{Y} \subset \mathcal{Y}')$를 대수적 스택의
두껍게 하기의 사상이라 하자. 그러면
$\mathcal{X} \times_\mathcal{Y} \mathcal{X} \to
\mathcal{X}' \times_{\mathcal{Y}'} \mathcal{X}'$
는 두껍게 하기이고, 표준 도식
$$
\xymatrix{
\mathcal{X} \ar[r]_-\Delta \ar[d] &
\mathcal{X} \times_\mathcal{Y} \mathcal{X} \ar[d] \\
\mathcal{X}' \ar[r]^-{\Delta'} &
\mathcal{X}' \times_{\mathcal{Y}'} \mathcal{X}'
}
$$
은 카르테시안이다.
\end{lemma}

\begin{proof}
$\mathcal{X} \to \mathcal{Y}'$가 닫힌 부분스택 $\mathcal{Y}$를 통해
분해되므로
$\mathcal{X} \times_\mathcal{Y} \mathcal{X} =
\mathcal{X} \times_{\mathcal{Y}'} \mathcal{X}$이다. 따라서
$\mathcal{X} \times_\mathcal{Y} \mathcal{X} \to
\mathcal{X}' \times_{\mathcal{Y}'} \mathcal{X}'$
는 합성
$$
\mathcal{X} \times_{\mathcal{Y}'} \mathcal{X} \to
\mathcal{X} \times_{\mathcal{Y}'} \mathcal{X}' \to
\mathcal{X}' \times_{\mathcal{Y}'} \mathcal{X}'
$$
과 동형이다. 이 두 사상은 모두 두껍게 하기의 밑변환이므로
두껍게 하기이다(보조정리 \ref{stacks-more-morphisms-lemma-base-change-thickening}).
따라서 그 합성도 그러하다(보조정리
\ref{stacks-more-morphisms-lemma-composition-thickening}).
$\mathcal{X} \to \mathcal{X}'$가 단사상이므로, 보조정리의 마지막 명제는
$\mathcal{X} \to \mathcal{X}' \to \mathcal{Y}'$에
『스택의 성질』의 보조정리
\ref{stacks-properties-lemma-monomorphism-diagonal}를 적용하여 얻는다.
\end{proof}

\begin{lemma}
\label{stacks-more-morphisms-lemma-thickening-diagonals}
$(f, f') : (\mathcal{X} \subset \mathcal{X}') \to
(\mathcal{Y} \subset \mathcal{Y}')$를 대수적 스택의 두껍게 하기의
사상이라 하자. $\Delta : \mathcal{X} \to \mathcal{X} \times_\mathcal{Y} \mathcal{X}$와
$\Delta' : \mathcal{X}' \to \mathcal{X}' \times_{\mathcal{Y}'} \mathcal{X}'$을
각각 대응하는 대각사상이라 하자. 그러면 다음 목록의 각 성질을
$\Delta$가 만족할 필요충분조건은 $\Delta'$가 만족하는 것이다.
(a) 스킴으로 표현가능, (b) 아핀, (c) 전사, (d) 준콤팩트,
(e) 보편 닫힘, (f) 정수적, (g) 준분리, (h) 분리,
(i) 보편 단사, (j) 보편 열림, (k) 국소 준유한,
(l) 유한, (m) 비분기, (n) 단사상, (o) 몰입,
(p) 닫힌 몰입, (q) 고유.
\end{lemma}

\begin{proof}
다음이 두껍게 하기의 사상임에 유의하자(보조정리
\ref{stacks-more-morphisms-lemma-reduced-diagonal}).
$$
(\Delta, \Delta') :
(\mathcal{X} \subset \mathcal{X}')
\longrightarrow
(\mathcal{X} \times_\mathcal{Y} \mathcal{X} \subset
\mathcal{X}' \times_{\mathcal{Y}'} \mathcal{X}')
$$
더욱이 『스택의 사상』의 보조정리
\ref{stacks-morphisms-lemma-properties-diagonal}에 따라
$\Delta$와 $\Delta'$는 대수공간으로 표현가능하다.
따라서 『스택의 성질』의 절
\ref{stacks-properties-section-properties-morphisms}의 논의를 거쳐,
(a), (b), (c), (d), (e), (f), (g), (h), (i), (j)의 경우는
『대수공간의 사상에 관한 추가 내용』의 보조정리
\ref{spaces-more-morphisms-lemma-thicken-property-morphisms}에서 따른다.

\medskip\noindent
보조정리 \ref{stacks-more-morphisms-lemma-reduced-diagonal}에 따르면
$\mathcal{X} = (\mathcal{X} \times_\mathcal{Y} \mathcal{X})
\times_{(\mathcal{X}' \times_{\mathcal{Y}'} \mathcal{X}')} \mathcal{X}'$이다.
더욱이 앞서 언급한 『스택의 사상』의 보조정리
\ref{stacks-morphisms-lemma-properties-diagonal}에 따라
$\Delta$와 $\Delta'$는 국소 유한형이다. 따라서 (k), (l), (m),
(n), (o), (p), (q)의 경우는 『대수공간의 사상에 관한 추가 내용』의
보조정리
\ref{spaces-more-morphisms-lemma-properties-that-extend-over-thickenings}에서 따른다.
\end{proof}

\noindent
그 결과 다음과 같은 보기 좋은 결과를 얻는다.

\begin{lemma}
\label{stacks-more-morphisms-lemma-thickening-properties}
\begin{reference}
\cite[Theorem 2.2.5]{Conrad-moduli}
\end{reference}
$\mathcal{X} \subset \mathcal{X}'$를 대수적 스택의 두껍게 하기라 하자.
그러면
\begin{enumerate}
\item $\mathcal{X}$가 대수공간일 필요충분조건은 $\mathcal{X}'$가
대수공간인 것이다.
\item $\mathcal{X}$가 스킴일 필요충분조건은 $\mathcal{X}'$가 스킴인 것이다.
\item $\mathcal{X}$가 DM일 필요충분조건은 $\mathcal{X}'$가 DM인 것이다.
\item $\mathcal{X}$가 준-DM일 필요충분조건은 $\mathcal{X}'$가 준-DM인 것이다.
\item $\mathcal{X}$가 분리일 필요충분조건은 $\mathcal{X}'$가 분리인 것이다.
\item $\mathcal{X}$가 준분리일 필요충분조건은 $\mathcal{X}'$가
준분리인 것이다.
\item 여기에 더 추가한다.
\end{enumerate}
\end{lemma}

\begin{proof}
각 경우를 대각사상에 관한 문제로 환원한 다음,
두껍게 하기의 사상
$$
(\mathcal{X} \subset \mathcal{X}') \to
\left(\Spec(\mathbf{Z}) \subset \Spec(\mathbf{Z})\right)
$$
에 보조정리 \ref{stacks-more-morphisms-lemma-thickening-diagonals}를 적용한다.
먼저 『대수적 스택』의 정의 \ref{algebraic-definition-viewed-as}를 이용하여
$\mathcal{X} \subset \mathcal{X}'$를 $\Spec(\mathbf{Z})$ 위의
대수적 스택의 두껍게 하기로 본다.

\medskip\noindent
경우 (1). 대수적 스택이 대수공간일 필요충분조건은 그 대각사상이
단사상인 것이다. 『스택의 사상』의 보조정리
\ref{stacks-morphisms-lemma-hierarchy}를 보라(이는 『대수적 스택』의 명제
\ref{algebraic-proposition-algebraic-stack-no-automorphisms}에서도
즉시 따른다).

\medskip\noindent
경우 (2). (1)에 의해 $\mathcal{X}$와 $\mathcal{X}'$가
대수공간이라고 가정해도 되며, 이제 『대수공간의 사상에 관한 추가 내용』의
보조정리 \ref{spaces-more-morphisms-lemma-thickening-scheme}를 사용할 수 있다.

\medskip\noindent
경우 (3)--(6). 이 각 경우는 대각사상에 관한 한 조건에 대응한다.
『스택의 사상』의 정의 \ref{stacks-morphisms-definition-separated}와
\ref{stacks-morphisms-definition-absolute-separated}를 보라.
\end{proof}








\section{두껍게 하기의 사상}
\label{stacks-more-morphisms-section-morphisms-thickenings}

\noindent
$(f, f') : (\mathcal{X} \subset \mathcal{X}') \to
(\mathcal{Y} \subset \mathcal{Y}')$가 대수적 스택의 두껍게 하기의
사상이면, 흔히 사상 $f$의 성질을 $f'$가 물려받는다.
여기에는 몇 가지 변형이 있다.

\begin{lemma}
\label{stacks-more-morphisms-lemma-thicken-property-morphisms}
$(f, f') : (\mathcal{X} \subset \mathcal{X}') \to
(\mathcal{Y} \subset \mathcal{Y}')$를 대수적 스택의
두껍게 하기의 사상이라 하자. 그러면
\begin{enumerate}
\item $f$가 아핀 사상일 필요충분조건은 $f'$가 아핀 사상인 것이다.
\item $f$가 전사일 필요충분조건은 $f'$가 전사인 것이다.
\item $f$가 준콤팩트일 필요충분조건은 $f'$가 준콤팩트인 것이다.
\item $f$가 보편 닫힘일 필요충분조건은 $f'$가 보편 닫힘인 것이다.
\item $f$가 정수적일 필요충분조건은 $f'$가 정수적인 것이다.
\item $f$가 보편 단사일 필요충분조건은 $f'$가 보편 단사인 것이다.
\item $f$가 보편 열림일 필요충분조건은 $f'$가 보편 열림인 것이다.
\item $f$가 준-DM일 필요충분조건은 $f'$가 준-DM인 것이다.
\item $f$가 DM일 필요충분조건은 $f'$가 DM인 것이다.
\item $f$가 (준)분리일 필요충분조건은 $f'$가 (준)분리인 것이다.
\item $f$가 표현가능할 필요충분조건은 $f'$가 표현가능한 것이다.
\item $f$가 대수공간으로 표현가능할 필요충분조건은 $f'$가
대수공간으로 표현가능한 것이다.
\item 여기에 더 추가한다.
\end{enumerate}
\end{lemma}

\begin{proof}
보조정리 \ref{stacks-more-morphisms-lemma-thickening}에 따라 사상
$\mathcal{X} \to \mathcal{X}'$와 $\mathcal{Y} \to \mathcal{Y}'$는
보편 위상동형이다. 따라서 $|f| : |\mathcal{X}| \to |\mathcal{Y}|$에
관한 임의의 조건은 $|f'| : |\mathcal{X}'| \to |\mathcal{Y}'|$에
관한 대응하는 조건과 동치이고, 사상
$\mathcal{Z}' \to \mathcal{Y}'$에 의한 임의의 밑변환 뒤에도
마찬가지이다. 이로써 (2), (3), (4), (6), (7)이 성립한다.

\medskip\noindent
(8), (9), (10), (12)의 경우에는 보조정리
\ref{stacks-more-morphisms-lemma-thickening-diagonals}에서처럼 $f$와 $f'$에 관한 조건을
대각사상 $\Delta$와 $\Delta'$에 관한 조건으로 옮길 수 있다.
『스택의 사상』의 정의 \ref{stacks-morphisms-definition-separated}와
보조정리 \ref{stacks-morphisms-lemma-hierarchy}를 보라.
따라서 이 경우들은 보조정리 \ref{stacks-more-morphisms-lemma-thickening-diagonals}에서 따른다.

\medskip\noindent
(11)의 증명. $f'$가 표현가능하면 $f$도 그러하다. 실제로 스킴 $T$와
사상 $T \to \mathcal{Y}$에 대해
$\mathcal{X} \times_\mathcal{Y} T =
\mathcal{X} \times_{\mathcal{X}'} (\mathcal{X}' \times_{\mathcal{Y}'} T)$이고
$\mathcal{X} \to \mathcal{X}'$가 닫힌 몰입(따라서 표현가능)이기 때문이다.
반대로 $f$가 표현가능하다고 가정하고, $T'$이 스킴인 사상
$T' \to \mathcal{Y}'$을 잡자. 그러면
$$
\mathcal{X} \times_{\mathcal{Y}}
(\mathcal{Y} \times_{\mathcal{Y}'} T') =
\mathcal{X} \times_{\mathcal{X}'}
(\mathcal{X}' \times_{\mathcal{Y}'} T') \to
\mathcal{X}' \times_{\mathcal{Y}'} T'
$$
는 두껍게 하기이고(보조정리 \ref{stacks-more-morphisms-lemma-base-change-thickening}),
그 원천은 스킴이다. 따라서 보조정리
\ref{stacks-more-morphisms-lemma-thickening-properties}에 의해 그 표적도 스킴이다.

\medskip\noindent
(1)과 (5)의 경우, $f$나 $f'$ 가운데 하나가 명시된 성질을 가지면
(11)에 의해 $f$와 $f'$가 모두 표현가능하다. 이 경우 대수공간 $V'$와
전사 매끄러운 사상 $V' \to \mathcal{Y}'$을 택하자.
$V = \mathcal{Y} \times_{\mathcal{Y}'} V'$,
$U' = \mathcal{X}' \times_{\mathcal{Y}'} V'$, 그리고
$U = \mathcal{X} \times_{\mathcal{Y}'} V'$로 두자.
그러면 원하는 결과는 『스택의 성질』의 보조정리
\ref{stacks-properties-lemma-check-property-covering}의 원리를 통해
대수공간의 두껍게 하기의 사상
$(U \subset U') \to (V \subset V')$에 관한 대응하는 결과에서 따른다.
대수공간의 경우에 해당하는 결과는
『대수공간의 사상에 관한 추가 내용』의 보조정리
\ref{spaces-more-morphisms-lemma-thicken-property-morphisms}를 보라.
\end{proof}





\section{대수적 스택의 무한소 변형}
\label{stacks-more-morphisms-section-deform}

\noindent
이 절은 『대수공간의 사상에 관한 추가 내용』의 절
\ref{spaces-more-morphisms-section-deform}에 대응한다.

\begin{lemma}
\label{stacks-more-morphisms-lemma-flatness-morphism-thickenings-fp-over-ft}
대수적 스택의 두껍게 하기들로 이루어진 다음 가환 도식을 생각하자.
$$
\xymatrix{
(\mathcal{X} \subset \mathcal{X}') \ar[rr]_{(f, f')} \ar[rd] & &
(\mathcal{Y} \subset \mathcal{Y}') \ar[ld] \\
& (\mathcal{B} \subset \mathcal{B}')
}
$$
다음을 가정하자.
\begin{enumerate}
\item $\mathcal{Y}' \to \mathcal{B}'$는 국소 유한형이다.
\item $\mathcal{X}' \to \mathcal{B}'$는 평탄하고 국소 유한 표시이다.
\item $f$는 평탄하다.
\item $\mathcal{X} = \mathcal{B} \times_{\mathcal{B}'} \mathcal{X}'$이고
$\mathcal{Y} = \mathcal{B} \times_{\mathcal{B}'} \mathcal{Y}'$이다.
\end{enumerate}
그러면 $f'$는 평탄하고, $|f'|$의 상에 속하는 모든
$y' \in |\mathcal{Y}'|$에 대해 사상
$\mathcal{Y}' \to \mathcal{B}'$가 $y'$에서 평탄하다.
\end{lemma}

\begin{proof}
대수공간 $U'$와 전사 매끄러운 사상 $U' \to \mathcal{B}'$를 택한다.
대수공간 $V'$와 전사 매끄러운 사상
$V' \to U' \times_{\mathcal{B}'} \mathcal{Y}'$을 택한다.
대수공간 $W'$와 전사 매끄러운 사상
$W' \to V' \times_{\mathcal{Y}'} \mathcal{X}'$을 택한다.
$U, V, W$를 각각 $\mathcal{B} \to \mathcal{B}'$에 의한
$U', V', W'$의 밑변환이라 하자. 그러면 $f'$의 평탄성은
$W' \to V'$의 평탄성과 동치이고, $W \to V$가 평탄하다는 것은
주어져 있다. 따라서 대수공간의 두껍게 하기들로 이루어진 도식
$$
\xymatrix{
(W \subset W') \ar[rr] \ar[rd] & & (V \subset V') \ar[ld] \\
& (U \subset U')
}
$$
에 대수공간의 경우의 보조정리를 적용할 수 있다.
『대수공간의 사상에 관한 추가 내용』의 보조정리
\ref{spaces-more-morphisms-lemma-flatness-morphism-thickenings-fp-over-ft}를 보라.
$|f'|$의 상에 속하는 점들에서 $\mathcal{Y}'/\mathcal{B}'$가
평탄하다는 명제도 같은 방식으로 따른다.
\end{proof}

\begin{lemma}
\label{stacks-more-morphisms-lemma-deform-property-fp-over-ft}
대수적 스택의 두껍게 하기들로 이루어진 다음 가환 도식을 생각하자.
$$
\xymatrix{
(\mathcal{X} \subset \mathcal{X}') \ar[rr]_{(f, f')} \ar[rd] & &
(\mathcal{Y} \subset \mathcal{Y}') \ar[ld] \\
& (\mathcal{B} \subset \mathcal{B}')
}
$$
$\mathcal{Y}' \to \mathcal{B}'$는 국소 유한형이고,
$\mathcal{X}' \to \mathcal{B}'$는 평탄하고 국소 유한 표시이며,
$\mathcal{X} = \mathcal{B} \times_{\mathcal{B}'} \mathcal{X}'$이고
$\mathcal{Y} = \mathcal{B} \times_{\mathcal{B}'} \mathcal{Y}'$이라고
가정하자. 그러면
\begin{enumerate}
\item $f$가 평탄할 필요충분조건은 $f'$가 평탄인 것이다.
\label{stacks-more-morphisms-item-flat-fp-over-ft}
\item $f$가 동형사상일 필요충분조건은 $f'$가 동형사상인 것이다.
\label{stacks-more-morphisms-item-isomorphism-fp-over-ft}
\item $f$가 열린 몰입일 필요충분조건은 $f'$가 열린 몰입인 것이다.
\label{stacks-more-morphisms-item-open-immersion-fp-over-ft}
\item $f$가 단사상일 필요충분조건은 $f'$가 단사상인 것이다.
\label{stacks-more-morphisms-item-monomorphism-fp-over-ft}
\item $f$가 국소 준유한일 필요충분조건은 $f'$가 국소 준유한인 것이다.
\label{stacks-more-morphisms-item-quasi-finite-fp-over-ft}
\item $f$가 신토믹일 필요충분조건은 $f'$가 신토믹인 것이다.
\label{stacks-more-morphisms-item-syntomic-fp-over-ft}
\item $f$가 매끄러울 필요충분조건은 $f'$가 매끄러운 것이다.
\label{stacks-more-morphisms-item-smooth-fp-over-ft}
\item $f$가 비분기일 필요충분조건은 $f'$가 비분기인 것이다.
\label{stacks-more-morphisms-item-unramified-fp-over-ft}
\item $f$가 에탈일 필요충분조건은 $f'$가 에탈인 것이다.
\label{stacks-more-morphisms-item-etale-fp-over-ft}
\item $f$가 유한일 필요충분조건은 $f'$가 유한인 것이다.
\label{stacks-more-morphisms-item-finite-fp-over-ft}
\item 여기에 더 추가한다.
\end{enumerate}
\end{lemma}

\begin{proof}
경우 (\ref{stacks-more-morphisms-item-flat-fp-over-ft})는 보조정리
\ref{stacks-more-morphisms-lemma-flatness-morphism-thickenings-fp-over-ft}에서 따른다.

\medskip\noindent
경우 (\ref{stacks-more-morphisms-item-syntomic-fp-over-ft})와
(\ref{stacks-more-morphisms-item-smooth-fp-over-ft})는 보조정리
\ref{stacks-more-morphisms-lemma-flatness-morphism-thickenings-fp-over-ft}의 증명에서 쓴 방법으로
증명할 수 있다. 즉, 대수공간 $U'$와 전사 매끄러운 사상
$U' \to \mathcal{B}'$를 택한다. 대수공간 $V'$와 전사 매끄러운 사상
$V' \to U' \times_{\mathcal{B}'} \mathcal{Y}'$을 택한다.
대수공간 $W'$와 전사 매끄러운 사상
$W' \to V' \times_{\mathcal{Y}'} \mathcal{X}'$을 택한다.
$U, V, W$를 각각 $\mathcal{B} \to \mathcal{B}'$에 의한
$U', V', W'$의 밑변환이라 하자. 그러면 $f$ 및 $f'$에 대한 성질은
각각 $W' \to V'$ 및 $W \to V$에 대한 성질과 동치이다.
따라서 대수공간의 두껍게 하기들로 이루어진 도식
$$
\xymatrix{
(W \subset W') \ar[rr] \ar[rd] & & (V \subset V') \ar[ld] \\
& (U \subset U')
}
$$
에 대수공간의 경우의 보조정리를 적용할 수 있다.
『대수공간의 사상에 관한 추가 내용』의 보조정리
\ref{spaces-more-morphisms-lemma-deform-property-fp-over-ft}를 보라.

\medskip\noindent
경우 (\ref{stacks-more-morphisms-item-unramified-fp-over-ft})와
(\ref{stacks-more-morphisms-item-etale-fp-over-ft})에는 먼저 $f$ 또는 $f'$에 관한 가정으로부터
$f$와 $f'$가 모두 대수적 스택의 DM 사상임을 알 수 있다.
보조정리 \ref{stacks-more-morphisms-lemma-thicken-property-morphisms}를 보라.
이제 대수공간 $U'$와 전사 매끄러운 사상
$U' \to \mathcal{B}'$를 택한다. 대수공간 $V'$와 전사 매끄러운 사상
$V' \to U' \times_{\mathcal{B}'} \mathcal{Y}'$을 택한다.
대수공간 $W'$와 전사 에탈(!) 사상
$W' \to V' \times_{\mathcal{Y}'} \mathcal{X}'$을 택한다.
$U, V, W$를 각각 $\mathcal{B} \to \mathcal{B}'$에 의한
$U', V', W'$의 밑변환이라 하자. 그러면
$W \to V \times_\mathcal{Y} \mathcal{X}$도 전사 에탈이다.
따라서 $f$ 및 $f'$에 대한 성질은 각각
$W' \to V'$ 및 $W \to V$에 대한 성질과 동치이다.
그러므로 대수공간의 두껍게 하기들로 이루어진 도식
$$
\xymatrix{
(W \subset W') \ar[rr] \ar[rd] & & (V \subset V') \ar[ld] \\
& (U \subset U')
}
$$
에 대수공간의 경우의 보조정리를 적용할 수 있다.
『대수공간의 사상에 관한 추가 내용』의 보조정리
\ref{spaces-more-morphisms-lemma-deform-property-fp-over-ft}를 보라.

\medskip\noindent
경우 (\ref{stacks-more-morphisms-item-isomorphism-fp-over-ft}),
(\ref{stacks-more-morphisms-item-open-immersion-fp-over-ft}),
(\ref{stacks-more-morphisms-item-monomorphism-fp-over-ft}),
(\ref{stacks-more-morphisms-item-finite-fp-over-ft})에는 먼저 보조정리
\ref{stacks-more-morphisms-lemma-thicken-property-morphisms}에 의해 $f$와 $f'$가
대수공간으로 표현가능하다고 결론짓는다. 따라서 대수공간 $U'$와
전사 매끄러운 사상 $U' \to \mathcal{B}'$, 대수공간 $V'$와
전사 매끄러운 사상
$V' \to U' \times_{\mathcal{B}'} \mathcal{Y}'$을 택할 수 있고,
그러면 $W' = V' \times_{\mathcal{Y}'} \mathcal{X}'$도 대수공간이다.
$U, V, W$를 각각 $\mathcal{B} \to \mathcal{B}'$에 의한
$U', V', W'$의 밑변환이라 하자. 이때
$W = V \times_\mathcal{Y} \mathcal{X}$이기도 하다.
이제 $W' \to V'$가 각각 동형사상, 열린 몰입, 단사상, 유한일
필요충분조건이 $W \to V$가 같은 성질을 가지는 것임을 보여야 한다.
『스택의 성질』의 보조정리
\ref{stacks-properties-lemma-check-property-covering}를 보라.
따라서 위와 같이 대수공간에 대한 결과를 적용하여 결론을 얻는다.

\medskip\noindent
경우 (\ref{stacks-more-morphisms-item-quasi-finite-fp-over-ft})에는 먼저
『스택의 사상』의 보조정리
\ref{stacks-morphisms-lemma-finite-type-permanence}에 의해
$f$와 $f'$가 국소 유한형임에 유의한다. 한편 보조정리
\ref{stacks-more-morphisms-lemma-thicken-property-morphisms}에 의해 $f$가 준-DM일
필요충분조건은 $f'$가 준-DM인 것이다. $f$ 또는 $f'$가
국소 준유한인지 알아보기 위해 마지막으로 확인할 것
(『스택의 사상』의 정의
\ref{stacks-morphisms-definition-locally-quasi-finite})은
바탕 위상공간에 관한 조건이다. 증명의 첫 문단의 논의에 의해
이 조건을 $f$가 만족할 필요충분조건은 $f'$가 만족하는 것이다.
\end{proof}








\section{아핀의 올림}
\label{stacks-more-morphisms-section-lifting-affines}

\noindent
다음 실선 도식을 생각하자.
$$
\xymatrix{
W \ar[d] \ar@{..>}[r] & W' \ar@{..>}[d] \\
\mathcal{X} \ar[r] & \mathcal{X}'
}
$$
여기서 $\mathcal{X} \subset \mathcal{X}'$는 대수적 스택의
두껍게 하기이고, $W$는 아핀 스킴이며,
$W \to \mathcal{X}$는 매끄럽다. 이 절에서 다루는 문제는
정사각형이 카르테시안이고 $W' \to \mathcal{X}'$가 매끄럽도록
$W'$와 점선 화살표들을 찾을 수 있는가 하는 것이다.
일반적으로는 그 답을 모르지만, $\mathcal{X} \subset \mathcal{X}'$가
일차 두껍게 하기이면 답이 그렇다는 것을 증명할 것이다.

\medskip\noindent
이 문제를 연구하기 위해 다음 범주를 도입한다.

\begin{remark}[올림의 범주]
\label{stacks-more-morphisms-remark-gerbe-of-lifts}
다음 도식을 생각하자.
$$
\xymatrix{
W \ar[d]_x \\
\mathcal{X} \ar[r] & \mathcal{X}'
}
$$
여기서 $\mathcal{X} \subset \mathcal{X}'$는 대수적 스택의
두껍게 하기이고, $W$는 대수공간이며,
$W \to \mathcal{X}$는 매끄럽다. 범주 $\mathcal{C}$와 함자
$$
p : \mathcal{C} \longrightarrow W_{spaces, \etale}
$$
를 다음과 같이 구성한다(표기법은 『대수공간의 성질』의 정의
\ref{spaces-properties-definition-spaces-etale-site}를 보라).
$\mathcal{C}$의 대상은 가환 도식
\begin{equation}
\label{stacks-more-morphisms-equation-object}
\vcenter{
\xymatrix{
U \ar[d]_a \ar[r]_i & U' \ar[dd]^{x'} \\
W \ar[d]_x & \\
\mathcal{X} \ar[r] & \mathcal{X}'
}
}
\end{equation}
을 이루는 계 $(U, U', a, i, x', \alpha)$이다. 여기서 가환성은
$2$-사상 $\alpha : x \circ a \to x' \circ i$로 입증되며,
$U$와 $U'$는 대수공간이고, $a : U \to W$는 에탈,
$x' : U' \to \mathcal{X}'$는 매끄러우며,
$U = \mathcal{X} \times_{\mathcal{X}'} U'$이다.
특히 $U \subset U'$는 두껍게 하기이다. 사상
$$
(U, U', a, i, x', \alpha) \to (V, V', b, j, y', \beta)
$$
은 $(f, f', \gamma)$로 주어진다. 여기서 $f : U \to V$는
$W$ 위의 사상이고, $f' : U' \to V'$는 $U$로 제한하면
$f$가 되는 사상이며, $\gamma : x' \circ f' \to y'$는
아래 도식의 오른쪽 삼각형의 가환성을 입증하는 $2$-사상이다.
\begin{equation}
\label{stacks-more-morphisms-equation-morphism}
\vcenter{
\xymatrix{
& V \ar[ld]_f \ar[ldd]^b \ar[rr]_j & & V' \ar[ld]_{f'} \ar[lddd]^{y'} \\
U \ar[d]_a \ar[rr]_i & & U' \ar[dd]_{x'} \\
W \ar[d]_x & \\
\mathcal{X} \ar[rr] & & \mathcal{X}'
}
}
\end{equation}
마지막으로 $\gamma$가 $\alpha$ 및 $\beta$와 양립할 것을 요구한다.
『범주』의 절 \ref{categories-section-formal-cat-cat}와
\ref{categories-section-2-categories}에 나오는 $2$-범주의 계산법으로 쓰면
이는
$$
\beta = (\gamma \star \text{id}_j) \circ (\alpha \star \text{id}_f)
$$
이다(더 간단히 쓰면 $\beta = j^*\gamma \circ f^*\alpha$이다).
달리 말하면, 대상은 몇 가지 추가 성질을 갖는 가환 도식
(\ref{stacks-more-morphisms-equation-object})이고, 사상은 가환 도식
(\ref{stacks-more-morphisms-equation-morphism})이며, 이는 『스택의 성질』의 주석
\ref{stacks-properties-remark-representable-over}에서 도입한 범주
$\textit{Spaces}/\mathcal{X}'$ 안의 도식이다. 이 서술로부터 $\mathcal{C}$가 범주이고,
$(U, U', a, i, x', \alpha)$를 $a : U \to W$로 보내는 규칙
$p : \mathcal{C} \to W_{spaces, \etale}$가 함자임이 분명하다.
\end{remark}

\begin{lemma}
\label{stacks-more-morphisms-lemma-morphisms-lifts-etale}
(\ref{stacks-more-morphisms-equation-morphism})의 임의의 사상에 대해 사상
$f' : V' \to U'$는 에탈이다.
\end{lemma}

\begin{proof}
실제로 $f : V \to U$는 $W_{spaces, \etale}$의 사상으로서 에탈이고,
$U' \to \mathcal{X}'$와 $V' \to \mathcal{X}'$가 매끄러우며
$U = \mathcal{X} \times_{\mathcal{X}'} U'$와
$V = \mathcal{X} \times_{\mathcal{X}'} V'$이므로
보조정리 \ref{stacks-more-morphisms-lemma-deform-property-fp-over-ft}를 적용할 수 있다.
\end{proof}



\begin{lemma}
\label{stacks-more-morphisms-lemma-gerbe-of-lifts-fibred}
주석 \ref{stacks-more-morphisms-remark-gerbe-of-lifts}에서 구성한 범주
$p : \mathcal{C} \to W_{spaces, \etale}$는 준군 섬유화 범주이다.
\end{lemma}

\begin{proof}
$p$의 섬유 범주들이 준군이라고 주장한다.
(\ref{stacks-more-morphisms-equation-morphism})에서와 같은 사상 $(f, f', \gamma')$에 대해
$f : U \to V$가 동형사상이면, 보조정리
\ref{stacks-more-morphisms-lemma-deform-property-fp-over-ft}에 의해 $f'$도 동형사상이므로
$(f, f', \gamma')$는 동형사상이다.

\medskip\noindent
$W_{spaces, \etale}$의 사상 $f : V \to U$와 $U$ 위의
$\mathcal{C}$의 대상 $\xi = (U, U', a, i, x', \alpha)$를 생각하자.
$V$ 위의 ``당김'' $f^*\xi$를 구성할 것이다.
$b = a \circ f$로 두자. 제한하면 $V$ 위에서 $f$가 되는 에탈 사상
$f' : V' \to U'$을 잡는다(『대수공간의 사상에 관한 추가 내용』의
보조정리 \ref{spaces-more-morphisms-lemma-topological-invariance}).
대응하는 두껍게 하기를 $j : V \to V'$라 쓰자.
$y' = x' \circ f'$와 $\gamma = \text{id} : x' \circ f' \to y'$로 두고,
$$
\beta = \alpha \star \text{id}_f :
x \circ b = x \circ a \circ f \to
x' \circ i \circ f = x' \circ f' \circ j = y' \circ j
$$
로 두자. 그러면
$(f, f', \gamma) : (V, V', b, j, y', \beta) \to
(U, U', a, i, x', \alpha)$가 (\ref{stacks-more-morphisms-equation-morphism})에서와 같은
사상임이 분명하다. 이렇게 구성한 사상 $(f, f', \gamma)$는
강 카르테시안이다(『범주』의 정의
\ref{categories-definition-cartesian-over-C}).
자세한 증명은 생략하지만, 그 핵심 이유는 다음과 같다.
$\mathcal{C}$의 사상
$(g, g', \epsilon) : (Y, Y', c, k, z', \delta) \to (U, U', a, i, x', \alpha)$가 주어지고 어떤 $h : Y \to V$에 대해
$g$가 $g = f \circ h$로 분해된다고 하자. 그러면
『대수공간의 사상에 관한 추가 내용』의 보조정리
\ref{spaces-more-morphisms-lemma-topological-invariance}로부터 유일한 분해
$g' = f' \circ h'$을 얻는다. 이어서
$(h, h', \zeta) :  (Y, Y', c, k, z', \delta) \to
(V, V', b, j, y', \beta)$가 $\mathcal{C}$의 사상이고
$(g, g', \epsilon) = (f, f', \gamma) \circ (h, h', \zeta)$가 되게 하는
필요한 $\zeta$를 만들 수 있다.

\medskip\noindent
따라서 $p : \mathcal{C} \to W_\etale$는 섬유화 범주이다
(『범주』의 정의 \ref{categories-definition-fibred-category}).
위에서 본 대로 섬유 범주들이 준군이라는 사실과 합치면,
『범주』의 보조정리 \ref{categories-lemma-fibred-groupoids}에 의해
$p : \mathcal{C} \to W_\etale$가 준군 섬유화 범주라고 결론짓는다.
\end{proof}

\begin{lemma}
\label{stacks-more-morphisms-lemma-gerbe-of-lifts-stack}
주석 \ref{stacks-more-morphisms-remark-gerbe-of-lifts}에서 구성한 범주
$p : \mathcal{C} \to W_{spaces, \etale}$는 준군 스택이다.
\end{lemma}

\begin{proof}
보조정리 \ref{stacks-more-morphisms-lemma-gerbe-of-lifts-fibred}에 의해 『스택』의 정의
\ref{stacks-definition-stack-in-groupoids}의 첫 번째 조건이 성립한다.
관례대로 대상의 내림을 확인하고 사상의 내림은 독자에게 맡긴다.
따라서 $W_{spaces, \etale}$의 $a : U \to W$,
$W_{spaces, \etale}$에서의 덮개 $\{U_k \to U\}_{k \in K}$,
각 $U_k$ 위의 $\mathcal{C}$의 대상
$\xi_k = (U_k, U'_k, a_k, i_k, x'_k, \alpha_k)$, 그리고 코사이클 조건을
만족하는 제한 사이의 사상
$$
\varphi_{kk'} = (f_{kk'}, f'_{kk'}, \gamma_{kk'}) :
\xi_k|_{U_k \times_U U_{k'}} \to
\xi_{k'}|_{U_k \times_U U_{k'}}
$$
이 주어졌다고 하자. 유효성을 증명할 때 먼저 덮개를 세분해도 된다.
따라서 각 $U_k$가 스킴이라고 가정할 수 있다(원한다면 아핀 스킴이라고
가정할 수도 있다). 다음과 같이 쓰자.
$$
\xi_k|_{U_k \times_U U_{k'}} =
(U_k \times_U U_{k'}, U'_{kk'}, a_{kk'}, x'_{kk'}, \alpha_{kk'})
$$
그러면 $\mathcal{C}$의 사상
$\xi_k|_{U_k \times_U U_{k'}} \to \xi_k$의 두 번째 성분으로부터
에탈 사상(보조정리 \ref{stacks-more-morphisms-lemma-morphisms-lifts-etale})
$s_{kk'} : U'_{kk'} \to U'_k$를 얻는다. 마찬가지로 합성
$$
\xi_k|_{U_k \times_U U_{k'}} \xrightarrow{\varphi_{kk'}}
\xi_{k'}|_{U_k \times_U U_{k'}} \to \xi_{k'}
$$
의 두 번째 성분을 살펴보아 에탈 사상
$t_{kk'} : U'_{kk'} \to U'_{k'}$를 얻는다. 다음 사상이
에탈 동치관계라고 주장한다.
$$
j :
\coprod\nolimits_{(k, k') \in K \times K} U'_{kk'}
\xrightarrow{(\coprod s_{kk'}, \coprod t_{kk'})}
(\coprod\nolimits_{k \in K} U'_k) \times (\coprod\nolimits_{k \in K} U'_k)
$$
먼저, 표시된 사상의 성분 $s, t$가 에탈이라는 것은 이미 보았다.
사상 $j$를
$(\coprod U_k) \times (\coprod U_k) \to (\coprod U'_k) \times (\coprod U'_k)$로 밑변환하면
$$
\coprod\nolimits_{(k, k') \in K \times K} U_k \times_U U_{k'}
\longrightarrow
(\coprod\nolimits_{k \in K} U_k) \times (\coprod\nolimits_{k \in K} U_k)
$$
가 되므로 단사상이다. 따라서 『사상에 관한 추가 내용』의 보조정리
\ref{more-morphisms-lemma-properties-that-extend-over-thickenings}에 의해
$j$는 단사상이다. 마지막으로, 관계 $j$의 대칭성은
$\varphi_{kk'}^{-1}$가 $\varphi_{k'k}$의 ``뒤집기''라는 사실에서 오고
(『스택』의 주석 \ref{stacks-remarks-definition-descent-datum}을 보라),
추이성은 코사이클 조건에서 온다(세부사항은 생략한다).
따라서 $j$에 의한 $\coprod U'_k$의 몫은 대수공간 $U'$이다
(『대수공간』의 정리 \ref{spaces-theorem-presentation}).
위에서 이미 $j$를 $\coprod U_k$에 제한하면
$(\coprod U_k) \times_U (\coprod U_k)$를 얻음을 보았으므로,
두껍게 하기 $i : U \to U'$가 존재한다는 것도 이미 보였다.
마지막으로 $1$-사상 $x'_k : U'_k \to \mathcal{X}'$를 잠시
$U'_k$ 위의 스택 $\mathcal{X}'$의 대상으로 보면,
이 대상들에는 $\mathcal{C}$의 사상 $\varphi_{kk'}$의 세 번째 성분
$\gamma_{kk'}$가 주는 에탈 덮개 $\{U'_k \to U'\}$에 관한
내림 자료가 딸려 있다. $\mathcal{X}'$가 스택이므로 이 내림 자료는
유효하고, 다시 사상의 언어로 옮기면 $1$-사상
$x' : U' \to \mathcal{X}'$를 얻는다. 이때 합성
$U'_k \to U' \to \mathcal{X}'$에는 $\gamma_{kk'}$와 양립하는
$x'_k$로의 동형사상이 딸려 있다. 이는 사상
$\alpha_k : x \circ a_k \to x'_k \circ i_k$가
$\alpha : x \circ a \to x' \circ i$로 붙는다는 뜻이다.
그러면 $\xi = (U, U', a, i, x', \alpha)$가 원하는 $U$ 위의 대상이다.
\end{proof}

\begin{lemma}
\label{stacks-more-morphisms-lemma-etale-local-lifts}
$\mathcal{X} \subset \mathcal{X}'$를 대수적 스택의 두껍게 하기라 하자.
$W$를 대수공간이라 하고 $W \to \mathcal{X}$를 매끄러운 사상이라 하자.
에탈 덮개 $\{W_i \to W\}_{i \in I}$와 각 $i$에 대해
$W_i' \to \mathcal{X}'$가 매끄러운 카르테시안 도식
$$
\xymatrix{
W_i \ar[r] \ar[d] & W_i' \ar[d] \\
\mathcal{X} \ar[r] & \mathcal{X}'
}
$$
이 존재한다.
\end{lemma}

\begin{proof}
스킴 $U'$와 전사 매끄러운 사상 $U' \to \mathcal{X}'$를 택한다.
평소와 같이 $U = \mathcal{X} \times_{\mathcal{X}'} U'$로 둔다.
그러면 $U \to \mathcal{X}$는 전사 매끄러운 사상이다. 따라서 밑변환
$$
V = W \times_{\mathcal{X}} U \longrightarrow W
$$
은 대수공간의 전사 매끄러운 사상이다. 『대수공간의 위상』의 보조정리
\ref{spaces-topologies-lemma-etale-dominates-smooth}에 의해
$W_i \to W$가 $V \to W$를 통해 분해되는 에탈 덮개
$\{W_i \to W\}$를 찾을 수 있다. $W_i$를 아핀들로 덮은 뒤
(『대수공간의 성질』의 보조정리
\ref{spaces-properties-lemma-cover-by-union-affines}), 각 $W_i$가
아핀이라고 가정할 수 있다. $W$를 $W_i$로 바꾸며, 그러면 문제는
다음 문단에서 다루는 상황으로 환원된다.

\medskip\noindent
$W$가 아핀이고 주어진 사상 $W \to \mathcal{X}$가 $U$를 통해
분해된다고 가정하자. 다음과 같다.
$$
W \xrightarrow{i} U \to \mathcal{X}
$$
$W$와 $U$가 $\mathcal{X}$ 위에서 매끄러우므로 $i$는 국소 유한형이다
(『스택의 사상』의 보조정리
\ref{stacks-morphisms-lemma-finite-type-permanence}).
$U$를 $\mathbf{A}^n_U$로 바꾼 뒤 $i$가 몰입이라고 가정할 수 있다.
『사상』의 보조정리
\ref{morphisms-lemma-quasi-affine-finite-type-over-S}를 보라.
『스택의 사상』의 보조정리
\ref{stacks-morphisms-lemma-lci-permanence}에 의해 사상 $i$는
국소 완전 교차이다. 따라서 $i$는 코쥘-정칙 몰입이다
(『인자』의 정의 \ref{divisors-definition-regular-immersion}에서
정의한 의미로). 이는 『사상에 관한 추가 내용』의 보조정리
\ref{more-morphisms-lemma-lci}에서 따른다.

\medskip\noindent
여전히 $W$를 아핀 열린 덮개로 바꾸어도 된다.
각 점 $w \in W$에 대해 다음 조건을 만족하는 아핀 열린집합
$U'_w \subset U'$를 택할 수 있다. 대응하는 아핀 열린집합을
$U_w \subset U$라 하면 $w \in i^{-1}(U_w)$이고,
$i^{-1}(U_w) \to U_w$는
$f_1, \ldots, f_r \in \Gamma(U_w, \mathcal{O}_{U_w})$라는
코쥘-정칙 수열이 잘라내는 닫힌 몰입이다. 이는 코쥘-정칙 몰입의 정의와
『인자』의 보조정리
\ref{divisors-lemma-regular-ideal-sheaf-scheme}에서 따른다.
$W_w = i^{-1}(U_w)$로 두자. 이는 $w \in W$의 아핀 열린 근방이다.
$f_1, \ldots, f_r$의 올림
$f'_1, \ldots, f'_r \in \Gamma(U'_w, \mathcal{O}_{U'_w})$를 택한다.
$U_w \to U'_w$가 아핀 스킴의 닫힌 몰입이므로 이는 가능하다.
$W'_w \subset U'_w$를 $f'_1, \ldots, f'_r$가 잘라내는 닫힌 부분스킴이라 하자.
$W'_w \to \mathcal{X}'$가 매끄럽다고 주장한다. 구성상
$W_w = \mathcal{X} \times_{\mathcal{X}'} W'_w$이므로 이 주장이
증명을 끝낸다.

\medskip\noindent
이 주장을 확인하려면 $\mathcal{X}'$ 위에서 매끄러운 모든 아핀 스킴
$X'$에 대해 밑변환
$W'_w \times_{\mathcal{X}'} X' \to X'$가 매끄러운지만 확인하면 충분하다.
$Y'$이 아핀인 에탈 사상
$$
Y' \to U'_w \times_{\mathcal{X}'} X'
$$
을 택하자. $U'_w \times_{\mathcal{X}'} X'$는 이러한 사상들의 상으로
덮이므로, $f'_1, \ldots, f'_r$가 잘라내는 $Y'$의 닫힌 부분스킴
$Z'$가 $X'$ 위에서 매끄럽다는 것을 보이면 충분하다. 다음과 같다.
$$
\xymatrix{
Z' \ar[r] \ar[d] & Y' \ar[d] \\
W'_w \times_{\mathcal{X}'} X' \ar[d] \ar[r] &
U'_w \times_{\mathcal{X}'} X' \ar[d] \ar[r] & X' \\
W'_w = V(f'_1, \ldots, f'_r) \ar[r] & U'_w
}
$$
$X = \mathcal{X} \times_{\mathcal{X}'} X'$,
$Y = X \times_{X'} Y' = \mathcal{X} \times_{\mathcal{X}'} Y'$, 그리고
$Z = Y \times_{Y'} Z' = X \times_{X'} Z' =
\mathcal{X} \times_{\mathcal{X}'} Z'$로 두자.
그러면
$(Z \subset Z') \to (Y \subset Y') \subset (X \subset X')$는
아핀 스킴의 두껍게 하기의 (카르테시안) 사상들이고,
$Z \to X$와 $Y' \to X'$가 매끄럽다는 것이 주어져 있다.
마지막으로 $Y' \to U'_w$가 매끄러워서 평탄하므로,
『대수학에 관한 추가 내용』의 보조정리
\ref{more-algebra-lemma-koszul-regular-flat-base-change}에 의해 함수 수열
$f'_1, \ldots, f'_r$의 $\Gamma(Y', \mathcal{O}_{Y'})$에서의 상은
코쥘-정칙 수열이다. 『대수학에 관한 추가 내용』의 보조정리
\ref{more-algebra-lemma-cut-by-koszul}에 의해(그리고
『대수학에 관한 추가 내용』의 보조정리
\ref{more-algebra-lemma-regular-koszul-regular},
\ref{more-algebra-lemma-koszul-regular-H1-regular},
\ref{more-algebra-lemma-H1-regular-quasi-regular}에 따라
코쥘-정칙 수열은 준정칙 수열이라는 사실에 의해), 원하는 대로
$Z' \to X'$가 매끄럽다고 결론짓는다.
\end{proof}

\begin{lemma}
\label{stacks-more-morphisms-lemma-etale-local-lifts-isomorphic}
$\mathcal{X} \subset \mathcal{X}'$를 대수적 스택의 두껍게 하기라 하자.
정사각형들이 카르테시안인 다음 가환 도식을 생각하자.
$$
\xymatrix{
W'' \ar[d]_{x''} & W \ar[l] \ar[r] \ar[d]_x & W' \ar[d]^{x'} \\
\mathcal{X}' & \mathcal{X} \ar[l] \ar[r] & \mathcal{X}'
}
$$
여기서 $W', W, W''$는 대수공간이고 세로 화살표들은 매끄럽다.
그러면 다음이 존재한다.
\begin{enumerate}
\item 에탈 덮개 $\{f'_k : W'_k \to W'\}_{k \in K}$,
\item 에탈 사상 $f''_k : W'_k \to W''$,
\item $2$-사상 $\gamma_k : x'' \circ f''_k \to x' \circ f'_k$.
\end{enumerate}
이들은 (a) $(f'_k)^{-1}(W) = (f''_k)^{-1}(W)$,
(b) $f'_k|_{(f'_k)^{-1}(W)} = f''_k|_{(f''_k)^{-1}(W)}$,
(c) $\gamma_k$를 (a)의 닫힌 부분스킴으로 당기면 처음 도식의
$W$ 위에서의 가환성이 주는 $2$-사상과 일치한다는 조건을 만족한다.
\end{lemma}

\begin{proof}
주어진 두껍게 하기를 $i : W \to W'$와 $i'' : W \to W''$라 쓰자.
명제의 도식이 가환이라는 것은
$2$-사상 $\delta : x' \circ i' \to x'' \circ i''$가 존재한다는 뜻이다.
이것이 명제의 (c)에서 말한 $2$-사상이다. 사영
$p' : I' \to W'$와 $q' : I' \to W''$를 갖는 대수공간
$$
I' = W' \times_{x', \mathcal{X}', x''} W''
$$
을 생각하자. ``보편'' $2$-사상
$\gamma : x' \circ p' \to x'' \circ q'$가 있음에 유의하자
(나중에 이를 사용할 것이다). $\delta$의 선택은 다음과 같은 사상을 정한다.
$$
\xymatrix{
W \ar[rr]_\delta & & I' \ar[ld]^{p'} \ar[rd]_{q'} \\
& W' & & W''
}
$$
이때 합성 $W \to I' \to W'$와 $W \to I' \to W''$는 각각
$i : W \to W'$와 $i' : W \to W''$이다. $x''$가 매끄러우므로
$p' : I' \to W'$는 $x''$의 밑변환으로서 매끄럽다.

\medskip\noindent
에탈 덮개 $\{f'_k : W'_k \to W'\}$와 사상
$\delta_k : W'_k \to I'$를 찾을 수 있고, $\delta_k$를
$W_k = (f'_k)^{-1}$에 제한하면 $\delta \circ f_k$와 같다고 가정하자.
여기서 $f_k = f'_k|_{W_k}$이다. 다음과 같다.
$$
\xymatrix{
W_k \ar[r]^{f_k} \ar[d] & W \ar[r]^\delta & I' \ar[d]^{p'} \\
W'_k \ar[rr]^{f'_k} \ar[rru]^{\delta_k} & & W'
}
$$
달리 말하면, 필요하다면 $W'$를 에탈 덮개로 바꾼 뒤
$p'$의 주어진 단면 $\delta : W \to I'$를 $W'$ 위의 단면으로
확장할 수 있기를 바란다.

\medskip\noindent
이것이 가능하다면 $f''_k = q' \circ \delta_k$와
$\gamma_k = \gamma \star \text{id}_{\delta_k}$로 둘 수 있다
(더 간단히 쓰면 $\gamma_k = \delta_k^*\gamma$이다).
실제로 이 단계에서 보일 것은 사상 $f''_k$가 에탈이라는 것뿐이다.
구성에 의해 사상 $x' \circ p'$는 $x'' \circ q'$와 $2$-동형이다.
따라서 $x'' \circ f''_k$는 $x' \circ f'_k$와 $2$-동형이다.
$x' \circ f'_k$가 매끄러우므로 합성
$$
W'_k \xrightarrow{f''_k} W'' \xrightarrow{x''} \mathcal{X}'
$$
은 매끄럽다고 결론짓는다. $f_k$가 에탈이므로 보조정리
\ref{stacks-more-morphisms-lemma-deform-property-fp-over-ft}에 의해 $f''_k$는 에탈이다.

\medskip\noindent
두껍게 하기가 일차 두껍게 하기이면, $W_k'$가 아핀인 임의의 에탈 덮개
$\{W'_k \to W'\}$를 택할 수 있다. 실제로 $p'$가 매끄러우므로
무한소 올림 판정법에 의해 $p'$는 형식적으로 매끄럽다
(『대수공간의 사상에 관한 추가 내용』의 보조정리
\ref{spaces-more-morphisms-lemma-smooth-formally-smooth}).
$W_k$가 아핀이고 $W_k \to W'_k$가 일차 두껍게 하기이므로
($\mathcal{X} \to \mathcal{X}'$의 밑변환이다. 보조정리
\ref{stacks-more-morphisms-lemma-base-change-thickening}을 보라), 원하는 $\delta_k$를 얻는다.

\medskip\noindent
일반적인 경우에는 덮개와 사상 $\delta_k$의 존재가
『대수공간의 사상에 관한 추가 내용』의 보조정리
\ref{spaces-more-morphisms-lemma-smooth-strong-lift}에서 따른다.
\end{proof}

\begin{lemma}
\label{stacks-more-morphisms-lemma-gerbe-of-lifts}
주석 \ref{stacks-more-morphisms-remark-gerbe-of-lifts}에서 구성한 범주
$p : \mathcal{C} \to W_{spaces, \etale}$는 gerbe이다.
\end{lemma}

\begin{proof}
보조정리 \ref{stacks-more-morphisms-lemma-gerbe-of-lifts-stack}에서 이것이 준군 스택임을 보았다.
따라서 『스택』의 정의 \ref{stacks-definition-gerbe}의 조건 (2)와 (3)을
확인하는 일만 남는다. 조건 (2)는 보조정리
\ref{stacks-more-morphisms-lemma-etale-local-lifts}에서 따르고, 조건 (3)은 보조정리
\ref{stacks-more-morphisms-lemma-etale-local-lifts-isomorphic}에서 따른다.
\end{proof}

\begin{lemma}
\label{stacks-more-morphisms-lemma-gerbe-of-lifts-first-order}
주석 \ref{stacks-more-morphisms-remark-gerbe-of-lifts}에서
$\mathcal{X} \subset \mathcal{X}'$가 일차 두껍게 하기라고 가정하자.
그러면
\begin{enumerate}
\item 주석 \ref{stacks-more-morphisms-remark-gerbe-of-lifts}에서 구성한 gerbe
$p : \mathcal{C} \to W_{spaces, \etale}$의 대상들의 자기동형 층은
아벨 층이다.
\item 『스택』의 보조정리
\ref{stacks-lemma-gerbe-abelian-auts}에서 구성한 군층
$\mathcal{G}$는 준연접 $\mathcal{O}_W$-가군이다.
\end{enumerate}
\end{lemma}

\begin{proof}
두 명제를 동시에 증명한다. 즉, 대상
$\xi = (U, U', a, i, x', \alpha)$가 주어지면 $\mathit{Aut}(\xi)$에
$U_{spaces, \etale}$ 위의 준연접 $\mathcal{O}_U$-가군 구조를 주고,
이 구조가 당김과 양립함을 보일 것이다. 층의 붙이기
(『사이트』의 절 \ref{sites-section-glueing-sheaves}),
『스택』의 보조정리 \ref{stacks-lemma-gerbe-abelian-auts}의 증명에서
자기동형 층 $\mathit{Aut}(\xi)$들을 붙여 $\mathcal{G}$를 구성하는 법,
그리고 가군이 준연접인지 여부는 에탈 덮개로 간 뒤 확인해도 충분하다는 사실
(『대수공간의 성질』의 보조정리
\ref{spaces-properties-lemma-characterize-quasi-coherent})에 의해
이것으로 충분하다.

\medskip\noindent
보조정리 \ref{stacks-more-morphisms-lemma-etale-local-lifts-isomorphic}의 증명에서 사용한 것과
같은 방법으로 층 $\mathit{Aut}(\xi)$를 서술한다. 사영
$p' : I' \to U'$와 $q' : I' \to U'$를 갖는 대수공간
$$
I' = U' \times_{x', \mathcal{X}', x'} U'
$$
을 생각하자. $I'$ 위에는 보편 $2$-사상
$\gamma : x' \circ p' \to x' \circ q'$가 있다.
항등사상 $x' \to x'$는 다음 대각사상을 정한다.
$$
\xymatrix{
U' \ar[rr]_{\Delta'} & & I' \ar[ld]^{p'} \ar[rd]_{q'} \\
& U' & & U'
}
$$
이때 두 합성 $U' \to I' \to U'$와 $U' \to I' \to U'$는
모두 항등사상이다.
$U', I', p', q', \Delta'$를 $\mathcal{X}$로 밑변환한 것을 각각
$U, I, p, q, \Delta$라 쓰겠다. $W' \to \mathcal{X}'$가 매끄러우므로
$p' : I' \to U'$는 밑변환으로서 매끄럽다.

\medskip\noindent
$U$ 위의 $\mathit{Aut}(\xi)$의 한 단면은
$\delta'|_U = \Delta$이고 $p' \circ \delta' = \text{id}_{U'}$인 사상
$\delta' : U' \to I'$이다. 명시적으로, 대응하는 자기동형사상은
$(\text{id}_U, q' \circ \delta', (\delta')^*\gamma) : \xi \to \xi$이다.
더 일반적으로 $f : V \to U$가 에탈 사상이면, 제한하면 $V$ 위에서
$f$가 되는 두껍게 하기 $j : V \to V'$와 에탈 사상
$f' : V' \to U'$가 존재하며, $f^*\xi$는
$(V, V', a \circ f, j, x' \circ f', f^*\alpha)$에 대응한다.
보조정리 \ref{stacks-more-morphisms-lemma-gerbe-of-lifts-fibred}의 증명을 보라.
$V$ 위의 $\mathit{Aut}(\xi)$의 한 단면은
$\delta'|_V = \Delta \circ f$이고 $p' \circ \delta' = f'$인 사상
$\delta' : V' \to I'$이다\footnote{대응하는 자기동형사상의 공식은
$(\text{id}_V, h', (\delta')^*\gamma)$이다. 여기서
$h' : V' \to V'$는 $h'|_V = \text{id}_V$이고
$$
\xymatrix{
V' \ar[r]_{h'} \ar[rd]_{q' \circ \delta'} & V' \ar[d]^{f'} \\
& U'
}
$$
가 가환이 되게 하는 유일한 (동형)사상이다. $h'$의 유일성과 존재는
에탈 사이트의 위상적 불변성에서 따른다. 『대수공간의 사상에 관한 추가 내용』의
정리 \ref{spaces-more-morphisms-theorem-topological-invariance}를 보라.
대신 $\delta'' \circ j = \Delta_{V'/\mathcal{X}'}$이고
$\text{pr}_1 \circ \delta'' = \text{id}_{V'}$인 사상
$\delta'' : V' \to V' \times_{\mathcal{X}'} V'$를 살펴봐야 한다고
독자는 생각할 수도 있다. 이 방법도 괜찮다.
$V' \times_{\mathcal{X}'} V' \to I'$가 에탈이므로, 같은 위상적 불변성에
따라 $\delta''$를
$\delta' = (V' \times_{\mathcal{X}'} V' \to I') \circ \delta''$로 보내는 것은
두 사상 집합 사이의 전단사이다.}.

\medskip\noindent
따라서 집합의 층으로서 $\mathit{Aut}(\xi)$는, 두껍게 하기
$(U \subset U')$와 $(I \subset I')$를 각각
$\text{id}_{U'}$와 $p'$를 통해 $(U \subset U')$ 위의 것으로 보아
『대수공간의 사상에 관한 추가 내용』의 주석
\ref{spaces-more-morphisms-remark-another-special-case}에서 정의한 층과 일치한다.
대각사상 $\Delta'$는 이 층의 한 단면이고,
『대수공간의 사상에 관한 추가 내용』의 보조정리
\ref{spaces-more-morphisms-lemma-action-sheaf}를 이용해 이 단면에
작용시키면 $U_{spaces, \etale}$ 위의 동형사상
\begin{equation}
\label{stacks-more-morphisms-equation-isomorphism}
\SheafHom_{\mathcal{O}_U}(\Delta^*\Omega_{I/U}, \mathcal{C}_{U/U'})
\longrightarrow
\mathit{Aut}(\xi)
\end{equation}
을 얻는다. 확인할 것이 세 가지 남았다.
\begin{enumerate}
\item (\ref{stacks-more-morphisms-equation-isomorphism})의 구성이 에탈 국소화와 가환한다.
\item $\SheafHom_{\mathcal{O}_U}(\Delta^*\Omega_{I/U}, \mathcal{C}_{U/U'})$는 $U$ 위의 준연접 가군이다.
\item $\mathit{Aut}(\xi)$에서의 합성은 이 준연접 가군의 단면들의
덧셈에 대응한다.
\end{enumerate}
차례대로 확인하겠다.

\medskip\noindent
(1)을 보이려면 $f : V \to U$가 에탈일 때 $U$ 위의 $\xi$를 이용하여
구성한 (\ref{stacks-more-morphisms-equation-isomorphism})을 제한한 사상이,
$V$ 위에서 $\xi|_V$를 이용해 $V_{spaces, \etale}$ 위에 구성한
(\ref{stacks-more-morphisms-equation-isomorphism})의 사상
$$
\SheafHom_{\mathcal{O}_V}(
\Delta_V^*\Omega_{V \times_\mathcal{X} V/V}, \mathcal{C}_{V/V'}) \to
\mathit{Aut}(\xi|_V)
$$
과 같음을 보여야 한다. 이는 위의 각주에 나온 논의와
『대수공간의 사상에 관한 추가 내용』의 보조정리
\ref{spaces-more-morphisms-lemma-action-by-derivations-etale-localization}에서 따른다.

\medskip\noindent
(2)의 증명. $p'$가 매끄러우므로 사상 $I \to U$는 매끄럽고,
따라서 상대 미분 가군 $\Omega_{I/U}$는 유한 국소 자유이다
(『대수공간의 사상에 관한 추가 내용』의 보조정리
\ref{spaces-more-morphisms-lemma-smooth-omega-finite-locally-free}).
한편 $\mathcal{C}_{U/U'}$는 준연접이다
(『대수공간의 사상에 관한 추가 내용』의 정의
\ref{spaces-more-morphisms-definition-conormal-sheaf}).
그러므로 『대수공간의 성질』의 보조정리
\ref{spaces-properties-lemma-properties-quasi-coherent}에 의해 결론을 얻는다.

\medskip\noindent
(3)의 증명. $(U', I', p', q', c')$가 항등원 $\Delta'$를 갖는
대수공간의 준군이 되게 하는 사상
$c' : I' \times_{p', U', q'} I' \to I'$가 존재한다.
예를 들어 『대수적 스택』의 보조정리
\ref{algebraic-lemma-map-space-into-stack}를 보라.
$\mathit{Aut}(\xi)$에서의 합성은 다음과 같이 사상 $c'$에 의해 유도된다.
위에서처럼 $U$ 위의 $\mathit{Aut}(\xi)$의 단면에 대응하는 두 사상
$$
\delta'_1, \delta'_2 : U' \longrightarrow I'
$$
이 주어졌다고 하자. 달리 말하면 $\delta'_i|U = \Delta_U$이고
$p' \circ \delta'_i = \text{id}_{U'}$이다. 그러면
$\mathit{Aut}(\xi)$에서의 합성은
$$
\delta'_1 \circ \delta'_2 = c'(\delta'_1 \circ q' \circ \delta'_2, \delta'_2)
$$
이다. 자세한 확인은 생략한다\footnote{섬유곱
$I' \times_{p', U', q'} I'$의 잘 정의된 $U'$-값 점을 얻으려면
$\delta'_1$에 $q' \circ \delta'_2$를 미리 합성해야 함을 독자는
즉시 알 수 있다.}. 따라서 우리는
『대수공간의 준군에 관한 추가 내용』의 절
\ref{spaces-more-groupoids-section-groupoid-sections}에 서술된 상황에 있고,
원하는 결과는 『대수공간의 준군에 관한 추가 내용』의 보조정리
\ref{spaces-more-groupoids-lemma-composition-is-addition}에서 따른다.
\end{proof}

\begin{proposition}[Emerton]
\label{stacks-more-morphisms-proposition-affine-smooth-lift-to-first-order}
\begin{reference}
Matthew Emerton이 2016년 4월 27일에 보낸 이메일.
\end{reference}
$\mathcal{X} \subset \mathcal{X}'$를 대수적 스택의
일차 두껍게 하기라 하자. $W$를 아핀 스킴이라 하고
$W \to \mathcal{X}$를 매끄러운 사상이라 하자. 그러면
$W' \to \mathcal{X}'$가 매끄럽고 $W'$가 아핀인 카르테시안 도식
$$
\xymatrix{
W \ar[d] \ar[r] & W' \ar[d] \\
\mathcal{X} \ar[r] & \mathcal{X}'
}
$$
이 존재한다.
\end{proposition}

\begin{proof}
주석 \ref{stacks-more-morphisms-remark-gerbe-of-lifts}에서 도입한 범주
$p : \mathcal{C} \to W_{spaces, \etale}$를 생각하자.
이 명제는 $W$ 위에 놓이는 $\mathcal{C}$의 대상이 존재한다는 말이다.
실제로 그러한 대상 $(W, W', a, i, y', \alpha)$가 있으면
$W = \mathcal{X} \times_{\mathcal{X}'} W'$이다. 따라서
$W \to W'$는 대수공간의 두껍게 하기이므로
『대수공간의 사상에 관한 추가 내용』의 보조정리
\ref{spaces-more-morphisms-lemma-thickening-scheme}와
『사상에 관한 추가 내용』의 보조정리
\ref{more-morphisms-lemma-thickening-affine-scheme}에 의해 $W'$는 아핀이다.

\medskip\noindent
보조정리 \ref{stacks-more-morphisms-lemma-gerbe-of-lifts}에 따르면 $\mathcal{C}$는
$W_{spaces, \etale}$ 위의 gerbe이다. 이는 에탈 국소적으로 해를 찾을 수
있고, 이 국소 해들이 에탈 국소적으로 동형이라는 뜻이다.
이 부분에는 두껍게 하기가 일차라는 가정이 필요하지 않다.
보조정리 \ref{stacks-more-morphisms-lemma-gerbe-of-lifts-first-order}에 의해 우리 gerbe의
대상들의 자기동형 층은 아벨 층이고, 함께 맞물려
$W_{spaces, \etale}$ 위의 준연접 가군 $\mathcal{G}$를 이룬다.
$W$ 위에 놓이는 $\mathcal{C}$의 대상이 존재한다고 결론짓기 위해
『사이트 위의 코호몰로지』의 보조정리
\ref{sites-cohomology-lemma-existence}의 조건 (1)과 (2)를 확인하겠다.
조건 (1)은 참이다. 각 $W_i$가 아핀인 에탈 덮개
$\{W_i \to W\}$들은 모든 덮개의 모임에서 공종이다.
그러한 덮개에 대해 $W_i$와 $W_i \times_W W_j$는 아핀이고,
$H^1(W_i, \mathcal{G})$와
$H^1(W_i \times_W W_j, \mathcal{G})$는 영이다. 예를 들어
『대수공간의 코호몰로지』의 명제
\ref{spaces-cohomology-proposition-vanishing}에 따라 아핀 대수공간 위의
준연접 가군의 코호몰로지는 영이기 때문이다. 마지막으로 조건 (2)는
우리 준연접 층 $\mathcal{G}$에 대해 $H^2(W, \mathcal{G}) = 0$이라는 것인데,
이 역시 『대수공간의 코호몰로지』의 명제
\ref{spaces-cohomology-proposition-vanishing}에서 따른다.
이로써 증명이 끝난다.
\end{proof}

\section{무한소 변형}
\label{stacks-more-morphisms-section-inf}

\noindent
『Artin의 공리』의 절 \ref{artin-section-inf}에서 시작한 논의를 계속한다.

\begin{lemma}
\label{stacks-more-morphisms-lemma-inf-quasi-coherent}
$\mathcal{X}$를 스킴 $S$ 위의 대수적 스택이라 하자.
$\mathcal{I}_\mathcal{X} \to \mathcal{X}$가 국소 유한 표시라고 가정하자.
$A \to B$를 평탄한 $S$-대수 준동형이라 하자.
$x$를 $A$ 위의 $\mathcal{X}$의 대상이라 하고 $y = x|_B$로 두자.
그러면
$\text{Inf}_x(M) \otimes_A B = \text{Inf}_y(M \otimes_A B)$이다.
\end{lemma}

\begin{proof}
$\text{Inf}_x(M)$은 $x$를 $A[M]$으로 자명하게 변형한 것의
자기동형사상 가운데 $A$ 위의 $x$의 항등 자기동형사상을 유도하는
것들의 집합임을 상기하자. 자명한 변형은
$\Spec(A[M]) \to \Spec(A)$를 통한 $x$의 $\Spec(A[M])$으로의 당김이다.
$G \to \Spec(A)$를 $x$의 자기동형 군 대수공간이라 하자
($\mathcal{X}$가 대수공간이므로 이것이 존재한다).
$e : \Spec(A) \to G$를 중립원이라 하자.
『대수공간의 사상에 관한 추가 내용』의 절
\ref{spaces-more-morphisms-section-action-by-derivations}의 논의에 따르면
$$
\text{Inf}_x(M) = \Hom_A(e^*\Omega_{G/A}, M)
$$
이다. 마찬가지로
$$
\text{Inf}_y(M \otimes_A B) = \Hom_B(e_B^*\Omega_{G_B/B}, M \otimes_A B)
$$
이다. 가정에 의해 $G \to \Spec(A)$가 국소 유한 표시이므로
$\Omega_{G/A}$는 국소 유한 표시이다. 『대수공간의 사상에 관한 추가 내용』의
보조정리 \ref{spaces-more-morphisms-lemma-finite-presentation-differentials}를
보라. 따라서 $e^*\Omega_{G/A}$는 유한 표시 $A$-가군이다.
더욱이 『대수공간의 사상에 관한 추가 내용』의 보조정리
\ref{spaces-more-morphisms-lemma-base-change-differentials}에 의해
$\Omega_{G_B/B}$는 $\Omega_{G/A}$의 당김이다. 그러므로
$e_B^*\Omega_{G_B/B} = e^*\Omega_{G/A} \otimes_A B$이다.
『대수학에 관한 추가 내용』의 보조정리
\ref{more-algebra-lemma-pseudo-coherence-and-base-change-ext}로 결론을 얻는다.
\end{proof}

\begin{lemma}
\label{stacks-more-morphisms-lemma-sheaf-of-infinitesimal-lifts}
$\mathcal{X}$를 밑스킴 $S$ 위의 대수적 스택이라 하자.
$\mathcal{I}_\mathcal{X} \to \mathcal{X}$가 국소 유한 표시라고 가정하자.
$(A' \to A, x)$를 변형 상황이라 하자. 그러면 함자
$$
F : B' \longmapsto
\{\text{올림 }x|_{B' \otimes_{A'} A}\text{ 에서 } B'\}/\text{동형사상}
$$
는 『위상』의 정의 \ref{topologies-definition-big-small-fppf}에 나오는
사이트 $(\textit{Aff}/\Spec(A'))_{fppf}$ 위의 층이다.
\end{lemma}

\begin{proof}
$\{T'_i \to T'\}_{i = 1, \ldots n}$을 $A'$ 위의 아핀 스킴의
표준 fppf 덮개라 하고 $T' = \Spec(B')$로 쓰자. 평소와 같이
$$
T'_{i_0 \ldots i_p} =
T'_{i_0} \times_{T'} \ldots \times_{T'} T'_{i_p} = \Spec(B'_{i_0 \ldots i_p})
$$
로 쓰며, 여기서 환은 적절한 텐서곱이다.
$B = B' \otimes_{A'} A$와
$B_{i_0 \ldots i_p} = B'_{i_0 \ldots i_p} \otimes_{A'} A$로 두자.
$y = x|_B$와 $y_{i_0 \ldots i_p} = x|_{B_{i_0 \ldots i_p}}$로 쓰자.
$\gamma_i \in F(B'_i)$들이 주어지고 $\gamma_{i_0}$와
$\gamma_{i_1}$의 $F(B'_{i_0i_1})$에서의 상이 같다고 하자.
$F(B'_i)$에서 $\gamma_i$로 가는 유일한
$\gamma \in F(B')$를 찾아야 한다.

\medskip\noindent
동형류 $\gamma_i$ 안에서 $\textit{Lift}(y_i, B'_i)$의 실제 대상
$y'_i$를 택한다. 범주 $\textit{Lift}(y_{i_0i_1}, B'_{i_0i_1})$ 안의
동형사상
$\varphi_{i_0i_1} : y'_{i_0}|_{B'_{i_0i_1}} \to y'_{i_1}|_{B'_{i_0i_1}}$를 택한다.
사상 $\varphi_{i_0i_1}$들이 코사이클 조건을 만족하면,
$\mathcal{X}$가 fppf 위상에서 스택이므로 원하는 대상 $\gamma$를 얻는다.
코사이클 조건은 합성
$$
y'_{i_0}|_{B'_{i_0i_1i_2}}
\xrightarrow{\varphi_{i_0i_1}|_{B'_{i_0i_1i_2}}}
y'_{i_1}|_{B'_{i_0i_1i_2}}
\xrightarrow{\varphi_{i_1i_2}|_{B'_{i_0i_1i_2}}}
y'_{i_2}|_{B'_{i_0i_1i_2}}
\xrightarrow{\varphi_{i_2i_0}|_{B'_{i_0i_1i_2}}}
y'_{i_0}|_{B'_{i_0i_1i_2}}
$$
이 항등사상이라는 것이다. 그렇지 않다면 이 사상들은 원소
$$
\delta_{i_0i_1i_2} \in
\text{Inf}_{y_{i_0i_1i_2}}(J_{i_0i_1i_2}) =
\text{Inf}_y(J) \otimes_B B_{i_0i_1i_2}
$$
를 준다. 여기서 $J = \Ker(B' \to B)$이고
$J_{i_0 \ldots i_p} = \Ker(B'_{i_0 \ldots i_p} \to B_{i_0 \ldots i_p})$이다.
표시된 등식은 $B' \to B'_{i_0 \ldots i_p}$와 $y$ 및
$y_{i_0 \ldots i_p}$에 보조정리 \ref{stacks-more-morphisms-lemma-inf-quasi-coherent}를
적용하여 얻는다. 여기에는 사상 $B' \to B'_{i_0 \ldots i_p}$의
평탄성도 사용하며, 이 평탄성은
$J_{i_0 \ldots i_p} = J \otimes_{B'} B'_{i_0 \ldots i_p}$도 보장한다.
계산(생략함)에 따르면 $\delta_{i_0i_1i_2}$는 {\v C}ech 복합체
$$
\prod \text{Inf}_y(J) \otimes_B B_{i_0} \to
\prod \text{Inf}_y(J) \otimes_B B_{i_0i_1} \to
\prod \text{Inf}_y(J) \otimes_B B_{i_0i_1i_2} \to \ldots
$$
의 $2$-코사이클을 준다. 『내림』의 보조정리
\ref{descent-lemma-standard-covering-Cech-quasi-coherent}에 의해
이 복합체는 양의 차수에서 비순환이고
$H^0 = \text{Inf}_y(J)$이다.
$\text{Inf}_{y_{i_0i_1}}(J_{i_0i_1})$가 사상들에 작용하므로
(『Artin의 공리』의 주석 \ref{artin-remark-automorphisms}),
$\delta_{i_0i_1i_2} = 0$인 경우가 되도록
$\varphi_{i_0i_1}$의 선택을 수정할 수 있다.

\medskip\noindent
유일성. 모든 $i$에 대해 $\gamma_i$로 제한되는 $\gamma$가 많아야
하나임을 여전히 보여야 한다. $\textit{Lift}(y, B')$의 두 대상
$y', z'$와 $\textit{Lift}(y_i, B'_i)$ 안의 동형사상
$\psi_i : y'|_{B'_i} \to z'|_{B'_i}$가 주어졌다고 하자. 그러면
$$
\psi_{i_1}^{-1} \circ \psi_{i_0} \in
\text{Inf}_{y_{i_0i_1}}(J_{i_0i_1}) =
\text{Inf}_y(J) \otimes_B B_{i_0i_1}
$$
을 생각할 수 있다. 앞에서와 같이 논증하면 $B'$ 위에서 $y'$와
$z'$ 사이의 동형사상을 찾는 데 대한 장애물은 위에 표시한 {\v C}ech 복합체의
$H^1$의 원소이며, 이는 영이다.
\end{proof}

\begin{lemma}
\label{stacks-more-morphisms-lemma-T-quasi-coherent}
$\mathcal{X}$를 스킴 $S$ 위의 대수적 스택이라 하고, 그 구조 사상
$\mathcal{X} \to S$가 국소 유한 표시라고 하자.
$A \to B$를 평탄한 $S$-대수 준동형이라 하자.
$x$를 $A$ 위의 $\mathcal{X}$의 대상이라 하자.
그러면 $T_x(M) \otimes_A B = T_y(M \otimes_A B)$이다.
\end{lemma}

\begin{proof}
스킴 $U$와 전사 매끄러운 사상 $U \to \mathcal{X}$를 택한다.
먼저 $x$가 $U$로 올라가는 경우로 보조정리를 환원한다.
$T_x(M)$은 $x$를 $A[M]$으로 올린 것들의 동형류 집합임을 상기하자.
따라서 보조정리 \ref{stacks-more-morphisms-lemma-sheaf-of-infinitesimal-lifts}\footnote{이 보조정리는
적용할 수 있다. $\Delta : \mathcal{X} \to \mathcal{X} \times_S \mathcal{X}$는
『스택의 사상』의 보조정리
\ref{stacks-morphisms-lemma-finite-presentation-permanence}와
$\mathcal{X} \to S$가 국소 유한 표시라는 가정에 의해 국소 유한 표시이다.
따라서 $\mathcal{I}_\mathcal{X} \to \mathcal{X}$는 $\Delta$의
밑변환으로서 국소 유한 표시이다.}에 따르면 규칙
$$
A_1 \mapsto T_{x|_{A_1}}(M \otimes_A A_1)
$$
은 $\Spec(A)$의 작은 에탈 사이트 위의 층이다. 텐서곱은
$A[M] \to A_1[M \otimes_A A_1]$을 평탄한 환 사상으로 만들기 위해 필요하다.
$x|_{A_1}$이 사상 $u_1 : \Spec(A_1) \to U$로 올라가게 하는
충실히 평탄한 에탈 환 사상 $A \to A_1$을 택할 수 있다.
예를 들어 『스택 위의 층』의 보조정리
\ref{stacks-sheaves-lemma-surjective-flat-locally-finite-presentation}를 보라.
$A_2 = A_1 \otimes_A A_1$로 쓰고
$B_1 = B \otimes_A A_1$, $B_2 = B \otimes_A A_2$로 두자. 도식
$$
\xymatrix{
0 \ar[r] &
T_y(M \otimes_A B) \ar[r] &
T_{y|_{B_1}}(M \otimes_A B_1) \ar[r] &
T_{y|_{B_2}}(M \otimes_A B_2) \\
0 \ar[r] &
T_x(M) \ar[r] \ar[u] &
T_{x|_{A_1}}(M \otimes_A A_1) \ar[r] \ar[u] &
T_{x|_{A_2}}(M \otimes_A A_2) \ar[u]
}
$$
을 생각하자. 층 조건에 의해 두 행은 완전하다.
$M \otimes_A B_i = (M \otimes_A A_i) \otimes_{A_i} B_i$이다.
따라서 가운데와 오른쪽 세로 화살표에 대해 결과를 증명하면
원하는 결과가 따른다. 이로써 다음 문단의 경우로 환원된다.

\medskip\noindent
$x$가 사상 $u : \Spec(A) \to U$의 상이라고 가정하자.
$U \to \mathcal{X}$가 매끄럽고 대수공간으로 표현가능하므로
$T_u(M) \to T_x(M)$은 전사이다. 『표현가능성의 판정 조건』의 보조정리
\ref{criteria-lemma-representable-by-spaces-formally-smooth}
(설명은 그 앞의 논의를 보라)와 『대수공간의 사상에 관한 추가 내용』의
보조정리 \ref{spaces-more-morphisms-lemma-smooth-formally-smooth}를 보라.
$R = U \times_\mathcal{X} U$로 두자. 대수공간의 준군
$(U, R, s, t, c, e, i)$를 얻고 $\mathcal{X} = [U/R]$임을 상기하자.
『Artin의 공리』의 보조정리 \ref{artin-lemma-ses-inf-and-T}에 의해
완전열
$$
T_{e \circ u}(M) \to T_u(M) \oplus T_u(M) \to T_x(M) \to 0
$$
을 얻으며, 오른쪽의 영은 위에서 보였다. $B$로 밑변환해도
비슷한 열이 성립한다. 따라서 $T_u(M)$과 $T_{e \circ u}(M)$에 대해
보조정리의 결과를 증명할 수 있으면 원하는 결과가 따른다.
이로써 다음 문단의 경우로 환원된다.

\medskip\noindent
$\mathcal{X} = X$가 $S$ 위에서 국소 유한 표시인 대수공간이라고 가정하자.
그러면 『대수공간의 사상에 관한 추가 내용』의 절
\ref{spaces-more-morphisms-section-action-by-derivations}의 논의에 의해
$$
T_x(M) = \Hom_A(x^*\Omega_{X/S}, M)
$$
이다. 마찬가지로
$$
T_y(M \otimes_A B) = \Hom_B(y^*\Omega_{X/S}, M \otimes_A B)
$$
이다. $X \to S$가 국소 유한 표시이므로 $\Omega_{X/S}$는
국소 유한 표시이다. 『대수공간의 사상에 관한 추가 내용』의 보조정리
\ref{spaces-more-morphisms-lemma-finite-presentation-differentials}를 보라.
따라서 $x^*\Omega_{X/S}$는 유한 표시 $A$-가군이다.
명백히 $y^*\Omega_{X/S} = x^*\Omega_{X/S} \otimes_A B$이다.
『대수학에 관한 추가 내용』의 보조정리
\ref{more-algebra-lemma-pseudo-coherence-and-base-change-ext}로 결론을 얻는다.
\end{proof}

\begin{lemma}
\label{stacks-more-morphisms-lemma-local-lift-enough}
$\mathcal{X}$를 스킴 $S$ 위의 대수적 스택이라 하고, 그 구조 사상
$\mathcal{X} \to S$가 국소 유한 표시라고 하자.
$(A' \to A, x)$를 변형 상황이라 하자. 충실히 평탄하고 유한 표시인
$A'$-대수 $B'$와 $x|_{B' \otimes_{A'} A}$를 올리는
$B'$ 위의 $\mathcal{X}$의 대상 $y'$가 존재한다면,
$x$를 올리는 $A'$ 위의 대상 $x'$가 존재한다.
\end{lemma}

\begin{proof}
$I = \Ker(A' \to A)$라 하자.
$B'_1 = B' \otimes_{A'} B'$와
$B'_2 = B' \otimes_{A'} B' \otimes_{A'} B'$로 두자.
$J = IB'$, $J_1 = IB'_1$, $J_2 = IB'_2$라 하고
$B = B'/J$, $B_1 = B'_1/J_1$, $B_2 = B'_2/J_2$라 하자.
$y = x|_B$, $y_1 = x|_{B_1}$, $y_2 = x|_{B_2}$로 두자.
$F$를 보조정리 \ref{stacks-more-morphisms-lemma-sheaf-of-infinitesimal-lifts}의 fppf 층이라 하자
(이 보조정리를 적용할 수 있음은 보조정리
\ref{stacks-more-morphisms-lemma-T-quasi-coherent}의 증명의 각주를 보라).
그러면 등화자 도식
$$
\xymatrix{
F(A') \ar[r] &
F(B') \ar@<1ex>[r] \ar@<-1ex>[r] &
F(B'_1)
}
$$
을 얻는다. 한편 『Artin의 공리』의 절 \ref{artin-section-inf}의
용어로 $F(B') = \text{Lift}(y, B')$,
$F(B'_1) = \text{Lift}(y_1, B'_1)$,
$F(B'_2) = \text{Lift}(y_2, B'_2)$이다.
이 집합들은 공집합이 아니며 각각 $T_y(J)$, $T_{y_1}(J_1)$,
$T_{y_2}(J_2)$의 (표준적인) 주동차 공간이다.
『Artin의 공리』의 보조정리
\ref{artin-lemma-properties-lift-RS-star}를 보라.
따라서 $F(B'_1)$에서 $y'$의 두 상의 차이는 원소
$$
\delta_1 \in T_{y_1}(J_1) = T_x(I) \otimes_A B_1
$$
이다. 표시된 등식은 $A' \to B'_1$과 $x$ 및 $y_1$에
보조정리 \ref{stacks-more-morphisms-lemma-T-quasi-coherent}를 적용하여 얻는다.
여기에는 $A' \to B'_1$의 평탄성도 사용하며, 이 평탄성은
$J_1 = I \otimes_{A'} B'_1$도 보장한다. $B'$와 $B'_2$에 대해서도
비슷한 등식들이 성립한다. 계산(생략함)에 따르면 $\delta_1$은
{\v C}ech 복합체
$$
T_x(I) \otimes_A B \to
T_x(I) \otimes_A B_1 \to
T_x(I) \otimes_A B_2 \to \ldots
$$
의 $1$-코사이클을 준다. 『내림』의 보조정리
\ref{descent-lemma-standard-covering-Cech-quasi-coherent}에 의해
이 복합체는 양의 차수에서 비순환이고 $H^0 = T_x(I)$이다.
따라서 경계가 $\delta_1$인 원소를
$T_x(I) \otimes_A B = T_y(J)$에서 택할 수 있다. $y'$를 이 원소가
작용한 결과로 바꾸면 $\delta_1 = 0$인 새로운 선택 $y'$를 얻는다.
따라서 $y'$는 두 사상 $F(B') \to F(B'_1)$ 아래에서 같은 원소로
가며, 층 조건에 의해 $F(A')$의 원소를 얻는다.
\end{proof}









\section{형식적으로 매끄러운 사상}
\label{stacks-more-morphisms-section-formally-smooth}

\noindent
이 절에서는 대수적 스택의 형식적으로 매끄러운 사상
$\mathcal{X} \to \mathcal{Y}$의 개념을 도입한다.
이러한 사상은 $T$가 아핀일 때 $\mathcal{X}$의 $T$-값 점이
$T$의 무한소 두껍게 하기로 올라간다는 성질로 특징지어진다.
주요 결과는 형식적으로 매끄럽고 국소 유한 표시인 사상이
매끄럽다는 것이다. 보조정리 \ref{stacks-more-morphisms-lemma-smooth-formally-smooth}를 보라.
이 판정법은 야코비 판정법보다 사용하기 쉬운 경우가 많다.

\begin{definition}
\label{stacks-more-morphisms-definition-formally-smooth}
대수적 스택의 사상 $f : \mathcal{X} \to \mathcal{Y}$가
{\it 형식적으로 매끄럽다}는 것은 『표현가능성의 판정 조건』의 절
\ref{criteria-section-formally-smooth}에서 설명한 대로,
준군 섬유화 범주들의 $1$-사상으로서 대상에 대해 형식적으로
매끄럽다는 뜻이다.
\end{definition}

\noindent
정의의 조건을 현재 사용하는 언어로 옮겨 쓰자
(『스택의 성질』의 절 \ref{stacks-properties-section-conventions}을 보라).
$f : \mathcal{X} \to \mathcal{Y}$를 대수적 스택의 사상이라 하자.
$i : T \to T'$가 아핀 스킴의 일차 두껍게 하기인 다음
$2$-가환 실선 도식을 생각하자.
\begin{equation}
\label{stacks-more-morphisms-equation-diagram}
\vcenter{
\xymatrix{
T \ar[r]_-x \ar[d]_i & \mathcal{X} \ar[d]^f \\
T' \ar[r]^-y \ar@{..>}[ru] & \mathcal{Y}
}
}
\end{equation}
도식의 $2$-가환성을 입증하는 $2$-사상을
$$
\gamma : y \circ i \longrightarrow f \circ x
$$
라 하자(『범주』의 절 \ref{categories-section-formal-cat-cat}와
\ref{categories-section-2-categories}의 표기법을 따른다).
(\ref{stacks-more-morphisms-equation-diagram})과 $\gamma$가 주어졌을 때,
{\it 점선 화살표}란 사상 $x' : T' \to \mathcal{X}$와 $2$-화살표
$\alpha : x' \circ i \to x$, $\beta : y \to f \circ x'$로 이루어진
삼중항 $(x', \alpha, \beta)$로서
$\gamma = (\text{id}_f \star \alpha) \circ (\beta \star \text{id}_i)$를 만족하는 것을 말한다. 달리 말하면 도식
$$
\xymatrix{
& f \circ x' \circ i \ar[rd]^{\text{id}_f \star \alpha} \\
y \circ i \ar[ru]^{\beta \star \text{id}_i} \ar[rr]^\gamma & &
f \circ x
}
$$
이 가환한다. {\it 점선 화살표의 사상}
$(x'_1, \alpha_1, \beta_1) \to (x'_2, \alpha_2, \beta_2)$란
$\alpha_1 = \alpha_2 \circ (\theta \star \text{id}_i)$와
$\beta_2 = (\text{id}_f \star \theta) \circ \beta_1$을 만족하는
$2$-화살표 $\theta : x'_1 \to x'_2$를 말한다.

\medskip\noindent
방금 서술한 점선 화살표의 범주는 『범주』의 정의
\ref{categories-definition-dotted-arrows}의 특수한 경우이다.

\begin{lemma}
\label{stacks-more-morphisms-lemma-reformulate-formal-smoothness}
대수적 스택의 사상 $f : \mathcal{X} \to \mathcal{Y}$가 형식적으로
매끄러울(정의 \ref{stacks-more-morphisms-definition-formally-smooth}) 필요충분조건은
모든 도식 (\ref{stacks-more-morphisms-equation-diagram})과 $\gamma$에 대해
점선 화살표의 범주가 공집합이 아닌 것이다.
\end{lemma}

\begin{proof}
서로 다른 언어 사이의 번역은 생략한다.
\end{proof}

\begin{lemma}
\label{stacks-more-morphisms-lemma-base-change-formal-smoothness}
대수적 스택의 형식적으로 매끄러운 사상을 대수적 스택의 임의의 사상으로
밑변환한 것은 형식적으로 매끄럽다.
\end{lemma}

\begin{proof}
『범주』의 보조정리
\ref{categories-lemma-cat-dotted-arrows-base-change}와 정의에서 따른다.
\end{proof}

\begin{lemma}
\label{stacks-more-morphisms-lemma-composition-formal-smoothness}
대수적 스택의 형식적으로 매끄러운 사상들의 합성은
형식적으로 매끄럽다.
\end{lemma}

\begin{proof}
『범주』의 보조정리
\ref{categories-lemma-cat-dotted-arrows-composition}과 정의에서 따른다.
\end{proof}

\begin{lemma}
\label{stacks-more-morphisms-lemma-formal-smoothness-representable}
$f : \mathcal{X} \to \mathcal{Y}$를 대수공간으로 표현가능한
대수적 스택의 사상이라 하자. 다음은 동치이다.
\begin{enumerate}
\item $f$는 형식적으로 매끄럽다.
\item 모든 스킴 $T$와 사상 $T \to \mathcal{Y}$에 대해 사상
$\mathcal{X} \times_\mathcal{Y} T \to T$가 대수공간의 사상으로서
형식적으로 매끄럽다.
\end{enumerate}
\end{lemma}

\begin{proof}
『범주』의 보조정리
\ref{categories-lemma-cat-dotted-arrows-base-change}와 정의에서 따른다.
\end{proof}

\begin{lemma}
\label{stacks-more-morphisms-lemma-lift-to-smooth}
$T \to T'$를 아핀 스킴의 일차 두껍게 하기라 하자.
$\mathcal{X}'$를 $T'$ 위의 대수적 스택이라 하고, 그 구조 사상
$\mathcal{X}' \to T'$가 매끄럽다고 하자.
$x : T \to \mathcal{X}'$를 $T'$ 위의 사상이라 하자.
그러면 $x'|_T = x$인 $T'$ 위의 사상
$x' : T' \to \mathcal{X}'$가 존재한다.
\end{lemma}

\begin{proof}
보조정리 \ref{stacks-more-morphisms-lemma-local-lift-enough}의 결과를 적용할 수 있다.
따라서 $W'$가 아핀이고 $x|_{T \times_{W'} T'}$가 $W'$로 올라가게 하는
매끄러운 전사 사상 $W' \to T'$를 구성하면 충분하다.
(대수공간에 대해 이미 확립한 비슷한 결과를 이용하여 이 사실을
직접 증명해 보기를 독자에게 권한다.) 스킴 $U'$와 전사 매끄러운 사상
$U' \to \mathcal{X}'$를 택한다. $U' \to T'$가 매끄럽고 사영
$T \times_{\mathcal{X}'} U' \to T$가 전사 매끄러움에 유의하자.
$W \to T$가 전사인 아핀 스킴 $W$와 에탈 사상
$W \to T \times_{\mathcal{X}'} U'$를 택한다. 그러면
$W \to T$는 아핀 스킴의 매끄러운 사상이다. $W$를 주 아핀 열린집합들의
서로소 합으로 바꾼 뒤, $W = T \times_{T'} W'$인 아핀들 사이의
매끄러운 사상 $W' \to T'$가 존재한다고 가정할 수 있다.
『대수학』의 보조정리 \ref{algebra-lemma-lift-smooth}를 보라.
『대수공간의 사상에 관한 추가 내용』의 보조정리
\ref{spaces-more-morphisms-lemma-smooth-formally-smooth}에 의해,
주어진 사상 $W \to U'$을 올리는 $T'$ 위의 사상 $W' \to U'$을
찾을 수 있다. 이로써 증명이 끝난다.
\end{proof}

\noindent
다음 보조정리가 이 절의 주요 결과이다. 이것을 『스택의 극한』의 명제
\ref{stacks-limits-proposition-characterize-locally-finite-presentation}와
합치면, 준군 스택의 $1$-사상 $\mathcal{X} \to \mathcal{Y}$의
``단순한'' 성질들을 이용하여 대수적 스택의 사상
$f : \mathcal{X} \to \mathcal{Y}$가 매끄러운지 판별할 수 있다.

\begin{lemma}[무한소 올림 판정법]
\label{stacks-more-morphisms-lemma-smooth-formally-smooth}
$f : \mathcal{X} \to \mathcal{Y}$를 대수적 스택의 사상이라 하자.
다음은 동치이다.
\begin{enumerate}
\item 사상 $f$는 매끄럽다.
\item 사상 $f$는 국소 유한 표시이고 형식적으로 매끄럽다.
\end{enumerate}
\end{lemma}

\begin{proof}
$f$가 매끄럽다고 가정하자. 그러면 『스택의 사상』의 보조정리
\ref{stacks-morphisms-lemma-smooth-locally-finite-presentation}에 의해
$f$는 국소 유한 표시이다. 따라서 도식 (\ref{stacks-more-morphisms-equation-diagram})과
$\gamma : y \circ i \to f \circ x$가 주어졌을 때 점선 화살표를 찾으면 충분하다
(보조정리 \ref{stacks-more-morphisms-lemma-reformulate-formal-smoothness}를 보라).
섬유곱을 만들면
$$
\xymatrix{
T \ar[d] \ar[r] &
T' \times_\mathcal{Y} \mathcal{X} \ar[d] \ar[r] &
\mathcal{X} \ar[d] \\
T' \ar[r] & T' \ar[r] & \mathcal{Y}
}
$$
를 얻는다. 따라서 왼쪽 정사각형에서 점선 화살표를 찾으면 충분하다.
$T' \times_\mathcal{Y} \mathcal{X} \to T'$가 매끄러우므로
(『스택의 사상』의 보조정리
\ref{stacks-morphisms-lemma-base-change-smooth}), 왼쪽 정사각형의
점선 화살표의 존재는 보조정리 \ref{stacks-more-morphisms-lemma-lift-to-smooth}로 보장된다.

\medskip\noindent
반대로 $f$가 국소 유한 표시이고 형식적으로 매끄럽다고 하자.
스킴 $U$와 전사 매끄러운 사상 $U \to \mathcal{X}$를 택한다.
그러면 $a : U \to \mathcal{X}$와 $b : U \to \mathcal{Y}$는
대수공간으로 표현가능하고 국소 유한 표시이다
(『스택의 사상』의 보조정리
\ref{stacks-morphisms-lemma-composition-finite-presentation}와,
위에서 본 매끄러운 사상은 국소 유한 표시라는 사실을 사용한다).
『대수적 스택』의 보조정리
\ref{algebraic-lemma-representable-transformations-property-implication}의
일반 원리를 적용하되, 『대수공간의 사상에 관한 추가 내용』의 보조정리
\ref{spaces-more-morphisms-lemma-smooth-formally-smooth}의 동치를
입력으로 삼고, 동시에 『표현가능성의 판정 조건』의 보조정리
\ref{criteria-lemma-representable-by-spaces-formally-smooth}의 번역을 사용한다.
먼저 이를 $a$에 적용하여 $a$가 대상에 대해 형식적으로 매끄러움을 안다.
다음으로 $f$가 가정에 의해 대상에 대해 형식적으로 매끄럽다는 사실
(보조정리 \ref{stacks-more-morphisms-lemma-reformulate-formal-smoothness}를 보라)과
『표현가능성의 판정 조건』의 보조정리
\ref{criteria-lemma-composition-formally-smooth}를 이용하여
$b = f \circ a$가 대상에 대해 형식적으로 매끄러움을 안다.
그런 다음 원리를 한 번 더 적용하여 $b$가 매끄럽다고 결론짓는다.
이는 대수적 스택의 사상에 대한 매끄러움의 정의에 의해 $f$가
매끄럽다는 뜻이고, 증명이 끝난다.
\end{proof}







\section{블로업과 평탄성}
\label{stacks-more-morphisms-section-blowup-flat}

\noindent
이 절에서는 대수공간 위의 대수적 스택에 관하여
『대수공간의 사상에 관한 추가 내용』의 절
\ref{spaces-more-morphisms-section-blowup-flat}과
\ref{spaces-more-morphisms-section-applications-flattening-by-blowing-up}에서
무엇을 추론할 수 있는지 간단히 논한다.

\begin{lemma}
\label{stacks-more-morphisms-lemma-flatten-stack}
$f : \mathcal{X} \to Y$를 대수적 스택에서 대수공간으로 가는 사상이라 하자.
$V \subset Y$를 열린 부분공간이라 하자. 다음을 가정하자.
\begin{enumerate}
\item $Y$는 준콤팩트이고 준분리이다.
\item $f$는 유한형이고 준분리이다.
\item $V$는 준콤팩트이다.
\item $\mathcal{X}_V$는 $V$ 위에서 평탄하고 국소 유한 표시이다.
\end{enumerate}
그러면 $V$-허용 블로업 $Y' \to Y$와
$\mathcal{X}'_V = \mathcal{X}_V$인 닫힌 부분스택
$\mathcal{X}' \subset \mathcal{X}_{Y'}$가 존재하여,
$\mathcal{X}' \to Y'$가 평탄하고 유한 표시이다.
\end{lemma}

\begin{proof}
$\mathcal{X}$가 준콤팩트임에 유의하자. 아핀 스킴 $U$와
전사 매끄러운 사상 $U \to \mathcal{X}$를 택한다.
$R = U \times_\mathcal{X} U$로 두면 $Y$ 위의 대수공간의 준군
$(U, R, s, t, c)$를 얻고 $\mathcal{X} = [U/R]$이다
(『대수적 스택』의 보조정리
\ref{algebraic-lemma-stack-presentation}).
$U \to Y$와 열린집합 $V \subset Y$에 『대수공간의 사상에 관한 추가 내용』의
보조정리 \ref{spaces-more-morphisms-lemma-flat-after-blowing-up}를
적용할 수 있다. 따라서 $V$-허용 블로업 $Y' \to Y$를 얻어,
고유변환 $U' \subset U_{Y'}$가 $Y'$ 위에서 평탄하고 유한 표시이다.
$R' \subset R_{Y'}$을 $R$의 고유변환이라 하자.
$s$와 $t$가 매끄럽고(특히 평탄하고), 『대수공간 위의 인자』의 보조정리
\ref{spaces-divisors-lemma-strict-transform-flat}에 의해
카르테시안 도식
$$
\vcenter{
\xymatrix{
R' \ar[r] \ar[d] & R_{Y'} \ar[d]^{s_{Y'}} \\
U' \ar[r] & U_{Y'}
}
}
\quad\text{그리고}\quad
\vcenter{
\xymatrix{
R' \ar[r] \ar[d] & R_{Y'} \ar[d]^{t_{Y'}} \\
U' \ar[r] & U_{Y'}
}
}
$$
을 얻는다. 달리 말하면 $U'$는 $U_{Y'}$의
$R_{Y'}$-불변 닫힌 부분공간이다. 따라서 『스택의 성질』의 보조정리
\ref{stacks-properties-lemma-substacks-presentation}에 의해 $U'$는
닫힌 부분스택 $\mathcal{X}' \subset \mathcal{X}_{Y'}$를 정한다.
$U' \to Y'$가 평탄하고 국소 유한 표시이므로
$\mathcal{X}' \to Y'$도 그러하다. 한편 $f$에 관한 가정과
닫힌 몰입에 대해서도 이 성질들이 성립한다는 사실에 의해,
$\mathcal{X}' \to Y'$가 준콤팩트이고 준분리라는 것도 이미 안다.
이로써 증명이 끝난다.
\end{proof}















\section{대수적 스택에 대한 저우 보조정리}
\label{stacks-more-morphisms-section-chow}

\noindent
이 절에서는 대수적 스택에 대한 저우 보조정리를 논한다.

\begin{lemma}
\label{stacks-more-morphisms-lemma-finite-cover-factor}
$Y$를 준콤팩트 준분리 대수공간이라 하자.
$V \subset Y$를 준콤팩트 열린집합이라 하자.
$f : \mathcal{X} \to V$가 전사이고 평탄하며 국소 유한 표시라고 하자.
그러면 유한 전사 사상 $g : Y' \to Y$가 존재하여
$V' = g^{-1}(V) \to Y$는 자리스키 국소적으로 $f$를 통해 분해된다.
\end{lemma}

\begin{proof}
먼저 $Y$가 스킴인 경우를 증명한다. 스킴 $U$와 전사 매끄러운 사상
$U \to \mathcal{X}$를 택할 수 있다. 그러면 $\{U \to V\}$는
스킴의 fppf 덮개이다. 『사상에 관한 추가 내용』의 보조정리
\ref{more-morphisms-lemma-dominate-fppf-finite}에 의해 유한 전사 사상
$V' \to V$가 존재하여 $V' \to V$는 자리스키 국소적으로 $U$를 통해
분해된다. 『사상에 관한 추가 내용』의 보조정리
\ref{more-morphisms-lemma-extend-finite-surjective-morphisms}에 의해,
원하는 대로 $V$로 제한하면 $V' \to V$가 되는 유한 전사 사상
$Y' \to Y$를 찾을 수 있다.

\medskip\noindent
$Y$가 대수공간이면, $Y'$가 스킴인 유한 전사 사상 $Y' \to Y$로
먼저 유한 밑변환을 함으로써 보조정리가 참임을 알 수 있다.
『대수공간의 극한』의 명제
\ref{spaces-limits-proposition-there-is-a-scheme-finite-over}를 보라.
\end{proof}

\begin{lemma}
\label{stacks-more-morphisms-lemma-make-section}
$f : \mathcal{X} \to Y$를 대수적 스택에서 대수공간으로 가는 사상이라 하자.
$V \subset Y$를 열린 부분공간이라 하자. 다음을 가정하자.
\begin{enumerate}
\item $f$는 분리이고 유한형이다.
\item $Y$는 준콤팩트이고 준분리이다.
\item $V$는 준콤팩트이다.
\item $\mathcal{X}_V$는 $V$ 위의 gerbe이다.
\end{enumerate}
그러면 가환 도식
$$
\xymatrix{
\overline{Z} \ar[rd]_{\overline{g}} &
Z \ar[l]^j \ar[d]_g \ar[r]_h & \mathcal{X} \ar[ld]^f \\
& Y
}
$$
이 존재한다. 여기서 $j$는 열린 몰입이고 $\overline{g}$와 $h$는
고유이며, $|V|$는 $|g|$의 상에 포함된다.
\end{lemma}

\begin{proof}
가환 도식
$$
\xymatrix{
\mathcal{X}' \ar[d]_{f'} \ar[r] & \mathcal{X} \ar[d]^f \\
Y' \ar[r] & Y
}
$$
과 준콤팩트 열린집합 $V' \subset Y'$가 주어졌다고 하자.
$Y' \to Y$는 대수공간의 고유 사상이고,
$\mathcal{X}' \to \mathcal{X}$는 대수적 스택의 고유 사상이며,
$V' \subset Y'$는 $V$ 위로 전사하고,
$\mathcal{X}'_{V'}$가 $V'$ 위의 gerbe라고 하자.
그러면 쌍 $(f' : \mathcal{X}' \to Y', V')$에 대해 보조정리를 증명하면
충분하다. 몇 가지 세부사항은 생략한다.

\medskip\noindent
증명의 전체 전략. 위의 관찰을 거듭 적용하여 $f$의 상이 열려 있고
$f$가 이 열린집합 위에서 단면을 갖는 경우로 환원한다.
각 단계는 간단하지만 단계가 꽤 많아서 증명이 조금 복잡하다.

\medskip\noindent
『대수공간의 극한』의 명제
\ref{spaces-limits-proposition-there-is-a-scheme-finite-over}를 이용하여
$Y$가 스킴인 경우로 환원한다. ($Y'$가 스킴인 유한 전사 사상
$Y' \to Y$를 잡고 $\mathcal{X}' = \mathcal{X}_{Y'}$로 둔 다음,
증명 첫머리의 관찰을 적용한다.)

\medskip\noindent
보조정리 \ref{stacks-more-morphisms-lemma-flatten-stack}을 이용하여
(gerbe가 평탄하고 국소 유한 표시임을 확인하려면
『스택의 사상』의 보조정리
\ref{stacks-morphisms-lemma-gerbe-fppf}도 사용한다),
$f$가 평탄하고 유한 표시인 경우로 환원한다.

\medskip\noindent
$f$가 평탄하고 국소 유한 표시이므로 $|f|$의 상은
열린집합 $W \subset Y$이다. $f$가 유한형이고 $Y$가 준콤팩트이므로
$\mathcal{X}$가 준콤팩트이고, 따라서 $W$도 준콤팩트이다.
보조정리 \ref{stacks-more-morphisms-lemma-finite-cover-factor}에 의해 유한 전사 사상
$g : Y' \to Y$를 찾아 $g^{-1}(W) \to Y$가 자리스키 국소적으로
$\mathcal{X} \to Y$를 통해 분해되게 할 수 있다.
$Y$를 $Y'$로, $\mathcal{X}$를 $\mathcal{X} \times_Y Y'$로 바꾸면
다음 문단에서 서술하는 상황으로 환원된다.

\medskip\noindent
$n \geq 0$, 준콤팩트 열린집합 $W_i \subset Y$ ($i = 1, \ldots, n$),
그리고 사상 $x_i : W_i \to \mathcal{X}$가 존재하여 다음을 만족한다고 하자.
(a) $f \circ x_i = \text{id}_{W_i}$,
(b) $W = \bigcup_{i = 1, \ldots, n} W_i$는 $V$를 포함하고,
(c) $W$는 $|f|$의 상이다.
$n$에 대한 귀납법을 사용한다. 시작 경우는 $n = 0$이다.
이는 $V = \emptyset$을 뜻하며, 이 경우 $\overline{Z} = \emptyset$으로
둘 수 있다. $n > 0$이면 각 $i = 1, \ldots, n$에 대해
바탕 위상공간이 $Y \setminus W_i$인 축약 닫힌 부분스킴 $Y_i$를 생각하자.
유한 사상
$$
Y' = Y \amalg \coprod\nolimits_{i = 1, \ldots, n} Y_i \longrightarrow Y
$$
와 준콤팩트 열린집합
$$
V' = (W_1 \cap \ldots \cap W_n \cap V) \amalg
\coprod_{i = 1, \ldots, n} (V \cap Y_i).
$$
를 생각하자. 증명 첫머리의 관찰에 의해, 쌍
$$
(\mathcal{X} \to Y, W_1 \cap \ldots \cap W_n \cap V)
\quad\text{그리고}\quad
(\mathcal{X} \times_Y Y_i \to Y_i, V \cap Y_i),\quad
i = 1, \ldots, n
$$
에 대해 보조정리를 증명할 수 있으면 결과가 참이다.
여기서 집합론적 등식
$V = (W_1 \cap \ldots \cap W_n \cap V) \cup
\bigcup\nolimits_{i = 1, \ldots n} (V \cap Y_i)$를 사용한다.
위의 두 번째 유형의 쌍에는 귀납 가정이 적용된다.
따라서 다음 문단의 상황으로 환원된다.

\medskip\noindent
$n \geq 0$, 준콤팩트 열린집합 $W_i \subset Y$ ($i = 1, \ldots, n$),
그리고 사상 $x_i : W_i \to \mathcal{X}$가 존재하여 다음을 만족한다고 하자.
(a) $f \circ x_i = \text{id}_{W_i}$,
(b) $W = \bigcup_{i = 1, \ldots, n} W_i$는 $V$를 포함하고,
(c) $W$는 $|f|$의 상이며,
(d) $V \subset W_1 \cap \ldots \cap W_n$이다.
사상
$$
T_{ij} = \mathit{Isom}_\mathcal{X}(x_i|_{W_i \cap W_j \cap V},
x_j|_{W_i \cap W_j \cap V}) \longrightarrow W_i \cap W_j \cap V
$$
은 전사이고 평탄하며 국소 유한 표시이다
(『스택의 사상』의 보조정리
\ref{stacks-morphisms-lemma-gerbe-isom-fppf}).
각 준콤팩트 열린집합 $W_i \cap W_j \cap V$와 사상
$T_{ij} \to W_i \cap W_j \cap V$에 보조정리
\ref{stacks-more-morphisms-lemma-finite-cover-factor}를 적용하여 유한 전사 사상
$Y'_{ij} \to Y$를 얻는다. $Y$를 모든 $Y'_{ij}$의 $Y$ 위의 섬유곱으로
바꾸면 다음 문단의 상황으로 환원된다.

\medskip\noindent
$n \geq 0$, 준콤팩트 열린집합 $W_i \subset Y$ ($i = 1, \ldots, n$),
그리고 사상 $x_i : W_i \to \mathcal{X}$가 존재하여 다음을 만족한다고 하자.
(a) $f \circ x_i = \text{id}_{W_i}$,
(b) $W = \bigcup_{i = 1, \ldots, n} W_i$는 $V$를 포함하고,
(c) $W$는 $|f|$의 상이며,
(d) $V \subset W_1 \cap \ldots \cap W_n$이고,
(e) $x_i$와 $x_j$는 $W_i \cap W_j \cap V$ 위에서 자리스키 국소적으로
동형이다. $y \in V$를 임의로 잡자. 보조정리가 쌍
$(\mathcal{X} \to Y, V_y)$에 대해 참이게 하는 준콤팩트 열린 근방
$y \in V_y \subset V$를 찾을 수 있다고 하자. 그 해를
$\overline{Z}_y, Z_y, \overline{g}_y, g_y, h_y$라 하자.
$V$가 준콤팩트이므로
$V = V_{y_1} \cup \ldots \cup V_{y_m}$인 유한 개의 점
$y_1, \ldots, y_m$을 찾을 수 있다. 그러면
$$
\overline{Z} = \coprod \overline{Z}_{y_j},\quad
Z = \coprod Z_{y_j},\quad
\overline{g} = \coprod \overline{g}_{y_j},\quad
g = \coprod g_{y_j},\quad
h = \coprod h_{y_j}
$$
로 두면 보조정리의 해가 됨을 알 수 있다. $y$가 주어지면 조건 (e)에 의해
준콤팩트 열린 근방 $y \in V_y \subset V$와 $i = 2, \ldots, n$에 대한
동형사상 $\varphi_i : x_1|_{V_y} \to x_i|_{V_y}$를 택할 수 있다.
$\varphi_{ij} = \varphi_j \circ \varphi_i^{-1}$로 두자.
이로써 다음 문단의 상황에 이른다.

\medskip\noindent
$n \geq 0$, 준콤팩트 열린집합 $W_i \subset Y$ ($i = 1, \ldots, n$),
그리고 사상 $x_i : W_i \to \mathcal{X}$가 존재하여 다음을 만족한다고 하자.
(a) $f \circ x_i = \text{id}_{W_i}$,
(b) $W = \bigcup_{i = 1, \ldots, n} W_i$는 $V$를 포함하고,
(c) $W$는 $|f|$의 상이며,
(d) $V \subset W_1 \cap \ldots \cap W_n$이고,
(f) $\varphi_{jk} \circ \varphi_{ij} = \varphi_{ik}$을 만족하는 동형사상
$\varphi_{ij} : x_i|_V \to x_j|_V$가 존재한다. 사상
$$
I_{ij} = \mathit{Isom}_\mathcal{X}(x_i|_{W_i \cap W_j},
x_j|_{W_i \cap W_j}) \longrightarrow W_i \cap W_j
$$
은 $f$가 분리이므로 고유이다
(『스택의 사상』의 보조정리
\ref{stacks-morphisms-lemma-separated-implies-isom}).
$\varphi_{ij}$는 $V$ 위에서 $I_{ij} \to W_i \cap W_j$의 단면
$V \to I_{ij}$를 정함에 유의하자. 『대수공간의 사상에 관한 추가 내용』의
보조정리 \ref{spaces-more-morphisms-lemma-get-section-after-blowup}에 의해
$s_{ij}$가 $p_{ij}^{-1}(W_i \cap W_j)$로 확장되게 하는
$V$-허용 블로업 $p_{ij} : Y_{ij} \to Y$를 찾을 수 있다.
$Y$를 모든 $Y_{ij}$의 $Y$ 위의 섬유곱으로 바꾸면 다음 문단의
상황에 이른다.

\medskip\noindent
$n \geq 0$, 준콤팩트 열린집합 $W_i \subset Y$ ($i = 1, \ldots, n$),
그리고 사상 $x_i : W_i \to \mathcal{X}$가 존재하여 다음을 만족한다고 하자.
(a) $f \circ x_i = \text{id}_{W_i}$,
(b) $W = \bigcup_{i = 1, \ldots, n} W_i$는 $V$를 포함하고,
(c) $W$는 $|f|$의 상이며,
(d) $V \subset W_1 \cap \ldots \cap W_n$이고,
(g) 동형사상
$\varphi_{ij} : x_i|_{W_i \cap W_j} \to x_j|_{W_i \cap W_j}$가 존재하여
$$
\varphi_{jk}|_V \circ \varphi_{ij}|_V = \varphi_{ik}|_V.
$$
를 만족한다. 필요하다면 $Y$를 다른 $V$-허용 블로업으로 바꾼 뒤,
$V$가 $Y$에서 조밀하고 스킴론적으로 조밀하다고 가정할 수 있고,
따라서 $V$를 포함하는 $Y$의 임의의 열린 부분공간에서도 그러하다.
이렇게 바꾼 뒤에는 『대수공간의 사상』의 보조정리
\ref{spaces-morphisms-lemma-equality-of-morphisms}와
$I_{ik} \to W_i \cap W_j$가 고유(따라서 분리)라는 사실에 호소하여
$$
\varphi_{jk}|_{W_i \cap W_j \cap W_k} \circ
\varphi_{ij}|_{W_i \cap W_j \cap W_k} =
\varphi_{ik}|_{W_i \cap W_j \cap W_k}
$$
를 얻는다. 물론 이는 $(x_i, \varphi_{ij})$가 내림 자료라는 뜻이고,
$\mathcal{X}$가 스택이므로 $W_i$ 위에서 $x_i$와 일치하는 사상
$x : W \to \mathcal{X}$를 얻는다. $x$는 분리 사상
$\mathcal{X} \to W$의 단면이므로 $x$는 고유이다
(『스택의 사상』의 보조정리
\ref{stacks-morphisms-lemma-section-immersion}).
따라서 이제 $\overline{Z} = Y$, $Z = W$,
$\overline{g} = \text{id}_Y$, $g = \text{id}_W$, $h = x$로 두면
보조정리가 성립한다.
\end{proof}

\begin{theorem}[저우 보조정리]
\label{stacks-more-morphisms-theorem-chow-finite-type}
\begin{reference}
이는 Ofer Gabber의 결과이다. \cite[Theorem 1.1]{olsson_proper}를 보라.
\end{reference}
$f : \mathcal{X} \to Y$를 대수적 스택에서 대수공간으로 가는 사상이라 하자.
다음을 가정하자.
\begin{enumerate}
\item $Y$는 준콤팩트이고 준분리이다.
\item $f$는 분리 유한형이다.
\end{enumerate}
그러면 가환 도식
$$
\xymatrix{
\mathcal{X} \ar[rd] & X \ar[l] \ar[d] \ar[r] & \overline{X} \ar[ld] \\
& Y
}
$$
이 존재한다. 여기서 $X \to \mathcal{X}$는 고유 전사이고,
$X \to \overline{X}$는 열린 몰입이며,
$\overline{X} \to Y$는 대수공간의 고유 사상이다.
\end{theorem}

\begin{proof}
대략적인 생각은 $\mathcal{X}$가 gerbe인 조밀한 열린집합을 갖는다는 사실
(『스택의 사상』의 명제
\ref{stacks-morphisms-proposition-open-stratum})을 이용하고
보조정리 \ref{stacks-more-morphisms-lemma-make-section}에 호소하는 것이다.
이 방법이 그대로는 작동하지 않는 이유는 그 열린집합이 준콤팩트가
아닐 수 있어 기술적인 문제에 부딪히기 때문이다.
따라서 먼저 뇌터 경우로 (표준적으로) 환원한다.

\medskip\noindent
먼저 $\mathcal{X}'$가 $Y$ 위에서 분리이고 유한형인 대수적 스택인
닫힌 몰입 $\mathcal{X} \to \mathcal{X}'$을 택한다.
『스택의 극한』의 보조정리
\ref{stacks-limits-lemma-separated-closed-in-finite-presentation}를 보라.
$\mathcal{X}'$에 대해 정리를 증명하면 충분함이 명백하므로,
$\mathcal{X} \to Y$가 분리이고 유한 표시라고 가정할 수 있다.

\medskip\noindent
$\mathcal{X} \to Y$가 분리이고 유한 표시라고 가정하자.
『대수공간의 극한』의 명제 \ref{spaces-limits-proposition-approximate}에 의해
$Y = \lim Y_i$를 아핀 전이 사상을 갖는 뇌터 대수공간 계의 유향 극한으로
쓸 수 있다. 『스택의 극한』의 보조정리
\ref{stacks-limits-lemma-descend-a-stack-down}에 의해, 어떤 $i$와
대수적 스택에서 $Y_i$로 가는 유한 표시 사상
$\mathcal{X}_i \to Y_i$가 존재하여
$\mathcal{X} = Y \times_{Y_i} \mathcal{X}_i$이다.
$i$를 키운 뒤 $\mathcal{X}_i \to Y_i$가 분리라고 가정할 수 있다.
『스택의 극한』의 보조정리
\ref{stacks-limits-lemma-eventually-separated}를 보라.
그러면 $\mathcal{X}_i \to Y_i$에 대해 정리를 증명하면 충분하다.
이로써 다음 문단의 경우로 환원된다.

\medskip\noindent
$Y$가 뇌터라고 가정하자. $\mathcal{X}$를 그 축약으로 바꾸어도 된다
(『스택의 성질』의 정의
\ref{stacks-properties-definition-reduced-induced-stack}).
이로써 다음 문단의 경우로 환원된다.

\medskip\noindent
$Y$가 뇌터이고 $\mathcal{X}$가 축약이라고 가정하자.
$\mathcal{X} \to Y$가 분리이고 $Y$가 준분리이므로
$\mathcal{X}$는 대수적 스택으로서 준분리이다. 따라서 관성 사상
$\mathcal{I}_\mathcal{X} \to \mathcal{X}$는 준콤팩트이다.
그러므로 『스택의 사상』의 명제
\ref{stacks-morphisms-proposition-open-stratum}에 의해 gerbe인 조밀한
열린 부분스택 $\mathcal{V} \subset \mathcal{X}$가 존재한다.
$\mathcal{V}$를 대수공간 $V$ 위의 gerbe로 나타내는 사상을
$\mathcal{V} \to V$라 하자. $\mathcal{V} \to V$의 구성은
『스택의 사상』의 보조정리
\ref{stacks-morphisms-lemma-gerbe-over-iso-classes}를 보라.
특히 이 구성은 사상 $\mathcal{V} \to Y$가
$\mathcal{V} \to V \to Y$로 분해됨을 보여준다. 다음과 같다.
$$
\xymatrix{
\mathcal{V} \ar[r] \ar[d] & \mathcal{X} \ar[d] \\
V \ar[r] & Y
}
$$
사상 $\mathcal{V} \to V$가 전사이고 평탄하며 유한 표시이고
(『스택의 사상』의 보조정리 \ref{stacks-morphisms-lemma-gerbe-fppf}),
$\mathcal{V} \to Y$가 국소 유한 표시이므로, 『스택의 사상』의 보조정리
\ref{stacks-morphisms-lemma-flat-finite-presentation-permanence}에 의해
$V \to Y$는 국소 유한 표시이다. $\mathcal{V} \to V$가
보편 위상동형임에 유의하자(『스택의 사상』의 보조정리
\ref{stacks-morphisms-lemma-gerbe-bijection-points}).
$\mathcal{V}$가 준콤팩트이므로(『스택의 사상』의 보조정리
\ref{stacks-morphisms-lemma-locally-closed-in-noetherian}를 보라),
$V$도 준콤팩트이다. 마지막으로 $\mathcal{V} \to Y$가 분리이므로
$\mathcal{V} \to V \to Y$에 『스택의 사상』의 보조정리
\ref{stacks-morphisms-lemma-check-separated-on-ui-cover}를 적용하면
$V \to Y$도 분리이다(그 가정들은 이미 확인했다).

\medskip\noindent
위의 모든 사실은 『대수공간의 극한』의 보조정리
\ref{spaces-limits-lemma-embedding-into-affine-over-qs}의 가정을
사상 $V \to Y$에 적용할 수 있다는 뜻이다. 따라서 조밀한 열린 부분공간
$V' \subset V$와 $Y$ 위의 몰입 $V' \to \mathbf{P}^n_Y$를 찾을 수 있다.
$V$를 $V'$로, $\mathcal{V}$를 $\mathcal{V}$에서 $V'$의 역상으로
바꾸어도 됨이 명백하다(위에서 보았듯이
$|\mathcal{V}| = |V|$임을 상기하자). 따라서 다음 도식이 있다고
가정할 수 있다.
$$
\xymatrix{
\mathcal{V} \ar[rr] \ar[d] & & \mathcal{X} \ar[d] \\
V \ar[r] & \mathbf{P}^n_Y \ar[r] & Y
}
$$
여기서 화살표 $V \to \mathbf{P}^n_Y$는 몰입이다. 사상
$$
j : \mathcal{V} \longrightarrow \mathbf{P}^n_Y \times_Y \mathcal{X}
$$
의 스킴론적 상을 $\mathcal{X}'$라 하고, 사상
$V \to \mathbf{P}^n_Y$의 스킴론적 상을 $Y'$라 하자.
가환 도식
$$
\xymatrix{
\mathcal{V} \ar[r] \ar[d] &
\mathcal{X}' \ar[r] \ar[d] &
\mathbf{P}^n_Y \times_Y \mathcal{X} \ar[d] \ar[r] &
\mathcal{X} \ar[d] \\
V \ar[r] &
Y' \ar[r] &
\mathbf{P}^n_Y \ar[r] &
Y
}
$$
을 얻는다(『스택의 사상』의 보조정리
\ref{stacks-morphisms-lemma-factor-factor}를 보라).
$\mathcal{V} = V \times_{Y'} \mathcal{X}'$이고, 보조정리
\ref{stacks-more-morphisms-lemma-make-section}을 사상 $\mathcal{X}' \to Y'$와
열린 부분공간 $V \subset Y'$에 적용할 수 있다고 주장한다.
이 주장이 참이면
$$
\xymatrix{
\overline{X} \ar[rd]_{\overline{g}} &
X \ar[l] \ar[d]_g \ar[r]_h & \mathcal{X}' \ar[ld]^f \\
& Y'
}
$$
을 얻으며, 여기서 $X \to \overline{X}$는 열린 몰입이고
$\overline{g}$와 $h$는 고유이며 $|V|$는 $|g|$의 상에 포함된다.
그러면 합성 $X \to \mathcal{X}' \to \mathcal{X}$는 고유 사상들의
합성으로서 고유이고, 그 상이 $|\mathcal{V}|$를 포함하므로 이 합성은
전사이다. 또한 $\overline{X} \to Y' \to Y$는 고유 사상들의 합성으로서
고유이다.

\medskip\noindent
마지막 단계는 이 주장을 증명하는 것이다.
$\mathcal{X}' \to Y'$가 분리이고 유한형이며,
$Y'$가 준콤팩트이고 준분리이고, $V$가 준콤팩트임에 유의하자
(모든 세부사항의 완전한 확인은 생략한다).
다음으로, 『스택의 사상』의 보조정리
\ref{stacks-morphisms-lemma-scheme-theoretic-image-of-partial-section}에 의해
$b : \mathcal{X}' \to \mathcal{X}$가 $\mathcal{V}$ 위에서
동형사상임에 유의하자. 특히 $\mathcal{V}$는 $\mathcal{X}'$의
열린 부분스택과 동일시된다. 사상 $j$는 준콤팩트이다
(원천은 준콤팩트이고 표적은 준분리이다). 따라서 『스택의 사상』의
보조정리 \ref{stacks-morphisms-lemma-existence-plus-flat-base-change}에 의해
$j$의 스킴론적 상을 만드는 것은 평탄 밑변환과 가환한다.
특히 $V \times_{Y'} \mathcal{X}'$는 사상
$\mathcal{V} \to V \times_{Y'} \mathcal{X}'$의 스킴론적 상이다.
그러나 『스택의 사상』의 보조정리
\ref{stacks-morphisms-lemma-universally-closed-permanence}에 의해
$|\mathcal{V}| \to |V \times_{Y'} \mathcal{X}'|$의 상은 닫혀 있다
($\mathcal{V} \to V$가 위에서 본 대로 보편 위상동형이고 따라서
보편 닫힘임을 사용한다). 또한 그 상은 조밀하다
(방금 말한 사실을 『스택의 사상』의 보조정리
\ref{stacks-morphisms-lemma-topology-scheme-theoretic-image}와 합친다).
따라서 $|\mathcal{V}| = |V \times_{Y'} \mathcal{X}'|$이라고 결론짓는다.
그러므로 $\mathcal{V} \to V \times_{Y'} \mathcal{X}'$는 동형사상이고,
주장의 증명이 끝난다.
\end{proof}



\section{뇌터 값매김 판정법}
\label{stacks-more-morphisms-section-Noetherian-valuative-criterion}

\noindent
이 절에서는 뇌터 상황에서 이산 값매김환만을 사용하여 대수적 스택의
사상에 대한 (정련된) 값매김 판정법을 논한다. 서로 다른 변형이 많으며,
시간이 지나면서 필요에 따라 여기에 더 추가할 것이다.

\medskip\noindent
$f : \mathcal{X} \to \mathcal{Y}$를 대수적 스택(또는 대수공간이나 스킴)의
사상이라 하자. {\it 정련된} 값매김 판정법이란 어떤 성질을 갖는 사상
$\mathcal{U} \to \mathcal{X}$가 주어지고, 다음 꼴의 실선 도식에서
점선 화살표의 존재성이나 유일성만을 살펴보는 판정법이다.
$$
\xymatrix{
\Spec(K) \ar[d] \ar[r] & \mathcal{U} \ar[r] & \mathcal{X} \ar[d] \\
\Spec(A) \ar[rr] \ar@{..>}[rru] & & \mathcal{Y}
}
$$
아래에서는 지금까지 얻은 결과를 서술할 때 이 용어를 사용한다.

\medskip\noindent
대수적 스택의 사상에 대한 비뇌터 값매김 판정법은 다음과 같다.
\begin{enumerate}
\item 『스택의 사상』의 절
\ref{stacks-morphisms-section-valuative-second}
(대각사상의 분리성에 관하여),
\item 『스택의 사상』의 절
\ref{stacks-morphisms-section-valuative-diagonal}
(분리성에 관하여),
\item 『스택의 사상』의 절
\ref{stacks-morphisms-section-valutive-criterion}
(보편 닫힘성에 관하여),
\item 『스택의 사상』의 절
\ref{stacks-morphisms-section-valutive-criterion-properness}
(고유성에 관하여).
\end{enumerate}
대수공간에 대해서는 다음 값매김 판정법들이 있다.
\begin{enumerate}
\item 『대수공간의 사상』의 절
\ref{spaces-morphisms-section-valuative-criterion-universally-closed}
(보편 닫힘성에 관하여),
\item 『대수공간의 사상』의 보조정리
\ref{spaces-morphisms-lemma-refined-valuative-criterion-universally-closed}
(보편 닫힘성에 대한 정련된 판정법),
\item 『대수공간의 사상』의 절
\ref{spaces-morphisms-section-valuative-separatedness}
(분리성에 관하여),
\item 『대수공간의 사상』의 절
\ref{spaces-morphisms-section-valuative-criterion-properness}
(고유성에 관하여),
\item 『양호한 대수공간』의 절
\ref{decent-spaces-section-valuative-criterion-universally-closed}
(양호한 대수공간의 보편 닫힘성에 관하여),
\item 『양호한 대수공간』의 보조정리
\ref{decent-spaces-lemma-re-characterize-universally-closed}
(대수공간 사이의 양호한 사상의 보편 닫힘성에 관하여),
\item 『대수공간의 코호몰로지』의 절
\ref{spaces-cohomology-section-Noetherian-valuative-criterion}에는
다음 뇌터 값매김 판정법들이 있다.
\begin{enumerate}
\item 『대수공간의 코호몰로지』의 보조정리
\ref{spaces-cohomology-lemma-check-separated-dvr}
(이산 값매김환을 사용하는 분리성 판정법),
\item 『대수공간의 코호몰로지』의 보조정리
\ref{spaces-cohomology-lemma-check-proper-dvr}
(이산 값매김환을 사용하는 고유성 판정법),
\item 『대수공간의 코호몰로지』의 주석
\ref{spaces-cohomology-remark-variant}
(완비 이산 값매김환의 경우로 환원하는 법을 논한다).
\end{enumerate}
\item 『대수공간의 극한』의 절
\ref{spaces-limits-section-Noetherian-valuative-criterion}에서는
다음 뇌터 값매김 판정법들을 논한다.
\begin{enumerate}
\item 『대수공간의 극한』의 보조정리
\ref{spaces-limits-lemma-Noetherian-dvr-valuative-separation}
(이산 값매김환과 일반점을 사용하는 분리성 판정법),
\item 『대수공간의 극한』의 보조정리
\ref{spaces-limits-lemma-Noetherian-dvr-valuative-proper}
(이산 값매김환과 일반점을 사용하는 고유성 판정법),
\item 『대수공간의 극한』의 보조정리
\ref{spaces-limits-lemma-check-universally-closed-Noetherian}
(이산 값매김환을 사용하는 보편 닫힘성 판정법).
\end{enumerate}
\item 『대수공간의 극한』의 절
\ref{spaces-limits-section-refined-valuative-criteria}에서는
다음 정련된 뇌터 값매김 판정법들을 논한다.
\begin{enumerate}
\item 『대수공간의 극한』의 보조정리
\ref{spaces-limits-lemma-refined-valuative-criterion-proper}와
\ref{spaces-limits-lemma-refined-valuative-criterion-universally-closed}
(이산 값매김환을 사용하는 고유성에 대한 정련된 판정법),
\item 『대수공간의 극한』의 보조정리
\ref{spaces-limits-lemma-refined-valuative-criterion-separated}
(이산 값매김환을 사용하는 분리성에 대한 정련된 판정법).
\end{enumerate}
\end{enumerate}
스킴에 대해서는 다음 값매김 판정법들이 있다.
\begin{enumerate}
\item 『스킴』의 절
\ref{schemes-section-valuative-criterion-universal-closedness}
(보편 닫힘성에 관하여),
\item 『스킴』의 절 \ref{schemes-section-valuative-separatedness}
(분리성에 관하여),
\item 『사상』의 절 \ref{morphisms-section-valuative-criteria}
(고유성에 관하여),
\item 『사상』의 보조정리
\ref{morphisms-lemma-refined-valuative-criterion-universally-closed}
(보편 닫힘성에 대한 정련된 판정법),
\item 『극한』의 절 \ref{limits-section-Noetherian-valuative-criterion}에서는
다음 뇌터 값매김 판정법들을 논한다.
\begin{enumerate}
\item 『극한』의 보조정리
\ref{limits-lemma-Noetherian-dvr-valuative-separation}
(이산 값매김환과 일반점을 사용하는 분리성 판정법),
\item 『극한』의 보조정리
\ref{limits-lemma-Noetherian-dvr-valuative-proper}
(이산 값매김환과 일반점을 사용하는 고유성 판정법),
\item 『극한』의 보조정리
\ref{limits-lemma-check-universally-closed-Noetherian}
(이산 값매김환을 사용하는 보편 닫힘성 판정법).
\end{enumerate}
\item 『극한』의 절 \ref{limits-section-refined-valuative-criteria}에서는
다음 정련된 뇌터 값매김 판정법들을 논한다.
\begin{enumerate}
\item 『극한』의 보조정리
\ref{limits-lemma-refined-valuative-criterion-proper}와
\ref{limits-lemma-refined-valuative-criterion-universally-closed}
(이산 값매김환을 사용하는 고유성에 대한 정련된 판정법),
\item 『극한』의 보조정리
\ref{limits-lemma-refined-valuative-criterion-separated}
(이산 값매김환을 사용하는 분리성에 대한 정련된 판정법).
\end{enumerate}
\item 『극한』의 절 \ref{limits-section-nagata-valuative}에서는
밑 위에서 본질적으로 유한형인 이산 값매김환을 얻을 수 있는 경우,
뇌터 밑 위의 값매김 판정법을 논한다.
\end{enumerate}
이로써 앞서 나온 결과들의 목록을 마친다.

\medskip\noindent
이 절의 결과 가운데 다수는 다음 보조정리에 호소하여 증명할 수 있고
(어쩌면 그렇게 증명해야 하지만), 항상 그렇게 하지는 않았다.

\begin{lemma}
\label{stacks-more-morphisms-lemma-reach-point-closure-Noetherian}
$f : \mathcal{X} \to \mathcal{Y}$를 대수적 스택의 사상이라 하자.
$f$가 유한형이고 $\mathcal{Y}$가 국소 뇌터라고 가정하자.
$y \in |\mathcal{Y}|$를 $|f|$의 상의 폐포에 속하는 점이라 하자.
그러면 대수적 스택의 가환 도식
$$
\xymatrix{
\Spec(K) \ar[r] \ar[d] & \mathcal{X} \ar[d]^f \\
\Spec(A) \ar[r] & \mathcal{Y}
}
$$
이 존재한다. 여기서 $A$는 이산 값매김환이고 $K$는 그 분수체이며,
$\Spec(A)$의 닫힌점은 $y$로 간다.
\end{lemma}

\begin{proof}
아핀 스킴 $V$, 점 $v \in V$, 그리고 $v$를 $y$로 보내는 매끄러운 사상
$V \to \mathcal{Y}$를 택한다. 사상 $|V| \to |\mathcal{Y}|$는 열리고,
『스택의 성질』의 보조정리
\ref{stacks-properties-lemma-points-cartesian}에 의해
$|\mathcal{X} \times_\mathcal{Y} V| \to |V|$의 상은 $|f|$의 상의 역상이다.
따라서 점 $v$는 $|\mathcal{X} \times_\mathcal{Y} V| \to |V|$의
상의 폐포에 속한다. $\mathcal{X} \times_\mathcal{Y} V \to V$와 점 $v$에
대해 보조정리를 증명하면 $f$와 $y$에 대한 보조정리가 따른다.
이로써 다음 문단의 상황으로 환원된다.

\medskip\noindent
$Y$가 뇌터 아핀 스킴인 경우 보조정리에서와 같은
$f : \mathcal{X} \to Y$와 $y \in |Y|$가 주어졌다고 하자.
$f$가 준콤팩트이므로 $\mathcal{X}$도 준콤팩트이다.
따라서 아핀 스킴 $W$와 전사 매끄러운 사상
$W \to \mathcal{X}$를 택할 수 있다. 그러면 $|f|$의 상은
$|W| \to |Y|$의 상과 같다. 이로써 『극한』의 보조정리
\ref{limits-lemma-reach-point-closure-Noetherian}인 스킴의 경우로 환원된다.
\end{proof}

\begin{lemma}
\label{stacks-more-morphisms-lemma-check-separated-dvr}
$f : \mathcal{X} \to \mathcal{Y}$를 대수적 스택의 사상이라 하자.
다음을 가정하자.
\begin{enumerate}
\item $\mathcal{Y}$는 국소 뇌터이다.
\item $f$는 국소 유한형이고 준분리이다.
\item 모든 가환 도식
$$
\xymatrix{
\Spec(K) \ar[r]_x \ar[d]_j & \mathcal{X} \ar[d]^f \\
\Spec(A) \ar[r]^y \ar@{-->}[ru] & \mathcal{Y}
}
$$
에서 $A$가 이산 값매김환이고 $K$가 그 분수체일 때,
그리고 임의의 $2$-화살표 $\gamma : y \circ j \to f \circ x$에 대해,
점선 화살표의 범주(『스택의 사상』의 정의
\ref{stacks-morphisms-definition-fill-in-diagram})는 공집합이거나
동형류가 정확히 하나인 세토이드이다.
\end{enumerate}
그러면 $f$는 분리이다.
\end{lemma}

\begin{proof}
$f$가 분리임을 증명하려면
$\Delta : \mathcal{X} \to \mathcal{X} \times_\mathcal{Y} \mathcal{X}$가
고유임을 보여야 한다. $\Delta$가 대수공간으로 표현가능하고
국소 유한형이라는 것은 이미 안다(『스택의 사상』의 보조정리
\ref{stacks-morphisms-lemma-properties-diagonal}). 또한 $f$가
준분리라는 정의에 의해 이는 준콤팩트이고 준분리이다.
스킴 $U$와 전사 매끄러운 사상
$U \to \mathcal{X} \times_\mathcal{Y} \mathcal{X}$를 택한다.
$$
V = \mathcal{X} \times_{\Delta, \mathcal{X} \times_\mathcal{Y} \mathcal{X}} U
$$
로 두자. 대수공간의 사상 $V \to U$가 고유임을 보이면 충분하다
(『스택의 성질』의 보조정리
\ref{stacks-properties-lemma-check-property-covering}).
$U$는 국소 뇌터이고(『스택의 사상』의 보조정리
\ref{stacks-morphisms-lemma-locally-finite-type-locally-noetherian}와
$U \to \mathcal{Y}$가 국소 유한형이라는 사실을 사용한다),
$V \to U$는 $\Delta$의 밑변환이므로 유한형이고 준분리임에 유의하자
(위에 열거한 $\Delta$의 성질도 사용한다). 『대수공간의 코호몰로지』의
보조정리 \ref{spaces-cohomology-lemma-check-proper-dvr}를 적용하면
다음을 보이면 충분하다. $A$가 이산 값매김환이고 $K$가 그 분수체인
가환 도식
$$
\xymatrix{
\Spec(K) \ar[r]_v \ar[d]_j &
V \ar[d]^g \ar[r] &
\mathcal{X} \ar[d]^\Delta \\
\Spec(A) \ar[r]^u \ar@{-->}[ru] \ar@{..>}[rru] &
U \ar[r] &
\mathcal{X} \times_\mathcal{Y} \mathcal{X}
}
$$
이 주어졌을 때 이를 가환으로 만드는 유일한 파선 화살표가 존재해야 한다.
『스택의 사상』의 보조정리
\ref{stacks-morphisms-lemma-cat-dotted-arrows-base-change}에 의해
파선 화살표의 범주와 점선 화살표의 범주는 동치이다.
가정 (3)에 의해 동형사상을 제외하면 유일한 점선 화살표가 존재한다.
『스택의 사상』의 보조정리
\ref{stacks-morphisms-lemma-helper-diagonal}를 보라.
따라서 원하는 대로 유일한 파선 화살표가 존재한다고 결론짓는다.
\end{proof}

\begin{lemma}
\label{stacks-more-morphisms-lemma-refined-valuative-criterion-proper}
$f : \mathcal{X} \to \mathcal{Y}$와 $h : \mathcal{U} \to \mathcal{X}$를
대수적 스택의 사상이라 하자. $\mathcal{Y}$가 국소 뇌터이고,
$f$와 $h$가 유한형이며, $f$가 분리이고,
$|h| : |\mathcal{U}| \to |\mathcal{X}|$의 상이
$|\mathcal{X}|$에서 조밀하다고 가정하자. $A$가 분수체 $K$를 갖는
이산 값매김환이고
$\gamma : y \circ j \to f \circ h \circ u$인 임의의 $2$-가환 도식
$$
\xymatrix{
\Spec(K) \ar[r]_-u \ar[d]_j & \mathcal{U} \ar[r]_h & \mathcal{X} \ar[d]^f \\
\Spec(A) \ar[rr]^-y & & \mathcal{Y}
}
$$
이 주어졌을 때 다음이 성립한다고 하자. 체의 확대 $K'/K$와
$A$를 지배하는 값매김환 $A' \subset K'$가 존재하여, 유도된 도식
$$
\xymatrix{
\Spec(K') \ar[r]_-{x'} \ar[d]_{j'} & \mathcal{X} \ar[d]^f \\
\Spec(A') \ar[r]^-{y'} \ar@{..>}[ru] & \mathcal{Y}
}
$$
및 유도된 $2$-화살표
$\gamma' : y' \circ j' \to f \circ x'$에 대한 점선 화살표의 범주가
공집합이 아니다(『스택의 사상』의 정의
\ref{stacks-morphisms-definition-fill-in-diagram}). 그러면 $f$는 고유이다.
\end{lemma}

\begin{proof}
$f$가 보편 닫힘임을 증명하면 충분하다. $V$가 아핀 스킴인 매끄러운 사상
$V \to \mathcal{Y}$를 잡자. 『스택의 성질』의 보조정리
\ref{stacks-properties-lemma-points-cartesian}에 의해
$|\mathcal{U} \times_\mathcal{Y} V| \to |\mathcal{X} \times_\mathcal{Y} V|$의 상 $I$는 $|h|$의 상의 역상이다.
$|\mathcal{X} \times_\mathcal{Y} V| \to |\mathcal{X}|$가 열리므로
(『스택의 사상』의 보조정리
\ref{stacks-morphisms-lemma-fppf-open}),
$I$는 $|\mathcal{X} \times_\mathcal{Y} V|$에서 조밀하다.
또한 점선 화살표의 범주는 밑변환에 대해 잘 행동하므로
(『스택의 사상』의 보조정리
\ref{stacks-morphisms-lemma-cat-dotted-arrows-base-change}),
확대 뒤의 점선 화살표의 존재에 관한 가정은 사상들
$\mathcal{U} \times_\mathcal{Y} V \to \mathcal{X} \times_\mathcal{Y} V \to V$에 물려진다.
따라서 사상들
$\mathcal{U} \times_\mathcal{Y} V \to \mathcal{X} \times_\mathcal{Y} V \to V$는 보조정리의 가정을 만족한다.
그러므로 $\mathcal{Y}$가 아핀 스킴이라고 가정해도 된다.

\medskip\noindent
$\mathcal{Y} = Y$가 아핀 스킴이라고 가정하자.
(이제부터 $2$-화살표 $\gamma$와 $\gamma'$는 더 이상 신경 쓰지 않아도 된다.
『스택의 사상』의 보조정리
\ref{stacks-morphisms-lemma-cat-dotted-arrows-independent}를 보라.)
그러면 $\mathcal{U}$는 준콤팩트이다. 아핀 스킴 $U$와
전사 매끄러운 사상 $U \to \mathcal{U}$를 택한다.
이제 $\mathcal{U}$를 $U$로 바꾼다. 따라서 $\mathcal{U}$가
아핀 스킴이라고 가정할 수 있다.

\medskip\noindent
$\mathcal{Y} = Y$와 $\mathcal{U} = U$가 아핀 스킴이라고 가정하자.
저우 보조정리(정리 \ref{stacks-more-morphisms-theorem-chow-finite-type})에 의해
$X$가 대수공간인 전사 고유 사상 $X \to \mathcal{X}$를 택할 수 있다.
아래에서는 $X \to Y$가 분리 사상들의 합성으로서 분리라는 사실을 사용한다.
대수공간 $W = X \times_\mathcal{X} U$를 생각하자. 사영 $W \to X$는
유한형이다. $X$를 $W \to X$의 스킴론적 상으로 바꾸어도 되므로,
$|W|$의 $|X|$에서의 상이 $|X|$에서 조밀하다고 가정할 수 있다
(여기서 $|h|$의 상이 $|\mathcal{X}|$에서 조밀하다는 사실을 사용하며,
이렇게 바꾼 뒤에도 $X \to \mathcal{X}$는 전사이다).
$A$가 분수체 $K$를 갖는 이산 값매김환인 모든 실선 가환 도식
$$
\xymatrix{
\Spec(K) \ar[r] \ar[d] & W \ar[r] & X \ar[d] \\
\Spec(A) \ar[rr] \ar@{..>}[rru] & & Y
}
$$
에 대해 도식을 가환으로 만드는 점선 화살표가 존재한다고 주장한다.
먼저 $K$를 확대체로 바꾸고 $A$를 이 확대체 안에서 $A$를 지배하는
값매김환으로 바꾼 뒤 점선 화살표가 존재함을 보이면 충분하다.
『대수공간의 사상』의 보조정리
\ref{spaces-morphisms-lemma-push-down-solution}를 보라.
보조정리의 가정에 의해 확대 $K'/K$, $A$를 지배하는 값매김환
$A' \subset K'$, 그리고 합성 $\Spec(A') \to \Spec(A) \to Y$를 올리며
합성 $\Spec(K') \to \Spec(K) \to W \to X$와 양립하는 화살표
$\Spec(A') \to \mathcal{X}$를 얻는다. $X \to \mathcal{X}$가 고유이므로
고유성의 값매김 판정법(『스택의 사상』의 보조정리
\ref{stacks-morphisms-lemma-criterion-proper})을 사용하여 다음을 찾을 수 있다.
확대 $K''/K'$, $A'$를 지배하는 값매김환 $A'' \subset K''$, 그리고
합성 $\Spec(A'') \to \Spec(A') \to \mathcal{X}$를 올리며 합성
$\Spec(K'') \to \Spec(K') \to \Spec(K) \to X$와 양립하는 사상
$\Spec(A'') \to X$이다. 그러면 $K''/K$, $A'' \subset K''$, 사상
$\Spec(A'') \to X$가 위 문제의 해이고 주장이 성립한다.

\medskip\noindent
이 주장과 대수공간의 경우의 보조정리
(『대수공간의 극한』의 보조정리
\ref{spaces-limits-lemma-refined-valuative-criterion-proper})에 의해
$X \to Y$는 고유이다. 『스택의 사상』의 보조정리
\ref{stacks-morphisms-lemma-image-proper-is-proper}에 의해 원하는 대로
$\mathcal{X} \to Y$가 고유라고 결론짓는다.
\end{proof}

\begin{lemma}
\label{stacks-more-morphisms-lemma-refined-valuative-criterion-separated}
$f : \mathcal{X} \to \mathcal{Y}$와 $h : \mathcal{U} \to \mathcal{X}$를
대수적 스택의 사상이라 하자. $\mathcal{Y}$가 국소 뇌터이고,
$f$가 국소 유한형이고 준분리이며, $h$가 유한형이고,
$|h| : |\mathcal{U}| \to |\mathcal{X}|$의 상이
$|\mathcal{X}|$에서 조밀하다고 가정하자.
$A$가 분수체 $K$를 갖는 이산 값매김환이고
$\gamma : y \circ j \to f \circ h \circ u$인 임의의 $2$-가환 도식
$$
\xymatrix{
\Spec(K) \ar[r]_-u \ar[d]_j & \mathcal{U} \ar[r]_h & \mathcal{X} \ar[d]^f \\
\Spec(A) \ar[rr]^-y \ar@{..>}[rru] & & \mathcal{Y}
}
$$
이 주어졌을 때 점선 화살표의 범주가 공집합이거나 동형류가
정확히 하나인 세토이드라고 하자. 그러면 $f$는 분리이다.
\end{lemma}

\begin{proof}
$\Delta$가 고유 사상임을 증명해야 한다. 먼저 $\Delta$가 분리라고 하자.
그러면 사상 $\mathcal{U} \to \mathcal{X}$와
$\Delta : \mathcal{X} \to \mathcal{X} \times_\mathcal{Y} \mathcal{X}$에 보조정리
\ref{stacks-more-morphisms-lemma-refined-valuative-criterion-proper}를 적용할 수 있다.
$f$가 준분리이므로 $\Delta$가 준콤팩트임에 유의하자.
물론 $\Delta$는 국소 유한형이다(모든 대각사상에 대해 참이다.
『스택의 사상』의 보조정리
\ref{stacks-morphisms-lemma-properties-diagonal}를 보라).
마지막으로 $A$가 분수체 $K$를 갖는 이산 값매김환이고
$\gamma : y \circ j \to \Delta \circ h \circ u$인 $2$-가환 도식
$$
\xymatrix{
\Spec(K) \ar[r]_-u \ar[d]_j &
\mathcal{U} \ar[r]_h &
\mathcal{X} \ar[d]^\Delta \\
\Spec(A) \ar[rr]^-y \ar@{..>}[rru] & &
\mathcal{X} \times_\mathcal{Y} \mathcal{X}
}
$$
이 주어졌다고 하자. 『스택의 사상』의 보조정리
\ref{stacks-morphisms-lemma-helper-diagonal}와 이 보조정리의 가정에 의해
유일한 점선 화살표가 존재함을 안다. 이로써 보조정리
\ref{stacks-more-morphisms-lemma-refined-valuative-criterion-proper}의 마지막 가정이 성립하며,
결과가 따른다.

\medskip\noindent
일반적인 경우에는 $\Delta$가 분리임을 증명하면 충분하다.
그러면 앞의 경우로 돌아가기 때문이다. 실제로 보조정리의 가정이
$$
\mathcal{U} \to \mathcal{X}
\quad\text{그리고}\quad
\Delta :
\mathcal{X} \to
\mathcal{X} \times_\mathcal{Y} \mathcal{X}
$$
에 대해 성립한다고 주장한다. 실제로 $\Delta$가 대수공간으로
표현가능하므로 앞 문단과 같은 도식의 점선 화살표의 범주는 세토이드이다
(예를 들어 『스택의 사상』의 보조정리
\ref{stacks-morphisms-lemma-cat-dotted-arrows}를 보라).
앞 문단의 논증은 이 범주들이 공집합이거나 동형류를 하나만 가짐을 보여준다.
따라서 $\Delta$는 분리이다.
\end{proof}

\begin{lemma}
\label{stacks-more-morphisms-lemma-valuative-criterion-universally-closed}
$f : \mathcal{X} \to \mathcal{Y}$를 대수적 스택의 사상이라 하자.
$\mathcal{Y}$가 국소 뇌터이고 $f$가 유한형이라고 가정하자.
$A$가 분수체 $K$를 갖는 이산 값매김환이고
$\gamma : y \circ j \to f \circ x$인 임의의 $2$-가환 도식
$$
\xymatrix{
\Spec(K) \ar[r]_-x \ar[d]_j & \mathcal{X} \ar[d]^f \\
\Spec(A) \ar[r]^-y & \mathcal{Y}
}
$$
이 주어졌을 때 다음이 성립한다고 하자. 체의 확대 $K'/K$와
$A$를 지배하는 값매김환 $A' \subset K'$가 존재하여 유도된 도식
$$
\xymatrix{
\Spec(K') \ar[r]_-{x'} \ar[d]_{j'} & \mathcal{X} \ar[d]^f \\
\Spec(A') \ar[r]^-{y'} \ar@{..>}[ru] & \mathcal{Y}
}
$$
및 유도된 $2$-화살표 $\gamma' : y' \circ j' \to f \circ x'$에 대한
점선 화살표의 범주는 공집합이 아니다
(『스택의 사상』의 정의
\ref{stacks-morphisms-definition-fill-in-diagram}).
그러면 $f$는 보편 닫힘이다.
\end{lemma}

\begin{proof}
$V$가 아핀 스킴인 매끄러운 사상 $V \to \mathcal{Y}$를 잡자.
점선 화살표의 범주는 밑변환에 대해 잘 행동한다
(『스택의 사상』의 보조정리
\ref{stacks-morphisms-lemma-cat-dotted-arrows-base-change}).
따라서 확대 뒤의 점선 화살표의 존재에 관한 가정은 사상
$\mathcal{X} \times_\mathcal{Y} V \to V$에 물려진다.
그러므로 사상 $\mathcal{X} \times_\mathcal{Y} V \to V$는
보조정리의 가정을 만족하며,
$\mathcal{Y}$가 아핀 스킴이라고 가정해도 된다.

\medskip\noindent
$\mathcal{Y} = Y$가 뇌터 아핀 스킴이라고 가정하자.
(이제부터 $2$-화살표 $\gamma$와 $\gamma'$는 더 이상 신경 쓰지 않아도 된다.
『스택의 사상』의 보조정리
\ref{stacks-morphisms-lemma-cat-dotted-arrows-independent}를 보라.)
$f$가 보편 닫힘임을 증명하려면 『스택의 극한』의 보조정리
\ref{stacks-limits-lemma-test-universally-closed}에 의해 모든 $n$에 대해
$|\mathcal{X} \times \mathbf{A}^n| \to |Y \times \mathbf{A}^n|$가
닫힌 사상임을 보이면 충분하다. 보조정리의 가정은 곱 사상
$\mathcal{X} \times \mathbf{A}^n \to Y \times \mathbf{A}^n$에 물려지므로
(세부사항은 생략한다), $|\mathcal{X}| \to |Y|$가 닫힌 사상임을
증명하는 것으로 환원된다.

\medskip\noindent
$Y$가 뇌터 아핀 스킴이라고 가정하자.
$T \subset |\mathcal{X}|$를 닫힌 부분집합이라 하자.
$T$의 $|Y|$에서의 상이 닫혀 있음을 보여야 한다.
$\mathcal{X}$를 $T$ 위의 유도된 축약 닫힌 부분공간 구조로 바꾸어도 된다.
이렇게 바꾸어도 점선 화살표의 존재에 관한 성질이 보존된다는 확인은 생략한다.
따라서 $|\mathcal{X}| \to |Y|$의 상이 닫혀 있음을 증명하는 것으로 환원된다.

\medskip\noindent
$y \in |Y|$를 $|\mathcal{X}| \to |Y|$의 상의 폐포에 속하는 점이라 하자.
보조정리 \ref{stacks-more-morphisms-lemma-reach-point-closure-Noetherian}에 의해 가환 도식
$$
\xymatrix{
\Spec(K) \ar[r] \ar[d] & \mathcal{X} \ar[d]^f \\
\Spec(A) \ar[r] & Y
}
$$
을 택할 수 있다. 여기서 $A$는 이산 값매김환이고 $K$는 그 분수체이며,
$\Spec(A)$의 닫힌점은 $y$로 간다. 이 보조정리의 가정에서 곧바로
$y$가 $|\mathcal{X}| \to |Y|$의 상에 속함을 얻으며, 증명이 끝난다.
\end{proof}







\section{모듈라이 공간}
\label{stacks-more-morphisms-section-moduli-spaces}

\noindent
이 절에서는 대수적 스택에서 대수공간으로 가는 사상
$f : \mathcal{X} \to Y$를 논한다. 적절한 가정 아래 $Y$를
$\mathcal{X}$의 {\it 모듈라이 공간}이라 한다.
$\mathcal{X} = [U/R]$이 표시이면 $R$-불변 사상 $U \to Y$를 얻고,
(적절한 가정 아래) $Y$는 준군 $(U, R, s, t, c)$의 {\it 몫}이다.
서로 다른 유형의 몫에 관한 논의는 『준군의 몫』의 절
\ref{groupoids-quotients-section-introduction}부터 볼 수 있다.

\begin{definition}
\label{stacks-more-morphisms-definition-categorical-quotient}
$\mathcal{X}$를 대수적 스택이라 하자.
$f : \mathcal{X} \to Y$를 대수공간 $Y$로 가는 사상이라 하자.
\begin{enumerate}
\item 대수공간 $W$로 가는 임의의 사상 $\mathcal{X} \to W$가
$f$를 통해 유일하게 분해되면 $f$를 {\it 범주적 모듈라이 공간}이라 한다.
\item 대수공간의 임의의 평탄 사상 $Y' \to Y$에 대해 밑변환
$f' : Y' \times_Y \mathcal{X} \to Y'$가 범주적 모듈라이 공간이면
$f$를 {\it 균일 범주적 모듈라이 공간}이라 한다.
\end{enumerate}
$\mathcal{C}$를 대수공간의 범주의 충만한 부분범주라 하자.
\begin{enumerate}
\item[(3)] $Y \in \Ob(\mathcal{C})$이고, $W \in \Ob(\mathcal{C})$인
임의의 사상 $\mathcal{X} \to W$가 $f$를 통해 유일하게 분해되면
$f$를 {\it $\mathcal{C}$에서의 범주적 모듈라이 공간}이라 한다.
\item[(4)] $Y \in \Ob(\mathcal{C})$이고, $\mathcal{C}$의 모든 평탄 사상
$Y' \to Y$에 대해 밑변환 $f' : Y' \times_Y \mathcal{X} \to Y'$가
$\mathcal{C}$에서의 범주적 모듈라이 공간이면 이를
{\it $\mathcal{C}$에서의 균일 범주적 모듈라이 공간}이라 한다.
\end{enumerate}
\end{definition}

\noindent
요네다 보조정리에 의해 범주적 모듈라이 공간이 존재한다면 유일하다.
이를 몫에 대해 도입한 언어와 맞추어 보자.

\begin{lemma}
\label{stacks-more-morphisms-lemma-quotient-compare}
$(U, R, s, t, c)$를 대수공간의 준군이라 하고
$s, t : R \to U$가 평탄하고 국소 유한 표시라고 하자.
대수적 스택 $\mathcal{X} = [U/R]$을 생각하자.
대수공간 $Y$가 주어지면, 사상 $f : \mathcal{X} \to Y$와
$R$-불변 사상 $\phi : U \to Y$ 사이에 $1$-대-$1$ 대응이 있다.
\end{lemma}

\begin{proof}
『표현가능성의 판정 조건』의 정리
\ref{criteria-theorem-flat-groupoid-gives-algebraic-stack}에 의해
$\mathcal{X}$는 대수적 스택이다. 사상 $f : \mathcal{X} \to Y$가 주어지면
$\phi : U \to Y$를 합성 $U \to \mathcal{X} \to Y$라 하자.
$R = U \times_\mathcal{X} U$이므로
(『대수공간의 준군』의 보조정리
\ref{spaces-groupoids-lemma-quotient-stack-2-cartesian}),
$\phi$가 $R$-불변임은 즉시 알 수 있다. 반대로 대수공간으로 가는
$R$-불변 사상 $\phi : U \to Y$가 주어지면 『대수공간의 준군』의 보조정리
\ref{spaces-groupoids-lemma-quotient-stack-2-coequalizer}에 의해 사상
$f : \mathcal{X} \to Y$를 얻는다. 『대수공간의 준군』의 보조정리
\ref{spaces-groupoids-lemma-quotient-stack-objects}에 나오는 몫 스택의
명시적 서술을 이용하여 $\phi$로부터 $f$를 구성할 수도 있다.
\end{proof}

\begin{lemma}
\label{stacks-more-morphisms-lemma-categorical-quotient-compare}
보조정리 \ref{stacks-more-morphisms-lemma-quotient-compare}의 가정과 표기법을 따르자.
$f$가 (균일) 범주적 모듈라이 공간일 필요충분조건은 $\phi$가
(균일) 범주적 몫인 것이다. 충만한 부분범주 안의 모듈라이 공간에
대해서도 마찬가지이다.
\end{lemma}

\begin{proof}
보조정리 \ref{stacks-more-morphisms-lemma-quotient-compare}에서 확립한 $1$-대-$1$ 대응으로부터,
$f$가 범주적 모듈라이 공간일 필요충분조건이 $\phi$가 범주적 몫인 것임을
즉시 알 수 있다(『준군의 몫』의 정의
\ref{groupoids-quotients-definition-categorical}).
$Y' \to Y$가 사상이면
$U' = Y' \times_Y U \to Y' \times_Y \mathcal{X} = \mathcal{X}'$는
$U \to \mathcal{X}$의 밑변환으로서 전사이고 평탄하며 국소 유한 표시이다
(『표현가능성의 판정 조건』의 보조정리
\ref{criteria-lemma-flat-quotient-flat-presentation}).
또한 섬유곱의 추이성에 의해 $R' = Y' \times_Y R$은
$U' \times_{\mathcal{X}'} U'$와 같다. 따라서
$\mathcal{X}' = [U'/R']$이다. 『대수적 스택』의 주석
\ref{algebraic-remark-flat-fp-presentation}을 보라.
그러므로 우리의 상황을 $Y'$로 밑변환하면 보조정리의 명제에서와 같은
또 하나의 상황을 얻는다. 이로부터 $f$가 균일 범주적 모듈라이 공간일
필요충분조건이 $\phi$가 균일 범주적 몫인 것임이 즉시 따른다.
\end{proof}

\begin{lemma}
\label{stacks-more-morphisms-lemma-check-uniform-categorical-quotient-on-affines}
$f : \mathcal{X} \to Y$를 대수적 스택에서 대수공간으로 가는 사상이라 하자.
모든 아핀 스킴 $Y'$와 평탄 사상 $Y' \to Y$에 대해 밑변환
$f' : Y' \times_Y \mathcal{X} \to Y'$가 범주적 모듈라이 공간이면,
$f$는 균일 범주적 모듈라이 공간이다.
\end{lemma}

\begin{proof}
$Y_i$가 아핀 스킴인 에탈 덮개 $\{Y_i \to Y\}$를 택한다.
각 $i$와 $j$에 대해 아핀 열린 덮개
$Y_i \times_Y Y_j = \bigcup Y_{ijk}$를 택한다.
$\mathcal{X}_i = Y_i \times_Y \mathcal{X}$와
$\mathcal{X}_{ijk} = Y_{ijk} \times_Y \mathcal{X}$로 두자.
$g : \mathcal{X} \to W$를 대수공간으로 가는 사상이라 하자.
도식
$$
\xymatrix{
\mathcal{X}_i \ar[r] \ar[d] & \mathcal{X} \ar[d] \ar[r]_g & W \\
Y_i \ar[r] \ar@{..>}[rru] & Y
}
$$
을 생각한다. $\mathcal{X}_i \to Y_i$가 범주적 모듈라이 공간이라는 가정은
유일한 점선 화살표 $h_i : Y_i \to W$를 만들어 낸다.
$\mathcal{X}_{ijk} \to Y_{ijk}$가 범주적 모듈라이 공간이라는 가정에 의해
$h_i$와 $h_j$를 $Y_{ijk}$로 제한하면 같다. 따라서 $h_i$와 $h_j$는
$Y_i \times_Y Y_j$에서 일치한다. 『대수공간』의 절
\ref{spaces-section-presentations}에 의해
$Y = \coprod Y_i / \coprod Y_i \times_Y Y_j$이므로, $g$가 그를 통해
분해되는 유일한 사상 $Y \to W$가 존재한다고 결론짓는다.
따라서 $f$는 범주적 모듈라이 공간이다. 평탄 밑변환 뒤에도 같은 논증을
적용할 수 있으므로 $f$는 균일 범주적 모듈라이 공간이다.
\end{proof}





\section{킬--모리 정리}
\label{stacks-more-morphisms-section-Keel-Mori}

\noindent
이 절에서는 대수적 스택의 맥락에서 킬과 모리의 정리를 논하기 시작한다.
관련 문헌에 관한 논의는 『문헌 안내』의 소절
\ref{guide-subsection-coarse-moduli-spaces}을 보라.

\begin{definition}
\label{stacks-more-morphisms-definition-well-nigh-affine}
$\mathcal{X}$를 대수적 스택이라 하자. 아핀 스킴 $U$와 전사이고
평탄하며 유한이고 유한 표시인 사상 $U \to \mathcal{X}$가 존재하면
$\mathcal{X}$가 {\it 거의 아핀}이라고 한다.
\end{definition}

\noindent
이 성질을 자주 사용할 생각이 없으므로 다소 우스꽝스러운 이름을 붙였다.

\begin{lemma}
\label{stacks-more-morphisms-lemma-well-nigh-affine}
$\mathcal{X}$를 대수적 스택이라 하자. 다음은 동치이다.
\begin{enumerate}
\item $\mathcal{X}$는 거의 아핀이다.
\item $U$와 $R$이 아핀이고 $s, t : R \to U$가 유한 국소 자유이며
$\mathcal{X} = [U/R]$인 준군 스킴 $(U, R, s, t, c)$가 존재한다.
\end{enumerate}
이들이 성립하면 $\mathcal{X}$는 준콤팩트이고 준-DM이며 분리이다.
\end{lemma}

\begin{proof}
$\mathcal{X}$가 거의 아핀이라고 하자. 아핀 스킴 $U$와 전사이고
평탄하며 유한이고 유한 표시인 사상 $U \to \mathcal{X}$를 택한다.
$R = U \times_\mathcal{X} U$로 두자. 그러면 대수공간의 준군
$(U, R, s, t, c)$와 동형사상 $[U/R] \to \mathcal{X}$를 얻는다.
『대수적 스택』의 보조정리
\ref{algebraic-lemma-map-space-into-stack}와 주석
\ref{algebraic-remark-flat-fp-presentation}을 보라.
$s, t : R \to U$는 $U \to \mathcal{X})$의 밑변환으로서 평탄하고 유한이며
유한 표시이다. 따라서 $s, t$는 유한 국소 자유이다
(『사상』의 보조정리 \ref{morphisms-lemma-finite-flat}).
유한 사상은 아핀이므로 이는 $R$이 아핀임을 뜻하고, 따라서 (2)가 성립한다.

\medskip\noindent
$U$와 $R$이 아핀이고 $s, t : R \to U$가 유한 국소 자유인
준군 스킴 $(U, R, s, t, c)$가 주어졌다고 하자.
$\mathcal{X} = [U/R]$로 두자. 『표현가능성의 판정 조건』의 정리
\ref{criteria-theorem-flat-groupoid-gives-algebraic-stack}에 의해
$\mathcal{X}$는 대수적 스택이다(엄밀히는 여기서 이 사실이 필요하지 않지만,
이것이 참이라는 점은 아무리 강조해도 지나치지 않다).
『표현가능성의 판정 조건』의 보조정리
\ref{criteria-lemma-flat-quotient-flat-presentation}에 의해
$U \to \mathcal{X}$는 전사이고 평탄하며 국소 유한 표시이다.
따라서 $U \to \mathcal{X}$가 유한인지 여부는 사영
$U \times_\mathcal{X} U \to U$가 이 성질을 갖는지 확인하여 판별할 수 있다.
『스택의 성질』의 보조정리
\ref{stacks-properties-lemma-check-property-covering}를 보라.
『대수공간의 준군』의 보조정리
\ref{spaces-groupoids-lemma-quotient-stack-2-cartesian}에 의해
$U \times_\mathcal{X} U = R$이므로 이는 참이다.
따라서 $\mathcal{X}$는 거의 아핀이다.

\medskip\noindent
마지막 명제의 증명. 『스택의 성질』의 보조정리
\ref{stacks-properties-lemma-quasi-compact-stack}에 의해
$\mathcal{X}$는 준콤팩트이다. 『스택의 사상』의 보조정리
\ref{stacks-morphisms-lemma-properties-diagonal-from-presentation}에 의해
$\mathcal{X} = [U/R]$은 준-DM이고 분리이다.
\end{proof}

\begin{lemma}
\label{stacks-more-morphisms-lemma-affine-over-well-nigh-affine}
대수적 스택 $\mathcal{X}$가 거의 아핀이라고 하자.
\begin{enumerate}
\item $\mathcal{X}$가 대수공간이면 아핀이다.
\item $\mathcal{X}' \to \mathcal{X}$가 대수적 스택의 아핀 사상이면
$\mathcal{X}'$는 거의 아핀이다.
\end{enumerate}
\end{lemma}

\begin{proof}
(1)은 『대수공간의 극한』의 보조정리
\ref{spaces-limits-lemma-affine}에서 즉시 따른다.
그러나 이는 필요 이상으로 강한 결과를 사용하는 것이다. (1)은
보조정리 \ref{stacks-more-morphisms-lemma-well-nigh-affine}와 『준군』의 명제
\ref{groupoids-proposition-finite-flat-equivalence}를 합쳐서도 따른다.

\medskip\noindent
(2)를 증명하기 위해 아핀 스킴 $U$와 전사이고 평탄하며 유한이고
유한 표시인 사상 $U \to \mathcal{X}$를 택한다.
그러면 $U' = \mathcal{X}' \times_\mathcal{X} U$는 $U$로 가는 아핀 사상을
갖는다(『스택의 사상』의 보조정리
\ref{stacks-morphisms-lemma-base-change-affine}).
따라서 $U'$는 아핀 스킴이다. 물론 $U' \to \mathcal{X}'$는
$U \to \mathcal{X}$의 밑변환으로서 전사이고 평탄하며 유한이고
유한 표시이다.
\end{proof}

\begin{lemma}
\label{stacks-more-morphisms-lemma-well-nigh-affine-moduli-space}
대수적 스택 $\mathcal{X}$가 거의 아핀이라고 하자.
아핀 스킴의 범주에서 균일 범주적 모듈라이 공간
$$
f : \mathcal{X} \longrightarrow M
$$
이 존재한다. 더욱이 $f$는 분리이고 준콤팩트이며 보편 위상동형이다.
\end{lemma}

\begin{proof}
보조정리 \ref{stacks-more-morphisms-lemma-well-nigh-affine}에서와 같이
$\mathcal{X} = [U/R]$로 쓰며 준군을 $(U, R, s, t, c)$라 하자.
$C$를 $U$ 위의 $R$-불변 함수들의 환이라 하자.
『준군』의 절 \ref{groupoids-section-finite-flat}를 보라.
$M = \Spec(C)$로 둔다. $R$-불변 사상 $U \to M$은 보조정리
\ref{stacks-more-morphisms-lemma-quotient-compare}에 의해 사상 $f : \mathcal{X} \to M$에 대응한다.
『스킴』의 보조정리 \ref{schemes-lemma-morphism-into-affine}에서 주어진
아핀 스킴으로 가는 사상의 특징짓기는 $\phi : U \to M$이 아핀 스킴의
범주에서 범주적 몫임을 즉시 보장한다. 따라서 $f$는 아핀 스킴의
범주에서 범주적 모듈라이 공간이다
(보조정리 \ref{stacks-more-morphisms-lemma-categorical-quotient-compare}).

\medskip\noindent
보조정리 \ref{stacks-more-morphisms-lemma-well-nigh-affine}에 의해 $\mathcal{X}$가 분리이므로,
『스택의 사상』의 보조정리
\ref{stacks-morphisms-lemma-compose-after-separated}에 의해 $f$는 분리이다.

\medskip\noindent
$U \to \mathcal{X}$가 전사이고 $U \to M$이 준콤팩트이므로,
『스택의 사상』의 보조정리
\ref{stacks-morphisms-lemma-surjection-from-quasi-compact}에 의해
$f$는 준콤팩트이다.

\medskip\noindent
『준군』의 보조정리 \ref{groupoids-lemma-integral-over-invariants}에 의해 합성
$$
U \to \mathcal{X} \to M
$$
은 아핀 스킴의 정수적 사상이다. 특히 이는 보편 닫힘이다
(『사상』의 보조정리 \ref{morphisms-lemma-integral-universally-closed}).
$U \to \mathcal{X}$가 전사이므로 $\mathcal{X} \to M$도 보편 닫힘이다
(『스택의 사상』의 보조정리
\ref{stacks-morphisms-lemma-image-proper-is-proper}).
$\mathcal{X} \to M$이 보편 위상동형이라고 결론짓기 위해서는
보편 전단사, 즉 전사이고 보편 단사임을 보이면 충분하다.

\medskip\noindent
『스택의 사상』의 보조정리
\ref{stacks-morphisms-lemma-points-presentation}에 의해
$|\mathcal{X}| = |U|/|R|$이다. 따라서 『준군』의 보조정리
\ref{groupoids-lemma-points}에 의해 $|f|$는 전사이고 나아가 전단사이다.

\medskip\noindent
$C \to C'$를 환 사상이라 하자. $(U', R', s', t', c')$를
$M' = \Spec(C') \to M$에 의한 $(U, R, s, t, c)$의 밑변환이라 하자.
$\mathcal{X}' = [U'/R']$로 두면 『준군의 몫』의 보조정리
\ref{groupoids-quotients-lemma-base-change-quotient-stack}에 의해
$M' \times_M \mathcal{X} = \mathcal{X}'$이다.
$C^1$을 $U'$ 위의 $R'$-불변 함수들의 환이라 하자.
$M^1 = \Spec(C^1)$로 두고 도식
$$
\xymatrix{
\mathcal{X}' \ar[d]^{f'} \ar[r] & \mathcal{X} \ar[dd]^f \\
M^1 \ar[d] \\
M' \ar[r] & M
}
$$
을 생각하자. 『준군』의 보조정리
\ref{groupoids-lemma-invariants-base-change}와 『대수학』의 보조정리
\ref{algebra-lemma-universally-bijective}에 의해 사상
$M^1 \to M'$는 위상동형이다. 한편 앞 문단을
$(U', R', s', t', c')$에 적용하면 $|f'|$가 전단사임을 알 수 있다.
따라서 $f$는 아핀 스킴에 의한 임의의 밑변환 뒤에도 점들 사이의
전단사를 유도한다. 그러므로 『스택의 사상』의 보조정리
\ref{stacks-morphisms-lemma-universally-injective-local}에 의해
$f$는 보편 단사이다.

\medskip\noindent
마지막으로 $f$가 아핀 스킴의 범주에서 균일 모듈라이 공간임을
보여야 한다. 이는 위의 논의와, 환 사상 $C \to C'$가 평탄이면
『준군』의 보조정리 \ref{groupoids-lemma-invariants-base-change}에 의해
$C' \to C^1$이 동형사상이라는 사실에서 따른다.
\end{proof}

\begin{lemma}
\label{stacks-more-morphisms-lemma-well-nigh-affine-moduli-space-etale}
$h : \mathcal{X}' \to \mathcal{X}$를 대수적 스택의 사상이라 하자.
$\mathcal{X}'$와 $\mathcal{X}$가 거의 아핀이고, $h$가 에탈이며,
$h$가 자기동형군 위의 동형사상을 유도한다고 가정하자
(『스택의 사상』의 주석
\ref{stacks-morphisms-remark-identify-automorphism-groups}).
그러면 카르테시안 도식
$$
\xymatrix{
\mathcal{X}' \ar[d] \ar[r] & \mathcal{X} \ar[d] \\
M' \ar[r] & M
}
$$
이 존재한다. 여기서 $M' \to M$은 에탈이고, 세로 화살표들은
보조정리 \ref{stacks-more-morphisms-lemma-well-nigh-affine-moduli-space}에서 구성한
모듈라이 공간이다.
\end{lemma}

\begin{proof}
『스택의 사상』의 보조정리
\ref{stacks-morphisms-lemma-aut-iso-unramified}와
\ref{stacks-morphisms-lemma-stabilizer-preserving}에 의해
$h$는 대수공간으로 표현가능하다. 아핀 스킴 $U$와 전사이고 평탄하며
유한이고 유한 표시인 사상 $U \to \mathcal{X}$를 택한다.
그러면 $U' = \mathcal{X}' \times_\mathcal{X} U$는 대수공간이고
유한(특히 아핀) 사상 $U' \to \mathcal{X}'$을 갖는다.
보조정리 \ref{stacks-more-morphisms-lemma-affine-over-well-nigh-affine}에 의해 $U'$는 아핀이다.
$R = U \times_\mathcal{X} U$와
$R' = U' \times_{\mathcal{X}'} U'$로 두면
$\mathcal{X} = [U/R]$이고 $\mathcal{X}' = [U'/R']$인 준군
$(U, R, s, t, c)$와 $(U', R', s', t', c')$를 얻는다.
보조정리 \ref{stacks-more-morphisms-lemma-well-nigh-affine}의 증명을 보라.
다음 도식들이 카르테시안임을 알 수 있다.
$$
\xymatrix{
R' \ar[d]_{s'} \ar[r]_f & R \ar[d]^s \\
U' \ar[r]^f & U
}
\quad
\quad
\xymatrix{
R' \ar[d]_{t'} \ar[r]_f & R \ar[d]^t \\
U' \ar[r]^f & U
}
\quad
\quad
\xymatrix{
G' \ar[d] \ar[r]_f & G \ar[d] \\
U' \ar[r]^f & U
}
$$
여기서 $G$와 $G'$는 안정자 군스킴이다. 앞의 두 도식은 섬유곱의
추이성에서 따르고, 마지막 도식은 동형사상
$\mathcal{I}_{\mathcal{X}'} \to
\mathcal{X}' \times_\mathcal{X} \mathcal{I}_\mathcal{X}$의 당김이므로
카르테시안이다(이미 사용한 『스택의 사상』의 보조정리
\ref{stacks-morphisms-lemma-aut-iso-unramified}).
$M$ 및 $M'$는 보조정리 \ref{stacks-more-morphisms-lemma-well-nigh-affine-moduli-space}에서
각각 $U$ 위의 $R$-불변 함수들의 환 및 $U'$ 위의 $R'$-불변 함수들의
환의 스펙트럼으로 구성했음을 상기하자. 따라서 『준군』의 보조정리
\ref{groupoids-lemma-etale}에 의해 $M' \to M$은 에탈이고
$U' = M' \times_M U$이다. 그 결과 $R' = M' \times_M U$이다.
달리 말하면 준군 $(U', R', s', t', c')$는
$M' \to M$에 의한 $(U, R, s, t, c)$의 밑변환이다.
『준군의 몫』의 보조정리
\ref{groupoids-quotients-lemma-base-change-quotient-stack}에 의해
이는 보조정리의 도식이 카르테시안임을 뜻한다.
\end{proof}

\begin{lemma}
\label{stacks-more-morphisms-lemma-moduli-space-finite-affine}
대수적 스택 $\mathcal{X}$가 거의 아핀이라고 하자.
보조정리 \ref{stacks-more-morphisms-lemma-well-nigh-affine-moduli-space}의 사상
$$
f : \mathcal{X} \longrightarrow M
$$
은 균일 범주적 모듈라이 공간이다.
\end{lemma}

\begin{proof}
$M$이 아핀 스킴의 범주에서 균일 범주적 모듈라이 공간이라는 것은
이미 안다. 보조정리
\ref{stacks-more-morphisms-lemma-check-uniform-categorical-quotient-on-affines}에 의해,
아핀 스킴의 임의의 평탄 사상 $M' \to M$에 대해 밑변환
$f' : M' \times_M \mathcal{X} \to M'$가 범주적 모듈라이 공간임을
보이면 충분하다. 보조정리
\ref{stacks-more-morphisms-lemma-affine-over-well-nigh-affine}에 의해
$\mathcal{X}' = M' \times_M \mathcal{X}$는 거의 아핀이다.
따라서 $\mathcal{X}$를 $\mathcal{X}'$로, $M$을 $M'$로 바꾸면
$f$가 범주적 모듈라이 공간임을 증명하는 것으로 환원된다.

\medskip\noindent
$g : \mathcal{X} \to Y$를 대수공간 $Y$로 가는 사상이라 하자.
유일한 사상 $h : M \to Y$에 대해 $g = h \circ f$임을 보여야 한다.

\medskip\noindent
유일성. $g = h_1 \circ f = h_2 \circ f$인 두 사상
$h_i : M \to Y$가 있다고 하자. $M' \subset M$을 $h_1$과 $h_2$의
등화자라 하자. 그러면 $M' \to M$은 단사상이고
$f : \mathcal{X} \to M$은 $M'$를 통해 분해된다. 따라서
$M' \to M$은 보편 위상동형이다. 『사상』의 보조정리
\ref{morphisms-lemma-universal-homeomorphism}에 의해 $M'$은 아핀이다.
그런데 $f : \mathcal{X} \to M$은 아핀 스킴의 범주에서 범주적
모듈라이 공간이므로 $M' = M$임을 알 수 있다.

\medskip\noindent
존재성. 아래에서 모든 $p \in M$에 대해 카르테시안 정사각형
$$
\xymatrix{
\mathcal{X}' \ar[r] \ar[d] & \mathcal{X} \ar[d] \\
M' \ar[r] & M
}
$$
이 존재함을 보일 것이다. 여기서 $M' \to M$은 $p$를 그 상에 포함하는
아핀들 사이의 에탈 사상이고, 합성 $\mathcal{X}' \to \mathcal{X} \to Y$는
$M'$를 통해 분해된다. 이는 사상 $h : M \to Y$를 $M$ 위에서
에탈 국소적으로 구성할 수 있다는 뜻이다. $Y$가 에탈 위상에 대한 층이고
위에서 유일성을 보였으므로 이것으로 충분하다(작은 세부사항은 생략한다).

\medskip\noindent
$y \in |Y|$를 $p$의 상이라 하자. $V$가 아핀인 에탈 사상
$(V, v) \to (Y, y)$를 택한다. $\mathcal{X}' = V \times_Y \mathcal{X}$를
생각하자. $\mathcal{X}' \to \mathcal{X}$는 $V \to Y$의 밑변환으로서
분리이고 에탈이다. 더욱이 $V \to Y$가 그러하므로
$\mathcal{X}' \to \mathcal{X}$는 자기동형군 위의 동형사상을 유도한다
(『스택의 사상』의 주석
\ref{stacks-morphisms-remark-identify-automorphism-groups}과 보조정리
\ref{stacks-morphisms-lemma-base-change-stabilizer-preserving}를 보라).
보조정리 \ref{stacks-more-morphisms-lemma-well-nigh-affine}에서와 같은 표시
$\mathcal{X} = [U/R]$을 택한다. $U' = \mathcal{X}' \times_\mathcal{X} U = V \times_Y U$로 두고, $p$와 $v$로 가는 점 $u' \in U'$를 택한다
(『대수공간의 성질』의 보조정리
\ref{spaces-properties-lemma-points-cartesian}에 의해 가능하다).
$U' \to U$가 분리이고 에탈이므로 $U'$의 모든 유한 점 집합은
어떤 아핀 열린집합에 포함된다. 『사상에 관한 추가 내용』의 보조정리
\ref{more-morphisms-lemma-separated-locally-quasi-finite-over-affine}를 보라.
한편 사상 $U' \to \mathcal{X}'$는 전사이고 유한이며 평탄하고
국소 유한 표시이다. $R' = U' \times_{\mathcal{X}'} U'$로 두면
$s', t' : R' \to U'$는 유한 국소 자유이다. 『준군』의 보조정리
\ref{groupoids-lemma-find-invariant-affine}에 의해 $u'$를 포함하는
$R'$-불변 아핀 열린 부분스킴 $U'' \subset U'$가 존재한다.
$\mathcal{X}'' \subset \mathcal{X}'$을 대응하는 열린 부분스택이라 하자.
그러면 $\mathcal{X}''$는 거의 아핀이다. 보조정리
\ref{stacks-more-morphisms-lemma-well-nigh-affine-moduli-space-etale}에 의해 카르테시안 정사각형
$$
\xymatrix{
\mathcal{X}'' \ar[r] \ar[d] & \mathcal{X} \ar[d] \\
M'' \ar[r] & M
}
$$
을 얻으며, 여기서 $M'' \to M$은 에탈이다.
$\mathcal{X}'' \to M''$가 아핀 스킴의 범주에서 범주적 모듈라이 공간이므로
합성 $\mathcal{X}'' \to \mathcal{X}' \to V$가 합성
$\mathcal{X}'' \to M'' \to V$와 같게 하는 사상 $M'' \to V$를 얻는다.
이로써 주장이 증명되고 증명이 끝난다.
\end{proof}

\begin{lemma}
\label{stacks-more-morphisms-lemma-etale-separated-over-well-nigh-affine}
$h : \mathcal{X}' \to \mathcal{X}$를 대수적 스택의 사상이라 하자.
$\mathcal{X}$가 거의 아핀이고, $h$가 에탈이고 $h$가 분리이며,
$h$가 자기동형군 위의 동형사상을 유도한다고 가정하자
(『스택의 사상』의 주석
\ref{stacks-morphisms-remark-identify-automorphism-groups}).
그러면 카르테시안 도식
$$
\xymatrix{
\mathcal{X}' \ar[d] \ar[r] & \mathcal{X} \ar[d] \\
M' \ar[r] & M
}
$$
이 존재한다. 여기서 $M' \to M$은 스킴의 분리 에탈 사상이고,
$\mathcal{X} \to M$은 보조정리
\ref{stacks-more-morphisms-lemma-well-nigh-affine-moduli-space}에서 구성한 모듈라이 공간이다.
\end{lemma}

\begin{proof}
아핀 스킴 $U$와 전사이고 평탄하며 유한이고 국소 유한 표시인 사상
$U \to \mathcal{X}$를 택한다. 『스택의 사상』의 보조정리
\ref{stacks-morphisms-lemma-aut-iso-unramified}와
\ref{stacks-morphisms-lemma-stabilizer-preserving}에 의해 $h$가
대수공간으로 표현가능하므로
$U' = \mathcal{X}' \times_\mathcal{X} U$는 대수공간이다.
$U' \to U$가 분리이고 에탈이므로 $U'$는 스킴이고, $U'$의 모든
유한 점 집합은 어떤 아핀 열린집합에 포함된다. 『대수공간의 사상』의 보조정리
\ref{spaces-morphisms-lemma-locally-quasi-finite-separated-representable}와
『사상에 관한 추가 내용』의 보조정리
\ref{more-morphisms-lemma-separated-locally-quasi-finite-over-affine}를 보라.
$R' = U' \times_{\mathcal{X}'} U'$로 두면
$s', t' : R' \to U'$는 유한 국소 자유이다. 『준군』의 보조정리
\ref{groupoids-lemma-find-invariant-affine}에 의해
$R'$-불변 아핀 열린 부분스킴 $U'_i \subset U'$들로 이루어진
열린 덮개 $U' = \bigcup U'_i$가 존재한다.
$\mathcal{X}'_i \subset \mathcal{X}'$을 대응하는 열린 부분스택이라 하자.
$U'_i \to \mathcal{X}'_i$가 전사이고 평탄하며 유한이고 유한 표시이므로
이들은 거의 아핀이다. 보조정리
\ref{stacks-more-morphisms-lemma-well-nigh-affine-moduli-space-etale}에 의해 카르테시안 도식
$$
\xymatrix{
\mathcal{X}'_i \ar[r] \ar[d] & \mathcal{X} \ar[d] \\
M'_i \ar[r] & M
}
$$
을 얻는다. 여기서 $M'_i \to M$은 아핀 스킴의 에탈 사상이고
세로 화살표들은 보조정리
\ref{stacks-more-morphisms-lemma-well-nigh-affine-moduli-space}에서와 같다.
$\mathcal{X}'_{ij} = \mathcal{X}'_i \cap \mathcal{X}'_j$는
$\mathcal{X}'_i$와 $\mathcal{X}'_j$의 열린 부분공간임에 유의하자.
따라서 대응하는 열린 부분스킴
$V_{ij} \subset M'_i$와 $V_{ji} \subset M'_j$를 얻는다.
보조정리 \ref{stacks-more-morphisms-lemma-moduli-space-finite-affine}의 결과에 의해
$\mathcal{X}'_{ij} \to V_{ij}$와
$\mathcal{X}'_{ji} \to V_{ji}$는 둘 다 범주적 모듈라이 공간이다!
따라서 도식
$$
\xymatrix{
\mathcal{X}'_i \ar[d] & &
\mathcal{X}'_i \cap \mathcal{X}'_j \ar[rr] \ar[ll] \ar[ld] \ar[rd] & &
\mathcal{X}'_j \ar[d] \\
M'_i &
V_{ij} \ar[l] \ar[rr]^{\varphi_{ij}} & &
V_{ji} \ar[r] &
M'_j
}
$$
을 가환으로 만드는 유일한 동형사상
$\varphi_{ij} : V_{ij} \to V_{ji}$를 얻는다. 이 동형사상들은
계산(그리고 앞 보조정리의 또 한 번의 적용)에 의해 『스킴』의 절
\ref{schemes-section-glueing-schemes}의 코사이클 조건을 만족하며,
계산은 생략한다. 따라서 『스킴』의 보조정리
\ref{schemes-lemma-glue}에 의해 아핀 스킴들을 스킴 $M'$으로 붙일 수 있다.
$M'_i$를 $M'$ 안에서의 상과 동일시하자. 카르테시안 도식
$$
\xymatrix{
\mathcal{X}'_i \ar[r] \ar[d] & \mathcal{X}' \ar[d] \\
M'_i \ar[r] & M'
}
$$
들에 들어맞는 사상 $\mathcal{X}' \to M'$가 존재한다고 주장한다.
이는 『스킴』의 보조정리 \ref{schemes-lemma-glue}에서 붙인 스킴
$M'$으로 가는 사상들을 서술한 방식과, 사상
$\mathcal{X}' \to M'$를 주는 것은 임의의 사상
$T \to \mathcal{X}'$에 대해 사상 $T \to M'$를 주는 것과 같다는
사실에서 명백하다. 마찬가지로 각 $M'_i$ 위에서 주어진 사상
$M'_i \to M$으로 제한되는 사상 $M' \to M$이 존재한다.
$M' \to M$은 에탈 덮개의 각 성분 위에서 에탈이므로 에탈이고,
섬유곱 성질은 (아핀) 열린 덮개 $M' = \bigcup M'_i$의 성분들 위에서
확인할 수 있으므로 성립한다. 마지막으로 합성
$U' \to \mathcal{X}' \to M'$가 전사이고 보편 닫힘이므로
『사상』의 보조정리
\ref{morphisms-lemma-image-universally-closed-separated}를 적용하여
$M' \to M$이 분리임을 얻는다.
\end{proof}

\begin{lemma}
\label{stacks-more-morphisms-lemma-etale-local-finite-inertia}
$\mathcal{X}$를 대수적 스택이라 하자.
$\mathcal{I}_\mathcal{X} \to \mathcal{X}$가 유한이라고 가정하자.
그러면 집합 $I$와 각 $i \in I$에 대해 대수적 스택의 사상
$$
g_i : \mathcal{X}_i \longrightarrow \mathcal{X}
$$
이 존재하여 다음 성질들을 만족한다.
\begin{enumerate}
\item $|\mathcal{X}| = \bigcup |g_i|(|\mathcal{X}_i|)$이다.
\item $\mathcal{X}_i$는 거의 아핀이다.
\item $\mathcal{I}_{\mathcal{X}_i} \to
\mathcal{X}_i \times_\mathcal{X} \mathcal{I}_\mathcal{X}$는 동형사상이다.
\item $g_i : \mathcal{X}_i \to \mathcal{X}$는 대수공간으로 표현가능하고
분리이며 에탈이다.
\end{enumerate}
\end{lemma}

\begin{proof}
임의의 $x \in |\mathcal{X}|$에 대해 『스택의 사상』의 보조정리
\ref{stacks-morphisms-lemma-etale-local-quasi-DM-at-x}에서와 같은
$g : \mathcal{U} \to \mathcal{X}$, $\mathcal{U} = [U/R]$, 그리고 $u$를
택할 수 있다. 그러면 『스택의 사상』의 보조정리
\ref{stacks-morphisms-lemma-stabilizer-preserving-unramified}에 의해
$u$를 포함하는 열린 부분스택 $\mathcal{U}' \subset \mathcal{U}$가 존재하여
$\mathcal{I}_{\mathcal{U}'} \to
\mathcal{U}' \times_\mathcal{X} \mathcal{I}_\mathcal{X}$
가 동형사상이다. $U' \subset U$를 열린 부분스택
$\mathcal{U}'$에 대응하는 $R$-불변 열린집합이라 하자.
$u$로 가는 $U'$의 점을 $u' \in U'$라 하자.
$s : R \to U$가 유한이므로 $t(s^{-1}(\{u'\}))$는 유한임에 유의하자.
『성질』의 보조정리 \ref{properties-lemma-ample-finite-set-in-affine}와
『준군』의 보조정리 \ref{groupoids-lemma-find-invariant-affine}에 의해
$u'$를 포함하는 $R$-불변 아핀 열린집합 $U'' \subset U'$를 찾을 수 있다.
$R''$을 $R$의 $U''$로의 제한이라 하자. 그러면
$\mathcal{U}'' = [U''/R'']$은 $u$를 포함하는 $\mathcal{U}'$의
열린 부분스택이고 거의 아핀이며,
$\mathcal{I}_{\mathcal{U}''} \to
\mathcal{U}'' \times_\mathcal{X} \mathcal{I}_\mathcal{X}$는 동형사상이고,
$\mathcal{U}'' \to \mathcal{X}$는 대수공간으로 표현가능하고 에탈이다.
마지막으로 $\mathcal{U}''$가 분리이고
(보조정리 \ref{stacks-more-morphisms-lemma-well-nigh-affine}), $\mathcal{X}$의 대각사상이
분리이므로(『스택의 사상』의 보조정리
\ref{stacks-morphisms-lemma-diagonal-diagonal}), 『스택의 사상』의 보조정리
\ref{stacks-morphisms-lemma-compose-after-separated}에 의해
$\mathcal{U}'' \to \mathcal{X}$는 분리이다.
점 $x \in |\mathcal{X}|$가 임의였으므로 증명이 끝난다.
\end{proof}

\begin{theorem}[Keel--Mori]
\label{stacks-more-morphisms-theorem-keel-mori}
$\mathcal{X}$를 대수적 스택이라 하자.
$\mathcal{I}_\mathcal{X} \to \mathcal{X}$가 유한이라고 가정하자.
그러면 균일 범주적 모듈라이 공간
$$
f : \mathcal{X} \longrightarrow M
$$
이 존재하고, $f$는 분리이고 준콤팩트이며 보편 위상동형이다.
\end{theorem}

\begin{proof}
보조정리 \ref{stacks-more-morphisms-lemma-etale-local-finite-inertia}에서와 같은 집합
$I$\footnote{집합론적 문제를 여전히 추적하고 있는 독자는
$I$가 지나치게 크지 않은지 확인해야 한다.}와 각 $i \in I$에 대한
대수적 스택의 사상 $g_i : \mathcal{X}_i \to \mathcal{X}$를 택한다.
이 보조정리에 열거된 모든 성질을 앞으로 별도의 언급 없이 사용한다.
$$
f_i : \mathcal{X}_i \to M_i
$$
를 보조정리 \ref{stacks-more-morphisms-lemma-well-nigh-affine-moduli-space}에서와 같이 잡자.
$i, j \in I$에 대해 스택
$$
\mathcal{X}_{ij} = \mathcal{X}_i \times_{g_i, \mathcal{X}, g_j} \mathcal{X}_j
$$
를 생각하자. 사영 $\mathcal{X}_{ij} \to \mathcal{X}_i$와
$\mathcal{X}_{ij} \to \mathcal{X}_j$는 『스택의 사상』의 보조정리
\ref{stacks-morphisms-lemma-base-change-separated}에 의해 분리이고,
보조정리 \ref{stacks-morphisms-lemma-base-change-etale}에 의해 에탈이며,
자기동형군 위의 동형사상을 유도한다(『스택의 사상』의 주석
\ref{stacks-morphisms-remark-identify-automorphism-groups}의 의미에서).
이는 보조정리
\ref{stacks-morphisms-lemma-base-change-stabilizer-preserving}에서 따른다.
따라서 보조정리 \ref{stacks-more-morphisms-lemma-etale-separated-over-well-nigh-affine}를
적용하여 카르테시안 정사각형들로 이루어진 가환 도식
$$
\xymatrix{
\mathcal{X}_i \ar[d]_{f_i} &
\mathcal{X}_{ij} \ar[d]_{f_{ij}} \ar[l] \ar[r] &
\mathcal{X}_j \ar[d]_{f_j} \\
M_i &
M_{ij} \ar[l] \ar[r] &
M_j
}
$$
을 찾을 수 있다. 여기서 $M_{ij} \to M_i$와 $M_{ij} \to M_j$는
스킴의 분리 에탈 사상이다. 또한 보조정리
\ref{stacks-more-morphisms-lemma-moduli-space-finite-affine}에 의해 $f_i$가 균일 범주적
몫이라는 사실도 사용한다. 다음을 주장한다.
$$
\coprod M_{ij} \longrightarrow \coprod M_i \times \coprod M_i
$$
는 에탈 동치관계이다.

\medskip\noindent
주장의 증명. $R = \coprod M_{ij}$와 $U = \coprod M_i$로 두자.
$t : R \to U$와 $s : R \to U$가 에탈이라는 것은 이미 보았다.
$\text{pr}_{13} : U \times U \times U \to U \times U$와 양립하는 사상
$c : R \times_{s, U, t} R \to R$을 구성하자. 실제로 $i, j, k \in I$에 대해
$$
\mathcal{X}_{ijk} =
\mathcal{X}_i \times_{g_i, \mathcal{X}, g_j} \mathcal{X}_j
\times_{g_j, \mathcal{X}, g_k} \mathcal{X}_k =
\mathcal{X}_{ij} \times_{\mathcal{X}_j} \mathcal{X}_{jk}
$$
를 생각한다. 앞 문단에서와 똑같이 논증하면
$M_{ijk} = M_{ij} \times_{M_j} M_{jk}$가
$\mathcal{X}_{ijk}$의 범주적 모듈라이 공간임을 알 수 있다.
특히 사영 $\mathcal{X}_{ijk} \to \mathcal{X}_{ik}$에서 오는 표준 사상
$M_{ijk} = M_{ij} \times_{M_j} M_{jk} \to M_{ik}$가 있다.
이 사상들을 합치면 사상 $c$를 얻는다. 비슷한 방식으로 대각사상
$\Delta : U \to U \times U$와 양립하는 사상 $e : U \to R$ 및
뒤집기 $U \times U \to U \times U$와 양립하는 사상 $i : R \to R$을
구성한다. $k$를 대수적으로 닫힌 체라 하자. 그러면
$$
\Mor(\Spec(k), \mathcal{X}_i) \to \Mor(\Spec(k), M_i) = M_i(k)
$$
는 동형류 위에서 전단사이고, 사상 $M' \to M$에 의한 임의의
밑변환 뒤에도 그러하다. 이는 $f_i$의 선택과 『스택의 사상』의 보조정리
\ref{stacks-morphisms-lemma-universally-injective} 및
\ref{stacks-morphisms-lemma-universally-injective-point}에서 따른다.
$2$-섬유곱의 구성에 의해 도식
$$
\xymatrix{
\Mor(\Spec(k), \mathcal{X}_{ij}) \ar[d] \ar[r] &
\Mor(\Spec(k), \mathcal{X}_j) \ar[d] \\
\Mor(\Spec(k), \mathcal{X}_i) \ar[r] &
\Mor(\Spec(k), \mathcal{X})
}
$$
은 범주들의 섬유곱이다. $g_i$의 선택에 의해 이 도식의 함자들은
자기동형군 위의 전단사를 유도한다. 따라서 이 도식은 동형류 집합들 위의
섬유곱 도식을 유도한다! 그러므로
$$
R(k) = U(k) \times_{|\Mor(\Spec(k), \mathcal{X})|} U(k)
$$
이다. 여기서 $|\Mor(\Spec(k), \mathcal{X})|$는 동형류 집합을 뜻한다.
특히 대수적으로 닫힌 임의의 체 $k$에 대해 $k$-값 점 위의 사상은
동치관계이다. 『준군』의 보조정리
\ref{groupoids-lemma-etale-equivalence-relation}에 의해 주장이 성립한다.

\medskip\noindent
$M = U/R$을 위 에탈 동치관계의 몫인 대수공간이라 하자.
『대수공간』의 정리 \ref{spaces-theorem-presentation}을 보라.
가환 도식
\begin{equation}
\label{stacks-more-morphisms-equation-fundamental-diagram}
\xymatrix{
\mathcal{X}_i \ar[r]_{g_i} \ar[d]_{f_i} & \mathcal{X} \ar[d]^f \\
M_i \ar[r] & M
}
\end{equation}
에 들어맞는 표준 사상 $f : \mathcal{X} \to M$이 있다.
실제로 이러한 사상 $f$는 임의의 스킴 $T$에 대해 밑변환과 양립하는 함자
$$
f : \Mor(T, \mathcal{X}) \longrightarrow \Mor(T, M)
$$
로 주어진다. 왼쪽의 대상 $a : T \to \mathcal{X}$를 잡자.
$T_i = \mathcal{X}_i \times_\mathcal{X} T$인 에탈 덮개
$\{T_i \to T\}$와 사상 $a_i : T_i \to \mathcal{X}_i$를 얻는다.
그러면 $b_i = f_i \circ a_i : T_i \to M_i$를 얻는다.
$T_i \times_T T_j = \mathcal{X}_{ij} \times_\mathcal{X} T$이므로
사상 $a_{ij} : T_i \times_T T_j \to \mathcal{X}_{ij}$도 얻는다.
$b_{ij} = f_{ij} \circ a_{ij}$로 두면 $b_i \times b_j$는 단사상
$M_{ij} \to M_i \times M_j$를 통해 분해된다. 따라서 사상들
$$
T_i \xrightarrow{b_i} M_i \to M
$$
은 $T_i \times_T T_j$에서 일치한다. $M$은 에탈 위상에 대한 층이므로
이 사상들은 유일한 사상 $b = f(a) : T \to M$으로 붙는다.
이 구성이 밑변환과 양립한다는 확인과 도식
(\ref{stacks-more-morphisms-equation-fundamental-diagram})들이 가환한다는 확인은 생략한다.

\medskip\noindent
주장: 도식 (\ref{stacks-more-morphisms-equation-fundamental-diagram})들은 카르테시안이다.
이를 보이기 위해 유도된 사상
$$
h_i : \mathcal{X}_i \longrightarrow M_i \times_M \mathcal{X}
$$
을 연구한다. 이는 $\mathcal{X}$ 위에서 에탈인 스택의 사상이므로
$h_i$는 에탈이다(『스택의 사상』의 보조정리
\ref{stacks-morphisms-lemma-etale-permanence}).
$g_i$가 분리이므로 $h_i$는 분리이다
(『스택의 사상』의 보조정리
\ref{stacks-morphisms-lemma-compose-after-separated}와 위에서 본
$\mathcal{X}$의 대각사상이 분리라는 사실을 사용한다).
$g_i$가 그러하므로 $h_i$는 자기동형군 위의 동형사상을 유도한다
(『스택의 사상』의 주석
\ref{stacks-morphisms-remark-identify-automorphism-groups}).
대수적으로 닫힌 체 $k$에 대해 도식
$$
\xymatrix{
\Mor(\Spec(k), M_i \times_M \mathcal{X}) \ar[r] \ar[d] &
\Mor(\Spec(k), \mathcal{X}) \ar[d] \\
M_i(k) \ar[r] &
M(k)
}
$$
은 범주들의 카르테시안 도식이고, 위쪽 화살표는 자기동형군 위의
전단사를 유도한다. 한편 위에서 말한 바에 의해
$$
M(k) = U(k)/R(k) = U(k)/
U(k) \times_{|\Mor(\Spec(k), \mathcal{X})|} U(k) =
|\Mor(\Spec(k), \mathcal{X})|
$$
이다. 따라서 위 카르테시안 도식의 오른쪽 세로 화살표는
동형류 위의 전단사이다. 그러므로
$|\Mor(\Spec(k), M_i \times_M \mathcal{X})| \to M_i(k)$가 전단사라고
결론짓는다. 정리하면, $h_i$는 분리이고 에탈이고,
자기동형군 위의 동형사상을 유도하며(『스택의 사상』의 주석
\ref{stacks-morphisms-remark-identify-automorphism-groups}의 의미에서),
대수적으로 닫힌 체 위의 섬유 범주에서 동치를 유도한다.
따라서 『스택의 사상』의 보조정리
\ref{stacks-morphisms-lemma-etale-iso}에 의해 이는 동형사상이다.

\medskip\noindent
이 주장으로부터 특히 다음을 얻는다. 전사 에탈 사상 $U \to M$이 존재하여
$f$의 밑변환이 분리이고 준콤팩트이며 보편 위상동형이다.
따라서 $f$는 분리이고 준콤팩트이며 보편 위상동형이다.
『스택의 사상』의 보조정리
\ref{stacks-morphisms-lemma-check-separated-covering},
\ref{stacks-morphisms-lemma-check-quasi-compact-covering}, 그리고
\ref{stacks-morphisms-lemma-check-universal-homeomorphism-covering}을 보라.

\medskip\noindent
증명을 끝내려면 $f : \mathcal{X} \to M$이 균일 범주적 모듈라이 공간임을
보여야 한다. 이를 위해서는 대수공간의 평탄 사상 $M' \to M$이 주어졌을 때
밑변환
$$
M' \times_M \mathcal{X} \longrightarrow M'
$$
이 범주적 모듈라이 공간임을 보이면 충분하다. 따라서 $E$가 대수공간인 사상
$$
\theta : M' \times_M \mathcal{X} \longrightarrow E
$$
을 생각하자. 각 $i$에 대해 $f_i$가 균일 범주적 모듈라이 공간임을 안다.
따라서
$$
\xymatrix{
M' \times_M \mathcal{X}_i \ar[d] \ar[r] &
M' \times_M \mathcal{X} \ar[d]^\theta \\
M' \times_M M_i \ar[r]^{\psi_i} &
E
}
$$
를 얻는다. $\{M' \times_M M_i \to M'\}$가 에탈 덮개이므로,
원하는 사상 $\psi : M' \to E$를 얻기 위해서는 $\psi_i$와 $\psi_j$가
$M' \times_M M_i \times_M M_j = M' \times_M M_{ij}$ 위에서 일치함을
보이면 충분하다. 이는
$f_{ij} : \mathcal{X}_{ij} =
\mathcal{X}_i \times_\mathcal{X} \mathcal{X}_j \to M_{ij}$가
균일 범주적 몫이라는 사실에서 쉽게 따른다. 세부사항은 생략한다.
마지막으로 $\psi$가 가환 도식
$$
\xymatrix{
M' \times_M \mathcal{X} \ar[d] \ar[rd]^\theta \\
M' \ar[r]^\psi &
E
}
$$
에 들어맞는다는 것을 보인다. 실제로
``$\{M' \times_M \mathcal{X}_i \to M' \times_M \mathcal{X}\}$는
에탈 덮개''이고, 구성에 의해 사상 $\psi_i$들은 대응하는 가환 도식들에
들어맞기 때문이다. 이로써 킬--모리 정리의 증명이 끝난다.
\end{proof}

\noindent
다음 보조정리는 이 구성이 에탈 국소적이라는 점을 강조한다.

\begin{lemma}
\label{stacks-more-morphisms-lemma-etale-separated-over-keel-mori}
$h : \mathcal{X}' \to \mathcal{X}$를 대수적 스택의 사상이라 하자.
다음을 가정하자.
\begin{enumerate}
\item $\mathcal{I}_\mathcal{X} \to \mathcal{X}$는 유한이다.
\item $h$는 에탈이고 분리이며 자기동형군 위의 동형사상을 유도한다
(『스택의 사상』의 주석
\ref{stacks-morphisms-remark-identify-automorphism-groups}).
\end{enumerate}
그러면 카르테시안 도식
$$
\xymatrix{
\mathcal{X}' \ar[d] \ar[r] &
\mathcal{X} \ar[d] \\
M' \ar[r] &
M
}
$$
이 존재한다. 여기서 $M' \to M$은 대수공간의 분리 에탈 사상이고,
세로 화살표들은 정리 \ref{stacks-more-morphisms-theorem-keel-mori}에서 구성한 모듈라이 공간이다.
\end{lemma}

\begin{proof}
『스택의 사상』의 보조정리
\ref{stacks-morphisms-lemma-aut-iso-unramified}에 의해
$\mathcal{I}_{\mathcal{X}'} \to
\mathcal{X}' \times_\mathcal{X} \mathcal{I}_\mathcal{X}$가 동형사상이다.
따라서 $\mathcal{I}_{\mathcal{X}'} \to \mathcal{X}'$는
$\mathcal{I}_\mathcal{X} \to \mathcal{X}$의 밑변환으로서 유한이다.
$f' : \mathcal{X}' \to M'$와 $f : \mathcal{X} \to M$을 정리
\ref{stacks-more-morphisms-theorem-keel-mori}에서와 같이 잡자. $f'$가 범주적 모듈라이 공간이므로
보조정리에서와 같은 가환 도식을 얻는다. 보조정리
\ref{stacks-more-morphisms-lemma-etale-local-finite-inertia}에서와 같은 $I$와
$g'_i : \mathcal{X}'_i \to \mathcal{X}'$를 택한다.
$g_i = h \circ g'_i$는 에탈이고 분리이며 자기동형군 위의 동형사상을
유도함에 유의하자(『스택의 사상』의 주석
\ref{stacks-morphisms-remark-identify-automorphism-groups}).
$f'_i : \mathcal{X}'_i \to M'_i$를 보조정리
\ref{stacks-more-morphisms-lemma-well-nigh-affine-moduli-space}에서와 같이 잡자.
정리 \ref{stacks-more-morphisms-theorem-keel-mori}의 증명에서 도식들
$$
\xymatrix{
\mathcal{X}'_i \ar[d]_{f'_i} \ar[r]_{g'_i} &
\mathcal{X}' \ar[d]^{f'} \\
M'_i \ar[r] &
M'
}
\quad\text{그리고}\quad
\xymatrix{
\mathcal{X}'_i \ar[d]_{f'_i} \ar[r]_{g_i} &
\mathcal{X} \ar[d]^f \\
M'_i \ar[r] &
M
}
$$
이 카르테시안이고 $M'_i \to M'$와 $M'_i \to M$이 에탈임을 보았다
(지금까지 열거한 사상 $g'_i, g_i, f', f'_i, f$의 성질만으로도 똑같이
논증하여 이를 직접 얻을 수 있다). 이로부터 첫째로 $M' \to M$이
에탈이고, 둘째로 보조정리의 도식이 카르테시안임을 얻는다.
$M' \to M$이 분리임은 여전히 보여야 한다. 이를 위해 도식
$$
\xymatrix{
\mathcal{X}' \ar[r] \ar[d] &
\mathcal{X}' \times_\mathcal{X} \mathcal{X}' \ar[d] \\
M' \ar[r] &
M' \times_M M'
}
$$
을 살펴본다. $\mathcal{X}' \to \mathcal{X}$가 분리이므로 위쪽 가로
화살표는 보편 닫힘이다. 세로 화살표들은
$\mathcal{X} \to M$의 평탄 밑변환으로서 정리
\ref{stacks-more-morphisms-theorem-keel-mori}에서와 같은 화살표이므로 보편 위상동형이다.
따라서 아래쪽 가로 화살표는 보편 닫힘이다. 이 사실과 이 화살표가
대수공간의 에탈 단사상이라는 사실을 합치면, 원하는 대로 이 화살표가
닫힌 몰입임을 얻는다.
\end{proof}

\section{모듈라이 공간의 성질}
\label{stacks-more-morphisms-section-properties-moduli-spaces}

\noindent
모듈라이 공간의 존재가 증명되고 나면 이 모듈라이 공간들의 성질을
확립할 수 있게 되며, 보통 이는 간단하다.

\begin{lemma}
\label{stacks-more-morphisms-lemma-keel-mori-finite-type}
$p : \mathcal{X} \to Y$를 대수적 스택에서 대수공간으로 가는 사상이라 하자.
다음을 가정하자.
\begin{enumerate}
\item $\mathcal{I}_\mathcal{X} \to \mathcal{X}$는 유한이다.
\item $Y$는 국소 뇌터이다.
\item $p$는 국소 유한형이다.
\end{enumerate}
$f : \mathcal{X} \to M$을 정리 \ref{stacks-more-morphisms-theorem-keel-mori}에서 구성한
모듈라이 공간이라 하자. 그러면 $M \to Y$는 국소 유한형이다.
\end{lemma}

\begin{proof}
$f$가 균일 범주적 모듈라이 공간이므로 사상 $M \to Y$를 얻는다.
$M$과 $Y$ 위에서 에탈 국소적으로 $M \to Y$가 국소 유한형임을
확인하면 충분하다. $f$가 균일 범주적 모듈라이 공간이므로 먼저
$Y$를 $Y$ 위에서 에탈인 아핀 스킴으로 바꾸어도 된다.
다음으로 보조정리 \ref{stacks-more-morphisms-lemma-etale-local-finite-inertia}에서와 같은
$I$와 $g_i : \mathcal{X}_i \to \mathcal{X}$를 택할 수 있다.
그러면 보조정리 \ref{stacks-more-morphisms-lemma-etale-separated-over-keel-mori}에 의해
$\mathcal{X} = \mathcal{X}_i$인 경우로 환원된다.
달리 말하면 $\mathcal{X}$가 거의 아핀이라고 가정할 수 있다.
이 경우 $Y = \Spec(A_0)$이고, $U = \Spec(A)$인
$\mathcal{X} = [U/R]$이 있으며, $C \subset A$가 $U$ 위의
$R$-불변 함수들의 집합일 때 $M = \Spec(C)$이다.
보조정리 \ref{stacks-more-morphisms-lemma-well-nigh-affine}와
\ref{stacks-more-morphisms-lemma-well-nigh-affine-moduli-space}를 보라.
그러면 $A_0$은 뇌터이고 $A_0 \to A$는 유한형이다.
더욱이 『준군』의 보조정리
\ref{groupoids-lemma-integral-over-invariants}에 의해 $A$는 $C$ 위에서
정수적이고, 따라서 $C$ 위에서 유한이다($A_0$ 위에서 유한형이기 때문이다).
이제 『대수학』의 보조정리 \ref{algebra-lemma-Artin-Tate}를 적용하여
결론을 얻는다.
\end{proof}

\begin{lemma}
\label{stacks-more-morphisms-lemma-keel-mori-diagonal}
$\mathcal{X}$를 대수적 스택이라 하자.
$\mathcal{I}_\mathcal{X} \to \mathcal{X}$가 유한이라고 가정하자.
$f : \mathcal{X} \to M$을 정리 \ref{stacks-more-morphisms-theorem-keel-mori}에서 구성한
모듈라이 공간이라 하자.
\begin{enumerate}
\item $\mathcal{X}$가 준분리이면 $M$은 준분리이다.
\item $\mathcal{X}$가 분리이면 $M$은 분리이다.
\item 예를 들어 위 결과들의 상대적 형태 등을 여기에 더 추가한다.
\end{enumerate}
\end{lemma}

\begin{proof}
이를 증명하기 위해 다음 도식을 생각하자.
$$
\xymatrix{
\mathcal{X} \ar[d]_f \ar[r]_{\Delta_\mathcal{X}} &
\mathcal{X} \times \mathcal{X} \ar[d]^{f \times f} \\
M \ar[r]^{\Delta_M} &
M \times M
}
$$
$f$가 보편 위상동형이므로 $f \times f$도 보편 위상동형이다.

\medskip\noindent
$\mathcal{X}$가 분리이면 $\Delta_\mathcal{X}$가 고유이고,
따라서 $\Delta_\mathcal{X}$가 보편 닫힘이고, 따라서 $\Delta_M$이
보편 닫힘이며, 따라서 『대수공간의 사상』의 보조정리
\ref{spaces-morphisms-lemma-separated-diagonal-proper}에 의해 $M$은 분리이다.

\medskip\noindent
$\mathcal{X}$가 준분리이면 $\Delta_\mathcal{X}$가 준콤팩트이고,
따라서 $\Delta_M$이 준콤팩트이며, 따라서 $M$은 준분리이다.
\end{proof}

\begin{lemma}
\label{stacks-more-morphisms-lemma-keel-mori-proper}
$p : \mathcal{X} \to Y$를 대수적 스택에서 대수공간으로 가는 사상이라 하자.
다음을 가정하자.
\begin{enumerate}
\item $\mathcal{I}_\mathcal{X} \to \mathcal{X}$는 유한이다.
\item $p$는 고유이다.
\item $Y$는 국소 뇌터이다.
\end{enumerate}
$f : \mathcal{X} \to M$을 정리 \ref{stacks-more-morphisms-theorem-keel-mori}에서 구성한
모듈라이 공간이라 하자. 그러면 $M \to Y$는 고유이다.
\end{lemma}

\begin{proof}
보조정리 \ref{stacks-more-morphisms-lemma-keel-mori-finite-type}에 의해
$M \to Y$는 국소 유한형이다. 보조정리
\ref{stacks-more-morphisms-lemma-keel-mori-diagonal}에 의해 $M \to Y$는 분리이다.
물론 $M \to Y$는 준콤팩트이고 보편 닫힘이다. 이들은 위상적 성질이고
$\mathcal{X} \to Y$가 이 성질들을 가지며
$\mathcal{X} \to M$이 보편 위상동형이기 때문이다.
\end{proof}





\section{스택과 fpqc 덮개}
\label{stacks-more-morphisms-section-fpqc}

\noindent
어떤 대수적 스택들은 fpqc 내림을 만족한다. 대수공간에 대해 이 절에
대응하는 내용은 『대수공간의 성질』의 절
\ref{spaces-properties-section-fpqc}이다.


\begin{proposition}
\label{stacks-more-morphisms-proposition-stack-fpqc}
\begin{reference}
Anatoly Preygel의 ``대수적 스택 입문''의 명제 3.3.6.
\end{reference}
$\mathcal{X}$를 대각사상이 준아핀인\footnote{ind-준아핀이라고만
가정해도 충분하다.} 대수적 스택이라 하자. 그러면
$\mathcal{X}$는 fpqc 덮개에 대한 내림을 만족한다.
\end{proposition}

\begin{proof}
우리의 규약에서 $\mathcal{X}$는 fppf 위상이 주어진 밑스킴 $S$ 위의
스킴의 범주 위에 놓인 준군 스택
$p : \mathcal{X} \to (\Sch/S)_{fppf}$이다. 명제의 뜻은 다음과 같다.
$S$ 위의 스킴의 fpqc 덮개
$\mathcal{U} = \{U_i \to U\}_{i \in I}$가 주어지면 함자
$$
\mathcal{X}_U \longrightarrow DD(\mathcal{U})
$$
가 동치이다. 왼쪽에는 $U$ 위의 $\mathcal{X}$의 대상들의 범주가 있고,
오른쪽에는 $\mathcal{U}$에 상대적인 $\mathcal{X}$ 안의 내림 자료의
범주가 있다. 『스택』의 절 \ref{stacks-section-descent-data}의 논의를 보라.

\medskip\noindent
완전충실성. $U$ 위의 $\mathcal{X}$의 두 대상 $x, y$가 있다고 하자.
그러면 $I = \mathit{Isom}(x, y)$는 $U$ 위의 대수공간이다.
$U_i$ 위의 $I$의 단면들이 주어지고 이들을
$U_i \times_U U_j$로 제한하면 일치한다고 하자. 대수공간에 대한 명제의
대응물, 즉 『대수공간의 성질』의 명제
\ref{spaces-properties-proposition-sheaf-fpqc}에 의해 이 단면들은
$U$ 위의 유일한 단면에서 온다. 따라서 우리 함자는 완전충실하다.

\medskip\noindent
본질적 전사성. 여기서는 각 $U_i$ 위의 대상 $x_i$와,
$U_i \times_U U_j$ 위에서 주어져 $U_i \times_U U_j \times_U U_k$ 위의
코사이클 조건을 만족하는 동형사상
$\varphi_{ij} : \text{pr}_0^*x_i \to \text{pr}_1^*x_j$가 주어져 있다.

\medskip\noindent
$W$를 아핀 스킴이라 하고 $W \to \mathcal{X}$를 사상이라 하자.
각 $i$에 대해
$$
W_i = U_i \times_{x_i, \mathcal{X}} W
$$
를 만들 수 있다. $\mathcal{X}$의 대각사상이 준아핀이므로
사영 $W_i \to U_i$는 준아핀이다. 각 쌍 $i, j \in I$에 대해
동형사상 $\varphi_{ij}$는 동형사상
$$
W_i \times_U U_j =
(U_i \times_U U_j) \times_{x_i \circ \text{pr}_0, \mathcal{X}} W
\to
(U_i \times_U U_j) \times_{x_j \circ \text{pr}_1, \mathcal{X}} W =
U_i \times_U W_j
$$
을 유도한다. 더욱이 이 동형사상들은
$U_i \times_U U_j \times_U U_k$ 위의 코사이클 조건을 만족한다.
달리 말하면 이 동형사상들은 $\mathcal{U}$에 상대적인 스킴
$W_i/U_i$ 위의 내림 자료를 정의한다. 『내림』의 보조정리
\ref{descent-lemma-quasi-affine}에 의해 이 내림 자료는 유효하다\footnote{
대각사상이 ind-준아핀인 경우에는 『준군에 관한 추가 내용』의 보조정리
\ref{more-groupoids-lemma-ind-quasi-affine}를 사용할 수도 있다.}.
따라서 준아핀 사상 $W' \to U$와, 두 정사각형이 카르테시안인 가환 도식
$$
\xymatrix{
W' \ar[d] &
W_i \ar[l] \ar[d] \ar[r] &
W \ar[d] \\
U &
U_i \ar[l] \ar[r]^{x_i} &
\mathcal{X}
}
$$
이 존재한다. $\{W_i \to W'\}_{i \in I}$는 $W' \to U$에 의한
$\mathcal{U}$의 밑변환이므로 fpqc 덮개이다.
$W$가 fpqc 덮개에 대한 층 조건을 만족하므로, 합성
$W_i \to W' \to W$가 주어진 사상 $W_i \to W$가 되게 하는 유일한 사상
$W' \to W$를 얻는다. 달리 말하면, $\varphi_{ij}$와 양립하고
정사각형과 직사각형이 카르테시안인 가환 도식
$$
\xymatrix{
W_i \ar[d] \ar[r] &
W' \ar[d] \ar[r] &
W \ar[d] \\
U_i \ar[r] \ar@/_1pc/[rr]_{x_i} &
U &
\mathcal{X}
}
$$
을 얻는다.

\medskip\noindent
아핀 스킴들의 모임 $W_\alpha$ ($\alpha \in A$)와 매끄러운 사상
$W_\alpha \to \mathcal{X}$들을 택하여
$\coprod W_\alpha \to \mathcal{X}$가 전사가 되게 한다.
앞 문단의 절차로 각 $\alpha$에 대해 도식
$$
\xymatrix{
W_{\alpha, i} \ar[d] \ar[r] &
W_\alpha' \ar[d] \ar[r] &
W_\alpha \ar[d] \\
U_i \ar[r] \ar@/_1pc/[rr]_{x_i} &
U &
\mathcal{X}
}
$$
을 만든다. 그러면 사상 $W_\alpha' \to U$들은 매끄럽고 함께 전사이다.

\medskip\noindent
$W_\alpha' \to W_\alpha \to \mathcal{X}$에 대응하는
$W_\alpha'$ 위의 $\mathcal{X}$의 대상을 $x_\alpha$라 쓰자.
$\mathcal{X}$가 fppf 스택이고 $\{W_\alpha' \to U\}$가 fppf 덮개이므로,
$W_\alpha' \times_U W'_\beta$ 위에서 코사이클 조건을 만족하는 동형사상
$\text{pr}_0^*x_\alpha \to \text{pr}_1^*x_\beta$가 존재함을 보이면 충분하다.
그런데 $x_\alpha$를 $W_{\alpha, i}$로 당기면 $x_i$를
$W_{\alpha, i}$로 당긴 것과 동형이므로, $W_{\alpha, i}$로 당긴 뒤에는
실제로
$W_{\alpha, i} \times_{U_i} W_{\beta, i} =
U_i \times_U (W_\alpha' \times_U W'_\beta)$
위에서 그러한 동형사상들이 존재한다.
$\{U_i \times_U (W_\alpha' \times_U W'_\beta) \to
W_\alpha' \times_U W'_\beta\}_{i \in I}$가 fpqc 덮개이고,
위 도식들이 $\varphi_{ij}$와 양립한다는 앞서 말한 사실에 의해
이 동형사상들은 $W_\alpha' \times_U W'_\beta$로 내려간다.
이로써 증명이 끝난다.
\end{proof}









\section{텐서 함자}
\label{stacks-more-morphisms-section-tensor-functors}

\noindent
$f : \mathcal{Y} \to \mathcal{X}$를 뇌터 대수적 스택의 사상이라 하자.
당김 함자
$$
f^* :
\textit{Coh}(\mathcal{O}_\mathcal{X})
\longrightarrow
\textit{Coh}(\mathcal{O}_\mathcal{Y})
$$
는 오른쪽 완전 텐서 함자이다. 즉, 가법이고 오른쪽 완전이며 연접 가군의
텐서곱과 가환한다. 임의의 오른쪽 완전 텐서 함자
$F : \textit{Coh}(\mathcal{O}_\mathcal{X}) \to
\textit{Coh}(\mathcal{O}_\mathcal{Y})$가 어느 정도까지 사상
$f : \mathcal{Y} \to \mathcal{X}$에서 오는지 물을 수 있다.
이런 종류의 매우 일반적인 결과는 \cite{Hall-Rydh-coherent}를 참고하라.
이 절의 목적은 이러한 생각의 입문으로서 정리
\ref{stacks-more-morphisms-theorem-main}의 짧은 증명을 제시하는 것이다.

\medskip\noindent
몇 가지 보조정리로 시작한다.

\begin{lemma}
\label{stacks-more-morphisms-lemma-extend}
$\mathcal{X}$와 $\mathcal{Y}$를 뇌터 대수적 스택이라 하자.
임의의 오른쪽 완전 텐서 함자
$F : \textit{Coh}(\mathcal{O}_\mathcal{X}) \to
\textit{Coh}(\mathcal{O}_\mathcal{Y})$는 모든 쌍대극한과 가환하는
오른쪽 완전 텐서 함자
$F : \QCoh(\mathcal{O}_\mathcal{X}) \to \QCoh(\mathcal{O}_\mathcal{Y})$로 유일하게 확장된다.
\end{lemma}

\begin{proof}
그러한 확장의 존재성과 유일성은 일반적인 사실이다.
『범주』의 보조정리 \ref{categories-lemma-extend-functor-by-colim}를 보라.
보조정리를 적용할 수 있음을 보이기 위해 다음을 관찰하자.
국소 뇌터 대수적 스택 위의 연접 가군은 정의상 유한 표시 가군이다.
『스택의 코호몰로지』의 정의
\ref{stacks-cohomology-definition-coherent}를 보라.
따라서 『스택의 코호몰로지』의 보조정리
\ref{stacks-cohomology-lemma-finite-presentation-quasi-compact-colimit}에 의해
$\mathcal{X}$ 위의 연접 가군은 $\QCoh(\mathcal{O}_\mathcal{X})$의
범주적으로 콤팩트한 대상이다. 마지막으로 『스택의 코호몰로지』의 보조정리
\ref{stacks-cohomology-lemma-directed-colimit-coherent}에 의해 모든
준연접 가군은 그 연접 부분가군들의 여과 쌍대극한이다.

\medskip\noindent
$F$가 가법이므로 $F$의 확장도 가법이다(세부사항은 생략한다).
$F$가 텐서 함자이고 가군의 쌍대극한이 텐서곱을 취하는 것과 가환하므로,
$F$의 확장도 텐서 함자이다(세부사항은 생략한다).

\medskip\noindent
이 문단에서는 그 확장이 임의의 직합과 가환함을 보인다.
$\mathcal{H}_j$들이 준연접이고
$\mathcal{F} = \bigoplus_{j \in J} \mathcal{H}_j$이면
$\mathcal{F} = \colim_{J' \subset J\text{ 유한}}
\bigoplus_{j \in J'} \mathcal{H}_j$이다.
$F$의 확장도 $F$로 쓰면
\begin{align*}
F(\mathcal{F})
& =
\colim_{J' \subset J\text{ 유한}}
F(\bigoplus\nolimits_{j \in J'} \mathcal{H}_j) \\
& =
\colim_{J' \subset J\text{ 유한}}
\bigoplus\nolimits_{j \in J'} F(\mathcal{H}_j) \\
& =
\bigoplus\nolimits_{j \in J} F(\mathcal{H}_j)
\end{align*}
를 얻는다. 따라서 $F$는 임의의 직합과 가환한다.

\medskip\noindent
이 문단에서는 그 확장이 오른쪽 완전임을 보인다.
$0 \to \mathcal{F} \to \mathcal{F}' \to \mathcal{F}'' \to 0$을
준연접 $\mathcal{O}_\mathcal{X}$-가군의 짧은 완전열이라 하자.
$\mathcal{F}' = \bigcup \mathcal{F}'_i$를 그 연접 부분가군들의 합집합으로
쓰자(위에서 제시한 참고문헌을 보라). $\mathcal{F}''_i \subset \mathcal{F}''$를 $\mathcal{F}'_i$의 상이라 하고,
$\mathcal{F}_i = \mathcal{F} \cap \mathcal{F}'_i =
\Ker(\mathcal{F}'_i \to \mathcal{F}''_i)$라 하자.
그러면 $\mathcal{F} = \bigcup \mathcal{F}_i$와
$\mathcal{F}'' = \bigcup \mathcal{F}''_i$임이 분명하고, 짧은 완전열
$$
0 \to \mathcal{F}_i \to \mathcal{F}_i' \to \mathcal{F}_i'' \to 0
$$
을 얻는다. 그 확장이 여과 쌍대극한과 가환하므로
$F(\mathcal{F}) = \colim_{i \in I} F(\mathcal{F}_i)$,
$F(\mathcal{F}') = \colim_{i \in I} F(\mathcal{F}'_i)$, 그리고
$F(\mathcal{F}'') = \colim_{i \in I} F(\mathcal{F}''_i)$이다.
가군의 층의 여과 쌍대극한은 완전이므로 $F$의 확장은 오른쪽 완전이다.

\medskip\noindent
모든 쌍대곱과 가환하는 오른쪽 완전 함자는 모든 쌍대극한과 가환하므로
증명이 끝난다. 『범주』의 보조정리
\ref{categories-lemma-colimits-coproducts-coequalizers}를 보라.
\end{proof}

\begin{lemma}
\label{stacks-more-morphisms-lemma-affine}
$\mathcal{X}$를 대각사상이 아핀인 대수적 스택이라 하자.
$B$를 환이라 하자. $F : \QCoh(\mathcal{O}_\mathcal{X}) \to \text{Mod}_B$를
직합과 가환하는 오른쪽 완전 텐서 함자라 하자.
$g : U \to \mathcal{X}$를 $U = \Spec(A)$가 아핀인 사상이라 하자.
그러면
\begin{enumerate}
\item $C = F(g_{\QCoh, *}\mathcal{O}_U)$는 가환 $B$-대수이고,
\item 환 사상 $A \to C$가 존재한다.
\end{enumerate}
이때 $F \circ g_{\QCoh, *} : \text{Mod}_A \to \text{Mod}_B$는
$M$을 $B$-가군으로 본 $M \otimes_A C$로 보낸다.
\end{lemma}

\begin{proof}
$g$는 준콤팩트이고 준분리임에 유의하자. 『스택의 사상』의 보조정리
\ref{stacks-morphisms-lemma-quasi-compact-quasi-separated-permanence}를 보라.
『스택의 코호몰로지』의 명제
\ref{stacks-cohomology-proposition-direct-image-quasi-coherent}에서 함자
$g_{\QCoh, *} : \QCoh(\mathcal{O}_U) \to \QCoh(\mathcal{O}_\mathcal{X})$를 구성했다.
『스택의 코호몰로지』의 주석
\ref{stacks-cohomology-remark-direct-image-quasi-coherent-tensor}와
\ref{stacks-cohomology-remark-QCoh-tensor}에 의해 곱셈
$$
\mu :
g_{\QCoh, *}\mathcal{O}_U
\otimes_{\mathcal{O}_\mathcal{X}}
g_{\QCoh, *}\mathcal{O}_U
\longrightarrow
g_{\QCoh, *}\mathcal{O}_U
$$
을 얻고, 이는 $g_{\QCoh, *}\mathcal{O}_U$를 가환
$\mathcal{O}_\mathcal{X}$-대수로 만든다. 따라서
$C = F(g_{\QCoh, *}\mathcal{O}_U)$는 $\text{Mod}_B$ 안의 가환 대수 대상,
달리 말하면 $C$는 가환 $B$-대수이다. 사상
$\kappa : A \to \text{End}_{\mathcal{O}_\mathcal{X}}(g_{\QCoh, *}\mathcal{O}_U)$가 존재하고, 임의의 $a \in A$에 대해 도식
$$
\xymatrix{
g_{\QCoh, *}\mathcal{O}_U \otimes_{\mathcal{O}_\mathcal{X}}
g_{\QCoh, *}\mathcal{O}_U
\ar[d]_{\kappa(r) \otimes 1} \ar[rr]_-\mu & &
g_{\QCoh, *}\mathcal{O}_U \ar[d]^{\kappa(r)} \\
g_{\QCoh, *}\mathcal{O}_U \otimes_{\mathcal{O}_\mathcal{X}}
g_{\QCoh, *}\mathcal{O}_U
\ar[rr]^-\mu & &
g_{\QCoh, *}\mathcal{O}_U
}
$$
이 가환함에 유의하자. 따라서 사상
$\kappa' = F(\kappa) : A \to \text{End}_B(C)$를 얻으며,
$\kappa'(a)(c) c' = \kappa'(a)(cc')$이다. 물론 이는
$a \mapsto \kappa'(a)(1)$이 환 사상 $A \to C$라는 뜻이다.

\medskip\noindent
사상 $g : U \to \mathcal{X}$는 아핀이다. 『스택의 사상』의 보조정리
\ref{stacks-morphisms-lemma-affine-permanence}를 보라.
따라서 『스택의 코호몰로지』의 보조정리
\ref{stacks-cohomology-lemma-quasi-coherent-pushforward-affine}에 의해
$g_{\QCoh, *}$는 완전이고 직합과 가환한다.
그러므로 $F \circ g_{\QCoh, *} : \text{Mod}_A \to \text{Mod}_B$는
직합과 가환하는 오른쪽 완전 함자이고 $A$를 $C$로 보낸다.
『함자와 사상』의 보조정리 \ref{functors-lemma-functor}에 의해
$F \circ g_{\QCoh, *}$는 $A$-가군 $M$을 $B$-가군으로 본
$M \otimes_A C$로 보낸다.
\end{proof}

\begin{lemma}
\label{stacks-more-morphisms-lemma-universally-injective}
보조정리 \ref{stacks-more-morphisms-lemma-affine}의 표기법을 따르자.
$\mathcal{X}$가 뇌터이고 $g$가 전사이고 평탄하다고 가정하자.
그러면 $B \to C$는 보편 단사이다.
\end{lemma}

\begin{proof}
$\QCoh(\mathcal{O}_\mathcal{X})$ 안의 자연스러운 사상
$1 : \mathcal{O}_\mathcal{X} \to g_{\QCoh, *}\mathcal{O}_U$를 생각하자.
$U$로 당기고 수반을 사용하면 합성
$$
\mathcal{O}_U =
g^*\mathcal{O}_\mathcal{X} \xrightarrow{g^*1}
g^*g_{\QCoh, *}\mathcal{O}_U \to \mathcal{O}_U
$$
이 $\QCoh(\mathcal{O}_U)$ 안의 항등사상임을 알 수 있다.
$g_{\QCoh, *}\mathcal{O}_U = \colim \mathcal{F}_i$를 연접
$\mathcal{O}_\mathcal{X}$-부분가군들의 여과 쌍대극한으로 쓰자.
『스택의 코호몰로지』의 보조정리
\ref{stacks-cohomology-lemma-directed-colimit-coherent}를 보라.
충분히 큰 $i$에 대해 사상
$1 : \mathcal{O}_\mathcal{X} \to g_{\QCoh, *}\mathcal{O}_U$는
$\mathcal{F}_i$를 통해 분해된다. 『스택의 코호몰로지』의 보조정리
\ref{stacks-cohomology-lemma-finite-presentation-quasi-compact-colimit}를 보라.
그 분해를 $s : \mathcal{O}_\mathcal{X} \to \mathcal{F}_i$라 하자.
그러면
$$
\mathcal{O}_U \xrightarrow{g^*s} g^*\mathcal{F}_i \to
g^*g_{\QCoh, *}\mathcal{O}_U \to \mathcal{O}_U
$$
는 항등사상이다. 달리 말하면 $s$를 $U$로 당기면 직합 인자의 포함이 된다.
『스택의 코호몰로지』의 보조정리
\ref{stacks-cohomology-lemma-internal-hom-fp-into-qcoh}의 표기법으로
$\mathcal{F}_i^\vee = hom(\mathcal{F}_i, \mathcal{O}_\mathcal{X})$로 두자.
특히 평가 사상
$ev : \mathcal{F}_i \otimes_{\mathcal{O}_\mathcal{X}} \mathcal{F}_i^\vee
\to \mathcal{O}_\mathcal{X}$가 있다.
$s$에서의 평가는 사상
$s^\vee : \mathcal{F}_i^\vee \to \mathcal{O}_\mathcal{X}$를 정의한다.
$s$에 관한 명제의 쌍대 명제에 따라 $g^*(s^\vee)$는 전사이다.
$hom$과 $\otimes$가 $U$로의 제한과 양립한다는 것은
『스택의 코호몰로지』의 절
\ref{stacks-cohomology-section-further-remarks}을 보라.
$g$가 전사이고 평탄하므로 $s^\vee$가 전사라고 결론짓는다
(앞서 인용한 곳을 보라). $F$가 오른쪽 완전이므로
$F(\mathcal{F}_i^\vee) \to F(\mathcal{O}_\mathcal{X}) = B$는 전사이다.
$1 \in B$로 가는 $\lambda \in F(\mathcal{F}_i^\vee)$를 택한다.
$F(s) : B = F(\mathcal{O}_\mathcal{X}) \to F(\mathcal{F}_i)$ 아래에서
$1$의 상을 $e = F(s)(1) \in F(\mathcal{F}_i)$라 쓰자. 그러면 사상
$$
F(ev) :
F(\mathcal{F}_i) \otimes_B F(\mathcal{F}_i^\vee) =
F(\mathcal{F}_i \otimes_{\mathcal{O}_\mathcal{X}} \mathcal{F}_i^\vee)
\longrightarrow
F(\mathcal{O}_\mathcal{X}) = B
$$
은 구성상 $e \otimes \lambda$를 $1$로 보낸다. 따라서 사상
$B \to F(\mathcal{F}_i)$, $b \mapsto be$는 보편 단사이다.
실제로 그 한쪽 역은
$F(\mathcal{F}_i) \to B$, $\xi \mapsto F(ev)(\xi \otimes \lambda)$이다.
이는 충분히 큰 모든 $i$에 대해 참이므로 결론을 얻는다.
\end{proof}

\begin{lemma}
\label{stacks-more-morphisms-lemma-flat}
$B \to C$를 환 사상이라 하자. 다음이 성립한다고 하자.
\begin{enumerate}
\item 쌍대사영 $C \to C \otimes_B C$는 평탄하다.
\item $B \to C$는 보편 단사이다.
\end{enumerate}
그러면 $B \to C$는 충실히 평탄하다.
\end{lemma}

\begin{proof}
$B \to C$가 보편 단사이므로 사상
$\Spec(C) \to \Spec(B)$는 전사이다. 따라서 $B \to C$가 평탄임을
보이면 충분하며, 이는 『내림』의 정리
\ref{descent-theorem-descend-module-properties}에서 따른다.
\end{proof}

\noindent
다음의 아주 간단한 퀸네트 정리는 대수적 스택 위의 준연접 가군의
코호몰로지에 대한 퀸네트 정리들을 다루는 절을 쓰고 나면
쓸모없어져야 한다.

\begin{lemma}
\label{stacks-more-morphisms-lemma-easy-kunneth}
$a : \mathcal{Y} \to \mathcal{X}$와 $b : \mathcal{Z} \to \mathcal{X}$가
스킴으로 표현가능하고 준콤팩트이고 준분리이며 평탄하다고 하자.
그러면
$a_{\QCoh, *}\mathcal{O}_\mathcal{Y}
\otimes_{\mathcal{O}_\mathcal{X}}
b_{\QCoh, *}\mathcal{O}_\mathcal{Z} =
f_{\QCoh, *}\mathcal{O}_{\mathcal{Y} \times_\mathcal{X} \mathcal{Z}}$이다.
여기서 $f : \mathcal{Y} \times_\mathcal{X} \mathcal{Z} \to \mathcal{X}$는
명백한 사상이다.
\end{lemma}

\begin{proof}
$\mathcal{P} = \mathcal{Y} \times_\mathcal{X} \mathcal{Z}$로 줄여 쓰자.
$a \circ \text{pr}_1 = f$이고 $b \circ \text{pr}_2 = f$이므로 사상
$a_*\mathcal{O}_\mathcal{Y} \to f_*\mathcal{O}_\mathcal{P}$와
$b_*\mathcal{O}_\mathcal{Z} \to f_*\mathcal{O}_\mathcal{P}$를 얻는다
(상대 당김 사상을 사용한다. 『사이트』의 절
\ref{sites-section-pullback}을 보라). 따라서 상대 컵곱
$$
\mu :
a_*\mathcal{O}_\mathcal{Y}
\otimes_{\mathcal{O}_\mathcal{X}}
b_*\mathcal{O}_\mathcal{Z} \longrightarrow
f_*\mathcal{O}_{\mathcal{Y} \times_\mathcal{X} \mathcal{Z}}
$$
을 얻는다. $Q$를 적용하고 텐서곱과의 양립성
(『스택의 코호몰로지』의 주석
\ref{stacks-cohomology-remark-QCoh-tensor})을 사용하여
$\QCoh(\mathcal{O}_\mathcal{X})$ 안의 화살표
$Q(\mu) : a_{\QCoh, *}\mathcal{O}_\mathcal{Y}
\otimes_{\mathcal{O}_\mathcal{X}}
b_{\QCoh, *}\mathcal{O}_\mathcal{Z} \to
f_{\QCoh, *}\mathcal{O}_{\mathcal{Y} \times_\mathcal{X} \mathcal{Z}}$
를 얻는다. 다음으로 스킴 $U$와 전사 매끄러운 사상
$U \to \mathcal{X}$를 택한다. $Q(\mu)$를 $U_\etale$로 제한한 것이
동형사상임을 보이면 충분하다. 『스택의 코호몰로지』의 절
\ref{stacks-cohomology-section-further-remarks}을 보라.
다시 같은 절의 내용에 의해 $\mu$를 $U_\etale$로 제한한 것이
동형사상임을 보이면 충분하다(여기서는 $\mu$의 원천과 표적이
밑변환 성질을 갖는 국소 준연접 가군이라는 사실을 사용한다).
더욱이 에탈 위상에서 순상을 계산할 수 있다.
『스택의 코호몰로지』의 명제
\ref{stacks-cohomology-proposition-loc-qcoh-flat-base-change}를 보라.
마지막으로 $V = U \times_\mathcal{X} \mathcal{Y}$일 때
$a_*\mathcal{O}_\mathcal{Y}|_{U_\etale} = (V \to U)_{small, *}\mathcal{O}_V$임을 알 수 있다.
$b_*$와 $f_*$에 대해서도 마찬가지이다. 따라서 결과는 평탄하고
준콤팩트이며 준분리인 스킴의 사상에 대한 퀸네트 공식에서 따른다.
『스킴의 유도범주』의 보조정리 \ref{perfect-lemma-kunneth}를 보라.
\end{proof}

\begin{lemma}
\label{stacks-more-morphisms-lemma-fully-faithful}
$\mathcal{X}$를 대각사상이 아핀인 대수적 스택이라 하자.
$B$를 환이라 하자. $f_i : \Spec(B) \to \mathcal{X}$ ($i = 1, 2$)를
두 사상이라 하자. $t : f_1^* \to f_2^*$를 텐서 함자
$f_i^* : \QCoh(\mathcal{O}_\mathcal{X}) \to \text{Mod}_B$의
동형사상이라 하자. 그러면 $t$를 유도하는 $2$-화살표
$f_1 \to f_2$가 존재한다.
\end{lemma}

\begin{proof}
아핀 스킴 $U = \Spec(A)$와 전사 매끄러운 사상
$g : U \to \mathcal{X}$를 택한다. 『스택의 성질』의 보조정리
\ref{stacks-properties-lemma-quasi-compact-stack}를 보라.
$\mathcal{X}$의 대각사상이 아핀이므로
$U_i = \Spec(B) \times_{f_i, \mathcal{X}, g} U$는 아핀이다.
$U_i = \Spec(C_i)$라 하자. 그러면 $C_i$는 함자 $F = f_i^*$를 사용하여
보조정리 \ref{stacks-more-morphisms-lemma-affine}에서 구성한 환 사상 $A \to C_i$를 갖는
$B$-대수이다. 따라서 $t$는 환 사상 $A \to C_1$과 $A \to C_2$와
양립하는 $B$-대수의 동형사상 $C_1 \to C_2$를 유도한다.
달리 말하면 다음 가환 도식들이 있다.
$$
\xymatrix{
U_i \ar[r] \ar[d] & U \ar[d]^g \\
\Spec(B) \ar[r]^{f_i} & \mathcal{X}
}
\quad\quad
\xymatrix{
&
U_2 \ar[ld] \ar[d]^{\cong} \ar[rd] \\
\Spec(B) &
U_1 \ar[l] \ar[r] &
U
}
$$
이로써 $\Spec(B)$ 위의 $\mathcal{X}$의 대상 $f_1$과 $f_2$가
매끄러운 덮개 $\{U_1 \to \Spec(B)\}$ 뒤에 동형이 됨을 이미 알 수 있다.
이 동형사상이 $\Spec(B)$ 위의 $f_1$과 $f_2$의 동형사상으로 내려감을
보이려면, 우리의 동형사상(위 가환 도식들에서 오는 것)이
$U_1 \times_{\Spec(B)} U_1$ 위의 내림 자료와 양립함을 보여야 한다.
이를 위해 $U \times_\mathcal{X} U$도 아핀이고, 사상
$g' : U \times_\mathcal{X} U \to \mathcal{X}$가 있으며,
$$
U_i \times_{\Spec(B)} U_i =
\Spec(B) \times_{f_i, \mathcal{X}, g'}
(U \times_\mathcal{X} U)
$$
임에 유의한다. 따라서 동형사상 $C_1 \to C_2$에서 오는 동형사상
$C_1 \otimes_B C_1 \to C_2 \otimes_B C_2$는 사상
$U_i \times_{\Spec(B)} U_i \to U \times_\mathcal{X} U$와 양립한다.
몇 가지 세부사항은 생략한다.
\end{proof}

\begin{lemma}
\label{stacks-more-morphisms-lemma-main}
$\mathcal{X}$를 대각사상이 아핀인 뇌터 대수적 스택이라 하자.
$B$를 환이라 하자.
$F : \QCoh(\mathcal{O}_\mathcal{X}) \to \text{Mod}_B$를 직합과 가환하는
오른쪽 완전 텐서 함자라 하자. 그러면 $F$는 유일한 사상
$\Spec(B) \to \mathcal{X}$에서 온다.
\end{lemma}

\begin{proof}
$U = \Spec(A)$가 아핀인 전사 매끄러운 사상
$g : U \to \mathcal{X}$를 택한다. 『스택의 성질』의 보조정리
\ref{stacks-properties-lemma-quasi-compact-stack}를 보라.
보조정리 \ref{stacks-more-morphisms-lemma-affine}를 적용하여 유한형 가환 $B$-대수
$C = F(g_{\QCoh, *}\mathcal{O}_U)$와 환 사상 $A \to C$를 얻는다.
보조정리 \ref{stacks-more-morphisms-lemma-universally-injective}에 의해 환 사상
$B \to C$는 보편 단사이다. 대수
$$
C \otimes_B C =
F(g_{\QCoh, *}\mathcal{O}_U
\otimes_{\mathcal{O}_\mathcal{X}}
g_{\QCoh, *}\mathcal{O}_U)
$$
를 생각하자. $g$가 평탄하고 준콤팩트이며 준분리이므로
보조정리 \ref{stacks-more-morphisms-lemma-easy-kunneth}에 의해 다음 식의 첫 번째 등식이 성립한다.
$$
g_{\QCoh, *}\mathcal{O}_U
\otimes_{\mathcal{O}_\mathcal{X}}
g_{\QCoh, *}\mathcal{O}_U
=
f_{\QCoh, *}\mathcal{O}_{U \times_\mathcal{X} U} =
g_{\QCoh, *}(\text{pr}_{2, *}\mathcal{O}_{U \times_\mathcal{X} U})
$$
여기서 $f : U \times_\mathcal{X} U \to \mathcal{X}$는 명백한 사상이고,
$\text{pr}_2 : U \times_\mathcal{X} U \to U$는 두 번째 사영이다.
두 번째 등식은 『스택의 코호몰로지』의 보조정리
\ref{stacks-cohomology-lemma-relative-leray}와
$f = g \circ \text{pr}_2$에서 따른다. $\mathcal{X}$의 대각사상이 아핀이므로
$U \times_\mathcal{X} U = \Spec(R)$은 아핀이다.
$\text{pr}_2 : A \to R$을 사용하여 $R$을 $A$-대수로 보자.
모두 합치면
$$
C \otimes_B C =
F(g_{\QCoh, *}\mathcal{O}_U
\otimes_{\mathcal{O}_\mathcal{X}}
g_{\QCoh, *}\mathcal{O}_U) =
F(g_{\QCoh, *}(\text{pr}_{2, *}\mathcal{O}_{U \times_\mathcal{X} U})) =
R \otimes_A C
$$
를 얻는다. 마지막 등식은 보조정리 \ref{stacks-more-morphisms-lemma-affine}의 마지막 명제에서
따른다. $\text{pr}_2$가 $U \to \mathcal{X}$의 밑변환으로서 평탄하므로
$A \to R$은 평탄이다. 따라서 $C \otimes_B C$는 $C$ 위에서 평탄이다.
보조정리 \ref{stacks-more-morphisms-lemma-flat}에 의해 $B \to C$는 충실히 평탄이다.

\medskip\noindent
다음 실선 가환 도식이 존재한다고 주장한다.
$$
\xymatrix{
\Spec(C \otimes_B C) \ar@<1ex>[d] \ar@<-1ex>[d] \ar[r] &
U \times_\mathcal{X} U \ar@<1ex>[d] \ar@<-1ex>[d] \\
\Spec(C) \ar[d] \ar[r] &
U \ar[d] \\
\Spec(B) \ar@{..>}[r] &
\mathcal{X}
}
$$
화살표 $\Spec(C) \to U = \Spec(A)$는 보조정리
\ref{stacks-more-morphisms-lemma-affine}의 명제에 나온 환 사상 $A \to C$에서 온다.
마찬가지로 화살표
$\Spec(C \otimes_B C) \to U \times_\mathcal{X} U$는 환 사상
$R \to C \otimes_B C$에서 온다. 위쪽 정사각형이 가환임을 확인하려면
보조정리 \ref{stacks-more-morphisms-lemma-fully-faithful}를 사용한다. 세부사항은 생략한다.
명제 \ref{stacks-more-morphisms-proposition-stack-fpqc}에 의해 점선 화살표
$\Spec(B) \to \mathcal{X}$를 얻는다고 결론짓는다.

\medskip\noindent
$F$가 점선 화살표에 의한 당김에 대응하는 함자라는 명제 역시
이 사실과 보조정리 \ref{stacks-more-morphisms-lemma-affine}의 대응하는 명제에서 명백하다.
세부사항은 생략한다.
\end{proof}

\noindent
환 $B$에 대해 유한 생성 $B$-가군(즉, 유한 $B$-가군)의 범주를
$\text{Mod}^{fg}_B$라 쓰자.

\begin{theorem}
\label{stacks-more-morphisms-theorem-main}
$\mathcal{X}$를 대각사상이 아핀인 뇌터 대수적 스택이라 하자.
$B$를 뇌터 환이라 하자. $F : \text{Coh}(\mathcal{O}_\mathcal{X}) \to \text{Mod}^{fg}_B$를 오른쪽 완전 텐서 함자라 하자.
그러면 $F$는 유일한 사상 $\Spec(B) \to \mathcal{X}$에서 온다.
\end{theorem}

\begin{proof}
보조정리 \ref{stacks-more-morphisms-lemma-extend}에 의해 $F$를 모든 직합과 가환하는
오른쪽 완전 텐서 함자
$F : \QCoh(\mathcal{O}_\mathcal{X}) \to \text{Mod}_B$로 유일하게
확장할 수 있다. 이제 보조정리 \ref{stacks-more-morphisms-lemma-main}을 적용할 수 있다.
\end{proof}
